\pdfoutput=1
\documentclass[11pt]{article}
\usepackage[OT1]{fontenc}
\usepackage{smile}
\usepackage{subcaption}
\usepackage[margin=1in]{geometry}
\usepackage[most]{tcolorbox}


%%%%%%%%%%%%-- Comments --%%%%%%%%%%%%
%%%%%%%%%%%%-- new command --%%%%%%%%%%%%
\renewcommand{\baselinestretch}{1.10}

% Custom commands matching your notation
\newcommand{\calN}{\mathcal{N}}
\newcommand{\calL}{\mathcal{L}}
\newcommand{\thetavec}{\bm{\theta}}
\newcommand{\calA}{\mathcal{A}}
\newcommand{\bbP}{\mathbb{P}}
\newcommand{\bbR}{\mathbb{R}}
\newcommand{\bbE}{\mathbb{E}}
\newcommand{\Grad}{\mathrm{Grad}}
\renewcommand{\d}{\mathrm{d}}
\newcommand{\wh}{\widehat}
\newcommand{\inv}{\mathrm{inv}}
\newcommand{\wt}{\widetilde}
\newcommand{\ol}{\overline}

\newcommand{\jg}[1]{{\color{red}#1}}
%%%%%%%%%%%%-- start --%%%%%%%%%%%%%%%
\title{\bf Learning How to Evolve: Variational Latent Dynamics for Long-Horizon Neural PDE Solvers}
% Possible title:
% (Stable by Construction:) Variational Latent Markov Operators for Long-Horizon PDE Prediction

\author{Junyi Liao}
\date{}
\begin{document}
\maketitle

\section{Introduction}
Partial differential equations (PDEs) provide a fundamental mathematical description of dynamical systems across science and engineering, but repeatedly solving complex PDEs with conventional numerical methods can be computationally demanding, particularly when solutions are required for many different initial conditions, coefficients, or forcing terms. This has motivated growing interest in data-driven PDE solvers, ranging from physics-informed neural networks to operator-learning approaches such as DeepONet, the Fourier neural operator (FNO), Operator Transformer (OFormer), and position-induced Transformer (PiT) (\citealp{raissi2019physics,lu2021learning,li2020fourier};\citealp{li2022transformer,chen2024positional}). Rather than solving each instance independently from first principles, operator-learning methods seek to learn mappings between function spaces from collections of PDE trajectories or solution pairs, enabling rapid surrogate evaluation through neural-network inference once training is complete. Depending on the underlying architecture, neural operators can further support properties such as transfer across spatial discretizations, nonlocal modeling of spatial interactions, or adaptation to irregular computational domains.

Despite these advances, reliably learning the temporal evolution of PDE systems over long prediction horizons remains challenging, particularly when a learned one-step evolution operator is deployed recursively far beyond the temporal regime represented during training. In a typical autoregressive formulation, the model is trained to advance ground-truth states by one or a few time steps, but at inference time its own predictions are repeatedly fed back as inputs. Consequently, even small one-step errors can perturb the predicted trajectory away from the states encountered during supervised training, after which subsequent predictions are made from increasingly imperfect inputs. Such discrepancies may accumulate or amplify through repeated application of the learned dynamics, leading to substantial degradation over long rollouts even when short-term prediction remains accurate. The difficulty becomes especially pronounced when training trajectories cover only a limited temporal horizon, while deployment requires evolving the system substantially beyond that horizon. Long-horizon neural PDE prediction therefore involves not only learning an accurate local approximation of the solution operator, but also maintaining stable evolution under the perturbations and distributional shifts induced by the model's own previous predictions.

Several strategies have been developed to mitigate this accumulation of autoregressive error. Multi-step and curriculum-based training schemes progressively expose the model to longer rollout horizons or to states generated by its own predictions, thereby reducing the discrepancy between training and recursive deployment \citep{li2022transformer}; recent approaches such as CALM-PDE adopt this perspective for long-term PDE forecasting \citep{hagnberger2025calm}. An alternative is to avoid recursive prediction altogether by learning a direct mapping from the input to an entire space-time solution trajectory, as in trajectory-level neural operators and related space-time architectures (\citealp{li2020fourier}; \citealp{li2022transformer}). These strategies can mitigate long-horizon prediction errors in some settings, but their effectiveness depends on the training regime and problem setting, and they address different aspects of the problem. Curriculum or multi-step training changes the distribution and horizon over which the predictor is optimized, while direct space-time prediction circumvents repeated application of a learned one-step map over a prescribed output interval, but does not naturally extend beyond the temporal regime represented during training. Neither approach, by itself, directly addresses how the learned local dynamics respond to perturbations around the observed trajectory distribution or provides an intrinsic mechanism for regularizing the stability of recursive evolution. This motivates us to consider a complementary question: whether the formulation of the learned dynamics itself can be modified to promote robustness under long-horizon autoregressive deployment.

Against this backdrop, most neural PDE solvers and operator-learning methods formulate the learned solution map or temporal evolution deterministically, while probabilistic and variational approaches remain comparatively less explored. Recent works such as Variational Autoencoding Neural Operators \citep{seidman2023variational} and Function-Space Autoencoders \citep{bunker2025autoencoders} have begun to develop probabilistic latent representations for operator learning, but have primarily focused on static mappings between functions rather than the recursive evolution of latent dynamics over time. A variational formulation is particularly appealing in the dynamic setting: representing each physical state by a distribution in latent space allows the learned transition to be trained over neighborhoods of the observed trajectory, while distributional alignment between predicted and encoded next states provides an additional mechanism for regularizing the latent evolution. These properties suggest that variational modeling may offer a principled way to improve robustness to the perturbations encountered during autoregressive rollout. However, how such probabilistic latent dynamics should be formulated for PDE evolution, and how their variational structure relates to long-horizon deterministic prediction, remain insufficiently understood.

In this work, we study variational modeling of time-dependent PDE dynamics with a particular focus on long-horizon autoregressive prediction. We formulate the evolution through probabilistic latent Markov dynamics, in which physical states are encoded into latent distributions, propagated through a stochastic transition model, and decoded back to the physical space. The resulting framework uses stochasticity during training to regularize the latent evolution, while prediction is performed through deterministic mean dynamics. We instantiate this formulation as the Variational Autoencoding Markov Operator (VAMO), using spatially resolved latent fields and structured Gaussian perturbations together with a neural-operator transition.

Our main contributions are:

\begin{itemize}[topsep=0pt,itemsep=0pt]
\item\textit{Variational latent dynamics.} We develop a function-space variational formulation for learning temporal PDE evolution through latent Markov dynamics.
\item\textit{Rollout analysis.} We characterize how latent perturbations and variational transition alignment enter the error propagation of deterministic autoregressive rollouts.
\item\textit{Long-horizon evaluation.} We instantiate the framework as VAMO and evaluate it on three fluid-dynamics benchmarks, where it improves long-horizon rollout stability over deterministic latent, direct neural-operator, and noise-injection baselines.
\end{itemize}

To summarize, our study positions variational modeling as a complementary approach to improving the robustness of learned PDE dynamics, rather than merely as a mechanism for probabilistic prediction. By connecting variational training with deterministic autoregressive evolution, we aim to broaden the role of probabilistic methods in neural PDE solvers toward long-horizon dynamical modeling.

\subsection{Related Works}
\paragraph{Neural PDE solvers and operator learning.} Neural networks have been widely used for PDE approximation through either physics-informed objectives or data-driven surrogate modeling. Physics-informed neural networks (PINNs) enforce governing equations and boundary or initial conditions through the training loss \citep{raissi2019physics, pang2019fpinns}, while operator-learning methods learn mappings between function spaces across families of PDE instances \citep{kovachki2023neural}. Representative approaches include DeepONet \citep{lu2021learning}, graph-based and general neural operators \citep{li2020neural,li2020multipole,kovachki2023neural}, and the Fourier neural operator (FNO) \citep{li2020fourier}. Subsequent extensions consider multiple-input operators, deeper architectures, physics-informed training, alternative spectral representations, and complex geometries \citep{jin2022mionet,rahman2023uno,li2024physics,gupta2021multiwavelet,li2023geo}. Another prominent line develops attention-based operator architectures, including the Galerkin Transformer \citep{cao2021choose}, LOCA \citep{kissas2022learning}, OFormer \citep{li2022transformer}, GNOT \citep{hao2023gnot}, ONO \citep{xiao2023improved}, PiT \citep{chen2024positional}, Transolver \citep{wu2024transolver}, and LNO \citep{wang2024latent}. Our work builds on this operator-learning paradigm and focuses on stable temporal evolution.

\paragraph{Learning physical dynamics.} A related line of work focuses explicitly on learning the temporal evolution of physical systems. Autoregressive simulators commonly learn a local update rule and repeatedly apply it to advance the physical state, including graph-based physical simulators, mesh-based models, neural PDE solvers, and convolutional encoder-decoder surrogates \citep{sanchez2020learning,pfaff2020learning,brandstetter2022message,stachenfeld2021learned,geneva2020modeling}. Since repeated prediction can amplify local errors and shift the model away from the state distribution encountered during training, explored perturbing training states, multi-step or curriculum training, and adaptive temporal stepping to improve rollout robustness \citep{sanchez2020learning,li2022transformer,hagnberger2025calm, wu2026tante}. Another prominent direction evolves the system in a learned latent space, using recurrent models, continuous latent dynamics, or Koopman-inspired representations \citep{wiewel2019latent,lusch2018deep,morton2018deep,pan2020physics,yin2023continuous}. Along this direction, \cite{geneva2022transformers} further combines Koopman-informed embeddings with Transformers for temporal prediction, while OFormer \citep{li2022transformer} performs recurrent time-marching on spatially resolved latent representations. More recent approaches such as LNO and CALM-PDE similarly model time-dependent PDEs through compressed latent representations \citep{wang2024latent,hagnberger2025calm}. Our work is most closely related to this latent-dynamics perspective, but differs in formulating the latent evolution variationally.

\paragraph{Probabilistic and variational modeling.} Probabilistic PDE surrogates have been studied primarily for uncertainty quantification. Early work includes Bayesian convolutional encoder-decoder models for stochastic PDEs and physics-constrained generative surrogates \citep{zhu2018bayesian,zhu2019physics}. Similar ideas have subsequently been extended to operator learning through Bayesian, probabilistic, and variational DeepONet formulations \citep{yang2022scalable,lin2023bdeeponet,garg2023vbdeeponet}. More closely related to our formulation, Generative Adversarial Neural Operators (GANO) learn distributions directly on function spaces using neural operators \citep{rahman2022generative}, while Variational Autoencoding Neural Operators (VANO) extend variational autoencoders to functional data \citep{seidman2023variational}. Autoencoders have also been formulated directly on function spaces, with variational extensions that address the well-posedness of infinite-dimensional objectives \citep{bunker2025autoencoders}. More recently, probabilistic generative models have been explored for PDE modeling and temporal prediction. Physics-Informed Diffusion Models incorporate PDE constraints into diffusion training \citep{bastek2025physics}, while PDE-Refiner uses diffusion-inspired denoising for long autoregressive rollouts \citep{lippe2023pde}, and AROMA employs diffusion-based latent dynamics for temporal prediction \citep{serrano2024latent,serrano2024aroma}. Our work complements them by developing a variational Markov formulation for latent PDE evolution and studying how its distributional training objective relates to autoregressive rollout.



\section{Problem Formulation}
\label{sec:problem_formulation}

We consider a class of time-dependent partial differential equations (PDEs) defined on a spatial domain $\Omega \subset \mathbb{R}^{d_x}$. Let
\[
    \mathbf{U}(t): \Omega \to \mathbb{R}^{c}
\]
denote the physical state of the system at time $t$, where $c$ is the number of physical channels. Depending on the governing equation, $\mathbf{U}$ may represent a scalar field, such as the vorticity field in two-dimensional Navier-Stokes equations, or a vector-valued field, such as the density,
velocity, and pressure variables in compressible flow.

We abstract the governing dynamics as an autonomous evolution equation
\begin{equation}
    \partial_t \mathbf{U}(t)
    =
    \mathcal{F}(\mathbf{U}(t)),
    \label{eq:abstract_pde}
\end{equation}
where $\mathcal{F}$ is a generally nonlinear differential operator. In concrete
PDEs, $\mathcal{F}$ may depend on spatial derivatives of $\mathbf{U}$, physical
parameters, external forcing fields, and boundary conditions. For example,
one may write more explicitly
\[
    \partial_t \mathbf{U}
    =
    \mathcal{F}
    \bigl(
        \mathbf{U},
        \nabla \mathbf{U},
        \nabla^2 \mathbf{U};
        \lambda, f, \Omega, \mathcal{B}
    \bigr),
\]
where $\lambda$ denotes physical parameters, $f$ denotes a forcing term, and $\mathcal{B}$ denotes the boundary condition. Throughout this work, we absorb these fixed problem-dependent quantities into the operator $\mathcal{F}$ and use the compact notation in \eqref{eq:abstract_pde}.

The continuous-time PDE induces a solution operator, or flow map,
\[
    \mathcal{S}_{h}: \mathbf{U}(t) \mapsto \mathbf{U}(t+h),
\]
which advances the physical state by a time interval $h$. Given a temporal discretization $t_n = nh$, we denote
\[
    \mathbf{U}_n := \mathbf{U}(t_n).
\]
The corresponding discrete-time dynamics are
\begin{equation}
    \mathbf{U}_{n+1}=\mathcal{S}_{h}(\mathbf{U}_n).\label{eq:flow_map}
\end{equation}
In data-driven PDE modeling, the goal is to learn an approximation to this flow map from observed trajectories $\{\mathbf{U}_0,\mathbf{U}_1,\ldots,\mathbf{U}_T\}$.

\subsection{Variational Autoencoders}
\label{sec:vae_review}

Here we briefly review the variational autoencoder (VAE) framework \citep{kingma2013auto, seidman2023variational} in a form that will be used throughout the paper. Let $\mathcal U$ denote the data space and $\mathcal Z$ the latent space. We denote a data sample by $\mathbf U\in\mathcal U$ and its latent representation by $\mathbf Z\in\mathcal Z$. A variational autoencoder consists of two main probabilistic components. The \emph{encoder} $Q_\phi$ defines a conditional distribution
\[
    \mathcal{U}\ni\mathbf U\,\mapsto\, Q_\phi(\d\mathbf Z\,|\,\mathbf U)\in\Delta(\mathcal Z)
\]
over latent representations $\mathbf Z\in\mathcal Z$ given an observation $\mathbf U\in\mathcal U$. Similarly, the \emph{decoder} $P_\psi$ is a conditional kernel
\[
    \mathcal{Z}\ni\mathbf Z\,\mapsto\, P_\psi(\d\mathbf U\,|\,\mathbf{Z})\in\Delta(\mathcal U)
\]
over observations conditioned on the latent representation. Here $\phi\in\Phi$ and $\psi\in\Psi$ denote the encoder and decoder parameters, respectively.

The generative model is completed by prescribing a reference latent distribution $\mathsf P_{\mathcal Z}\in\Delta(\mathcal Z)$. A latent sample is first drawn according to $\mathbf Z\sim\mathsf P_{\mathcal Z}$, and the corresponding observation is generated from $\mathbf U\sim P_\psi(\cdot\mid\mathbf Z)$. The encoder $Q_\phi(\cdot\mid\mathbf U)$ provides a tractable variational approximation to the generally intractable posterior distribution of $\mathbf Z$ conditioned on $\mathbf U$.

To train the model, one maximizes a variational lower bound on the data likelihood. Let $\mathsf W_{\mathcal U}$ be a reference measure on $\mathcal U$, and suppose that the decoder likelihood is dominated by $\mathsf W_{\mathcal U}$. Under the usual absolute-continuity conditions, the evidence lower bound (ELBO) takes the form
\begin{equation}
\begin{aligned}
    \log
    \frac{\d P_\psi^{\,\mathcal U}}
         {\d\mathsf W_{\mathcal U}}(\mathbf U)
    \geq
    &\;
    \mathbb E_{\mathbf Z\sim Q_\phi(\cdot\mid\mathbf U)}
    \left[
        \log
        \frac{\d P_\psi(\cdot\mid\mathbf Z)}
             {\d\mathsf W_{\mathcal U}}(\mathbf U)
    \right]-
    D_{\mathrm{KL}}
    \left(
        Q_\phi(\cdot\mid\mathbf U)
        \,\Vert\,
        \mathsf P_{\mathcal Z}
    \right),
\end{aligned}
\label{eq:vae_elbo}
\end{equation}
where
\[
P_\psi^{\,\mathcal U}(\d\mathbf U)=\int_{\mathcal Z} P_\psi(\d\mathbf U\mid\mathbf Z)\,\mathsf P_{\mathcal Z}(\d\mathbf Z)
\]
is the marginal distribution induced by the generative model.

The two terms in \eqref{eq:vae_elbo} play complementary roles. The first encourages latent samples produced by the encoder to retain sufficient information for reconstruction through the decoder, while the KL term regularizes the encoded distribution toward the prescribed latent prior. Thus, a variational autoencoder combines probabilistic autoencoding with distributional regularization of the latent representation.

\paragraph{Notations.}


\section{Variational Latent Dynamics on Function Spaces}\label{sec:var_latent_dyn}
In this section, we develop the theoretical foundation of variational latent dynamics for time-dependent PDEs. We first formulate the model abstractly on physical and latent function spaces and characterize the variational objectives governing latent representation and transition alignment in \S\ref{sec:abstract_vld}. Then in \S\ref{sec:functional_gaussian_latent_markov}, we specialize to functional Gaussian distributions, where Cameron-Martin theory provides tractable likelihood representations and the covariance structure induces a spectral geometry on latent perturbations and residuals. Finally, in \S\ref{sec:err_var_dyn}, we analyze long-horizon error propagation under deterministic mean rollout, highlighting two complementary mechanisms of variational training: noise injection promotes local stability, while variational alignment provides direct control of latent transition errors and their accumulation over the rollout horizon.

\subsection{Abstract Variational Latent Dynamics}\label{sec:abstract_vld}
We formulate the latent dynamics directly at the level of function spaces, before introducing a particular Gaussian parameterization or spatial discretization. Let $\mathcal U$ denote the space of physical PDE states and $\mathcal Z$ the corresponding latent state space.
\begin{assumption}[State spaces]
\label{assump:state_spaces}
The physical state space $\mathcal U$ and latent state space $\mathcal Z$ are real separable Banach spaces, equipped with their Borel $\sigma$-algebras.
\end{assumption}
This setting is sufficiently general to include many state spaces arising in PDE modeling, while providing the standard measurable structure required to define conditional probability distributions on $\mathcal U$ and $\mathcal Z$. In particular, separability ensures that these spaces have the regularity needed for the probabilistic constructions below.

Built on this setting, we study a variational latent formulation of the PDE flow map. Instead of directly learning a deterministic propagator in the physical space, we introduce a latent representation $\mathbf{Z}_n\in \mathcal Z$ for each physical state $\mathbf{U}_n\in \mathcal U$ and model the one-step evolution through the latent pathway
\begin{equation*}
    \mathbf{U}_n\xrightarrow{\text{encode}}
    \mathbf{Z}_n\xrightarrow{\text{evolve}}
    \mathbf{Z}_{n+1}\xrightarrow{\text{decode}}
    \mathbf{U}_{n+1}.
\end{equation*}
Along this pathway, an encoder maps the physical field to a latent representation, a latent transition model propagates this representation in time, and a decoder maps the propagated latent state back to the physical space. This formulation provides a probabilistic description of the learned dynamics while allowing the evolution to be regularized in a structured latent space.

Concretely, for the pathway above, we introduce encoder, transition and decoder kernels, parameterized by $\phi\in\Phi,\theta\in\Theta$ and $\psi\in\Psi$, respectively:
\begin{equation*}
\begin{aligned}
\text{(encoder)}\qquad &\mathcal{U}\ni\mathbf{U}_n\,\mapsto \,Q_\phi(\d\mathbf{Z}_n \,|\, \mathbf{U}_n)\in\Delta\left(\mathcal{Z}\right),\\
\text{(transition)}\qquad &\mathcal{Z}\ni\mathbf{Z}_n\,\mapsto\, P_{\,\theta}(\d\mathbf{Z}_{n+1} \,|\, \mathbf{Z}_n)\in\Delta(\mathcal{Z}),\\
\text{(decoder)}\qquad &\mathcal{Z}\ni\mathbf{Z}_n\,\mapsto \,P_\psi(\d\mathbf{U}_{n+1}\,|\, \mathbf{Z}_{n+1})\in\Delta(\mathcal{U}).
\end{aligned}
\end{equation*}
Together, these components define a probabilistic approximation of the one-step PDE flow map and yield a variational latent Markov model for time-dependent PDEs.

The formulation is independent of the particular physical variables used to represent the underlying PDE state. For example, $\mathbf U_n$ may denote the vorticity field in incompressible Navier-Stokes or the density, velocity, and pressure fields in compressible Euler. The same latent probabilistic framework therefore applies across different PDE systems through the corresponding physical-state space and data distribution.

Before specializing the model to functional Gaussian kernels, we first establish the variational foundation of the latent Markov formulation at the level of general function-space probability measures. We begin by deriving lower bounds on the one-step predictive likelihood, including a conditional functional ELBO that motivates the reconstruction and latent transition consistency terms. We then study the dynamics term at the population level and characterize the latent Markov transition targeted by its minimization. These results do not rely on any particular parameterization and therefore apply to the general formulation above.



\begin{proposition}[Functional lower bounds]
\label{prop:conditional_functional_elbo}
Let $\mathcal U$ and $\mathcal Z$ be Banach spaces. Let
$\mathsf W_{\mathcal U}$ be a fixed reference measure on $\mathcal U$. For a consecutive pair $(\mathbf U_n,\mathbf U_{n+1})\in\mathcal{U}\times\mathcal{U}$, suppose that the encoder, latent transition, and decoder are given by Markov kernels
\[
    Q_\phi(\d\mathbf Z_n\,|\, \mathbf U_n)\in\Delta(\mathcal{Z}),\quad
    P_\theta(\d\mathbf Z_{n+1}\,|\, \mathbf Z_n)\in\Delta(\mathcal{Z}),\quad
    P_\psi(\d\mathbf U_{n+1}\,|\, \mathbf Z_{n+1})\in\Delta(\mathcal{U}).
\]
Assume that the decoder kernel is absolutely continuous with respect to $\mathsf W_{\mathcal U}$:
\[
    P_\psi(\cdot\mid \mathbf Z)\ll \mathsf W_{\mathcal U}.
\]
Then the following lower bounds hold.
\begin{itemize}
\item[(a)] (Predictive decoder likelihood bound).
\begin{equation}
    \log \frac{\d P_{\theta,\psi,\phi}(\cdot\,|\,\mathbf{U}_n)}{\d\mathsf{W}_\mathcal{U}}(\mathbf{U}_{n+1})\geq\bbE_{\mathbf{Z}_{n+1}\sim P_\theta(\cdot|\mathbf{Z}_n),\mathbf{Z}_n\sim Q_\phi(\cdot|\mathbf{U}_n)}\left[\log \frac{\d P_\psi(\cdot\,|\,\mathbf{Z}_{n+1})}{\d\mathsf{W}_\mathcal{U}}(\mathbf{U}_{n+1})\right].\label{eq:predlb}
\end{equation}
\item[(b)] (Conditional functional ELBO). Assume that the relevant likelihood densities exist and that
\[
    Q_\phi(\cdot\,|\, \mathbf U_{n+1})
    \ll
    P_\theta(\cdot\,|\, \mathbf Z_n)
\]
for $Q_\phi(\cdot\,|\, \mathbf U_n)$-almost every $\mathbf Z_n$. Then
\[
\begin{aligned}
    \log \frac{\d P_{\theta,\psi,\phi}(\cdot\,|\,\mathbf{U}_n)}{\d\mathsf{W}_\mathcal{U}}(\mathbf{U}_{n+1})
    &\geq\mathbb{E}_{\mathbf Z_{n+1}\sim Q_\phi(\cdot\,|\, \mathbf U_{n+1})}
    \left[
        \log\frac{\d P_\psi(\cdot\,|\, \mathbf Z_{n+1})}{\d\mathsf{W}_\mathcal{U}}(\mathbf U_{n+1})
    \right]
    \\
    &\quad -\mathbb E_{\mathbf Z_n\sim Q_\phi(\cdot\,|\, \mathbf U_n)}
    \left[D_{\mathrm{KL}}\left(
    Q_\phi(\cdot\,|\, \mathbf U_{n+1})\,\Vert\,P_\theta(\cdot\,|\, \mathbf Z_n)
    \right)\right].
\end{aligned}
\]
\end{itemize}
\end{proposition}

\begin{remark}
Equivalently, the first bound gives the predictive negative log-likelihood objective
\begin{equation}
    \mathcal L_{\mathrm{NLL}}^{\mathrm{pred}}(\mathbf U_n,\mathbf U_{n+1})=\bbE_{\mathbf{Z}_{n+1}\sim P_\theta(\cdot|\mathbf{Z}_n),\mathbf{Z}_n\sim P_\psi(\cdot|\mathbf{U}_n)}\bigl[-\log P_\psi(\mathbf{U}_{n+1}\,|\,\mathbf{Z}_{n+1})\bigr].\label{eq:pred_nll}
\end{equation}
Similarly, the negative conditional ELBO is
\begin{equation}
\begin{aligned}
    \mathcal L_{\mathrm{ELBO}}(\mathbf U_n,\mathbf U_{n+1})
    =
    &
    -
    \mathbb E_{\mathbf Z_{n+1}\sim Q_\phi(\cdot\,|\, \mathbf U_{n+1})}
    \big[
        \log P_\psi(\mathbf U_{n+1}\,|\, \mathbf Z_{n+1})
    \big]
    \\
    &+
    \mathbb E_{\mathbf Z_n\sim Q_\phi(\cdot\,|\, \mathbf U_n)}
    \left[
    D_{\mathrm{KL}}
    \left(
        Q_\phi(\cdot\,|\, \mathbf U_{n+1})
        \,\Vert\,
        P_\theta(\cdot\,|\, \mathbf Z_n)
    \right)
    \right].
\end{aligned}\label{eq:condnegelbo}
\end{equation}
Notably, the ELBO can be understood in the extended-real sense. If $Q_\phi(\cdot\mid \mathbf U_{n+1})$ is not absolutely continuous with respect to $P_\theta(\cdot\mid\mathbf Z_n)$, the corresponding KL divergence is infinite and the lower bound becomes trivial. A finite, nontrivial ELBO therefore requires
\[
    Q_\phi(\cdot\mid \mathbf U_{n+1})
    \ll
    P_\theta(\cdot\mid \mathbf Z_n)
\]
for $Q_\phi(\cdot\mid \mathbf U_n)$-almost every $\mathbf Z_n$.
\end{remark}

The conditional ELBO in \eqref{eq:condnegelbo} separates two complementary learning objectives. The reconstruction term
\[
    -
    \mathbb E_{\mathbf Z_{n+1}\sim Q_\phi(\cdot\,|\,\mathbf U_{n+1})}
    \big[\log P_\psi(\mathbf U_{n+1}\,|\,\mathbf Z_{n+1})\big],
\]
encourages latent samples drawn from $Q_\phi(\cdot\mid\mathbf U_{n+1})$ to retain the information needed to reconstruct the physical state $\mathbf U_{n+1}$. It therefore promotes consistency between the encoder and decoder and ensures that the latent variables remain informative representations of the underlying PDE states. The second term,
\[
    \mathbb E_{\mathbf Z_n\sim Q_\phi(\cdot\,|\,\mathbf U_n)}
    \left[D_{\mathrm{KL}}
    \left(
        Q_\phi(\cdot\,|\,\mathbf U_{n+1})\,\Vert\,P_\theta(\cdot\,|\,\mathbf Z_n)
    \right)\right],
\]
enforces consistency of the latent dynamics. Given a latent state $\mathbf Z_n$ encoded from $\mathbf U_n$, the transition kernel $P_\theta(\cdot\mid\mathbf Z_n)$ predicts a distribution over the next latent state. The KL divergence aligns this prediction with the posterior distribution obtained by encoding the true next state $\mathbf U_{n+1}$. Thus, the variational objective couples two complementary requirements:
\[
    \text{physical reconstruction}\quad+\quad\text{latent transition alignment}.
\]
The preceding ELBO gives a sample-level variational objective. We next characterize its latent dynamics term at the population level and identify the transition kernel targeted by its minimization.


\begin{proposition}[Population decomposition of the ELBO dynamics term]
\label{prop:population_elbo_transition_consistency}
Let $\rho$ be a probability measure on consecutive PDE states $(\mathbf U,\mathbf U^+)\in\mathcal U\times\mathcal U$, and let $Q_\phi(\cdot\mid\mathbf U)$ be an encoder kernel. For a transition kernel $P$ from $\mathcal Z$ to $\mathcal Z$, define the population dynamics loss
\[
\mathcal R(P)=\mathbb E_{(\mathbf U,\mathbf U^+)\sim\rho} \mathbb E_{\mathbf Z\sim Q_\phi(\cdot\mid\mathbf U)} \left[D_{\mathrm{KL}}\left(Q_\phi(\cdot\mid\mathbf U^+)
\,\Vert\,P(\cdot\mid\mathbf Z)\right)\right].
\]
Let $\Gamma_\phi$ be the joint law of $(\mathbf{U},\mathbf{U}^+,\mathbf{Z},\mathbf{Z}^+)$ induced by
\[
    (\mathbf U,\mathbf U^+)\sim\rho,\qquad
    \mathbf Z\sim Q_\phi(\cdot\mid\mathbf U),\qquad
    \mathbf Z^+\sim Q_\phi(\cdot\mid\mathbf U^+),
\]
and let $\Pi_\phi(\cdot\mid\mathbf Z)$ be the conditional law of
$\mathbf Z^+$ given $\mathbf Z$ under $\Gamma_\phi$. Then
\begin{equation}
\mathcal R(P)=\bbE_{(\mathbf{Z},\mathbf{U}^+)\sim\Gamma_\phi}\left[D_{\mathrm{KL}}
    \left(Q_\phi(\cdot\,|\,\mathbf U^+)\,\Vert\,\Pi_\phi(\cdot\,|\,\mathbf Z)
    \right)\right]+
\bbE_{\mathbf{Z}\sim\Gamma_\phi}\left[D_{\mathrm{KL}}
    \left(\Pi_\phi(\cdot\,|\,\mathbf Z)\,\Vert\,P(\cdot\,|\,\mathbf Z)
    \right)\right].\label{eq:elbokldecomp}
\end{equation}
Consequently, among all Markov kernels
$P$, the population dynamics loss is minimized by
\[
    P^\star(\cdot\,|\,\mathbf Z)=\Pi_\phi(\cdot\,|\,\mathbf Z)
\]
for $\Gamma_\phi$-almost every $\mathbf Z$. 
\end{proposition}

Proposition~\ref{prop:population_elbo_transition_consistency} separates the population dynamics loss into an encoder-induced term and a transition-model mismatch:
\begin{itemize}
\item The first term in \eqref{eq:elbokldecomp} depends only on the encoder and the distribution of consecutive physical states. It quantifies the dispersion of the encoded next-state posterior $Q_\phi(\cdot\mid\mathbf U^+)$ around the aggregated conditional latent transition $\Pi_\phi(\cdot\mid\mathbf Z)$ and is therefore irreducible when optimizing over the transition kernel $P$.

\item The second term is the only component that depends on $P$. It measures the discrepancy between the learned Markov transition $P(\cdot\mid\mathbf Z)$ and the encoder-induced conditional transition $\Pi_\phi(\cdot\mid\mathbf Z)$. Hence, for a fixed encoder, minimizing the population dynamics objective is equivalent to matching the learned transition to $\Pi_\phi$ almost surely under the latent marginal distribution. In particular, if the model class $\{P_\theta\}_{\theta\in\Theta}$ contains
$\Pi_\phi$, any global population minimizer satisfies
\[
    P_{\theta^\star}(\cdot\mid\mathbf Z)
    =
    \Pi_\phi(\cdot\mid\mathbf Z)
\]
up to $\Gamma_\phi$-null sets. Therefore, rather than identifying the original PDE generator $\mathcal F$ directly, the ELBO dynamics term encourages the model to learn the Markov transition induced by the encoder on latent representations of true PDE trajectories.
\end{itemize}

The preceding results characterize the variational latent dynamics without imposing a specific form on the underlying kernels. We now specialize this framework to
functional Gaussian kernels and study the covariance-induced geometry of the latent dynamics through their Cameron-Martin spaces and spectral structure.

\subsection{Functional Gaussian Latent Markov Model}
\label{sec:functional_gaussian_latent_markov}

We now specialize the abstract variational formulation to Gaussian distributions on function spaces. This specialization serves two purposes. First, the Cameron-Martin theorem converts the Radon-Nikodym derivatives in the functional variational objective into tractable reconstruction and transition terms. Second, the covariance operators induce a geometry on the latent space that determines how perturbations and transition residuals are weighted across different functional directions. 

We briefly review the necessary background on Gaussian measures and Cameron-Martin spaces in Appendix~\S\ref{app:gaussian_measures_banach}. Recall that a functional Gaussian measure $\mathsf N(m,K)$ is characterized by its mean element $m$ and covariance operator $K$. Accordingly, we consider the functional Gaussian parameterization
\begin{equation*}
\begin{aligned}
    \text{(encoder)}\qquad Q_\phi(\cdot\mid\mathbf U)
    &=
    \mathsf N
    \bigl(
        m_\phi(\mathbf U),
        K_\phi(\mathbf U)
    \bigr),
    \\
    \text{(transition)}\ \qquad P_\theta(\cdot\mid\mathbf Z)
    &=
    \mathsf N
    \bigl(
        T_\theta(\mathbf Z),
        \alpha_\theta(\mathbf Z)^2 K_\theta
    \bigr),
    \\
    \text{(decoder)\,}\qquad P_\psi(\cdot\mid\mathbf Z)
    &=
    \mathsf N
    \bigl(
        D_\psi(\mathbf Z),
        K_\psi
    \bigr).
\end{aligned}
\end{equation*}
Here $m_\phi$ denotes the encoder mean, $T_\theta$ the mean latent transition, and $D_\psi$ the decoder mean. The state-dependent covariance $K_\phi(\mathbf U)$ describes uncertainty in the encoded representation, while $K$ specifies the spatial structure of the latent transition noise and $\alpha_\theta(\mathbf Z)>0$ controls its amplitude. The decoder covariance $K_\psi$ is fixed.

These Gaussian distributions are defined on the underlying function spaces rather than on a particular spatial discretization. The finite-dimensional parameterizations used in computation are introduced later in \S\ref{subsec:structured_latent_perturbations}. Here we focus on the function-space structure induced by the corresponding Gaussian measures.

\paragraph{Cameron-Martin representation of the likelihood.}
The likelihood terms in the functional variational bounds of Proposition~\ref{prop:conditional_functional_elbo} are expressed through Radon-Nikodym derivatives with respect to reference measures. For Gaussian kernels, the Cameron-Martin theorem gives an explicit representation of these derivatives in terms of the geometry induced by the covariance operator. We first record this consequence for the decoder likelihood.

\begin{proposition}[Cameron-Martin form of the decoder likelihood]
\label{prop:cm_decoder_likelihood}
Let $\mathcal U$ be a separable Banach space and let $\mathsf W_{\mathcal U}$ be a centered Gaussian measure on $\mathcal U$. Denote its Cameron-Martin space by $\mathcal H_{\mathcal U}$, continuously embedded in $\mathcal U$. Suppose the decoder kernel is given by the
Cameron-Martin shift
\[
    P_\psi(\d\mathbf U\mid \mathbf Z)
    =
    \mathsf W_{\mathcal U}^{D_\psi(\mathbf Z)}(\d\mathbf U),
\]
where $D_\psi(\mathbf Z)\in \mathcal H_{\mathcal U}$.
Then
\[
    P_\psi(\cdot\mid \mathbf Z)\ll \mathsf W_{\mathcal U},
\]
and the Radon-Nikodym derivative is
\begin{equation}
\frac{\d P_\psi(\cdot\mid \mathbf Z)}{\d\mathsf W_{\mathcal U}
}(\mathbf U)
=\exp\left(\langle D_\psi(\mathbf Z),\mathbf U\rangle^\sim
-\frac12\|D_\psi(\mathbf Z)\|_{\mathcal H_{\mathcal U}}^2
\right),\label{eq:cm_decoder_rn}
\end{equation}
where $\langle\cdot,\cdot\rangle^\sim$ denotes the Paley-Wiener map associated with $\mathsf W_{\mathcal U}$. Consequently, the decoder negative log-likelihood with respect to
$\mathsf W_{\mathcal U}$ is
\begin{equation}
-\log\frac{\d P_\psi(\cdot\mid \mathbf Z)}{\d\mathsf W_{\mathcal U}}(\mathbf U)
    =
    \frac12
    \|D_\psi(\mathbf Z)\|_{\mathcal H_{\mathcal U}}^2
    -
    \langle D_\psi(\mathbf Z),\mathbf U\rangle^\sim.
    \label{eq:cm_decoder_nll}
\end{equation}
In particular, if $\mathbf U\in \mathcal H_{\mathcal U}$, then
\[
    \langle D_\psi(\mathbf Z),\mathbf U\rangle^\sim=
    \langle D_\psi(\mathbf Z),\mathbf U\rangle_{\mathcal H_{\mathcal U}},
\]
and therefore
\begin{equation}
-\log\frac{\d P_\psi(\cdot\mid \mathbf Z)}{\d\mathsf W_{\mathcal U}}(\mathbf U)
=
\frac12\|\mathbf U-D_\psi(\mathbf Z)\|_{\mathcal H_{\mathcal U}}^2-\frac12\|\mathbf U\|_{\mathcal H_{\mathcal U}}^2.
\label{eq:cm_decoder_squared_norm}
\end{equation}
Thus, up to a term depending only on the observed field $\mathbf U$, the Gaussian decoder negative log-likelihood is equivalent to the squared Cameron-Martin reconstruction error.
\end{proposition}

\begin{remark}[Deterministic decoder mean]
The probabilistic decoder is parameterized through the deterministic mean map $D_\psi$, while the observation uncertainty is represented by the fixed Gaussian reference measure. Proposition~\ref{prop:cm_decoder_likelihood}
therefore connects the functional Gaussian likelihood to the deterministic decoder used in computation. In particular, after finite-dimensional discretization with isotropic covariance, the Cameron-Martin reconstruction
term reduces to the usual squared Euclidean  reconstruction loss.
\end{remark}

\paragraph{Spectral geometry of the latent transition.}
The same Gaussian structure also determines how latent transition errors are measured. We now assume that $\mathcal Z$ is a separable Hilbert space, so that the covariance operator admits a spectral decomposition. The resulting Cameron-Martin norm reveals how the covariance spectrum assigns different weights to different latent directions.

\begin{proposition}[Spectral geometry of the latent transition]
\label{prop:general_spectral_residual_control}
Let $\mathcal Z$ be a separable Hilbert space and let linear operator $K:\mathcal Z\to\mathcal Z$ be self-adjoint, positive-definite, and trace-class. Let $\{(\lambda_i,\mathbf E_i)\}_{i\geq1}$ be an eigensystem of $K$, where
$\{\mathbf E_i\}_{i\geq1}$ is an orthonormal basis of $\mathcal Z$, eigenvalues $\lambda_i>0$, and $\sum_{i=1}^{\infty}\lambda_i<\infty$. Fix
\[
    \mathsf W_{\mathcal Z}
    =
    \mathsf N(0,K)
\]
and consider the latent transition
\[
    P_\theta(\cdot\mid\mathbf Z)
    =
    \mathsf N
    \bigl(
        T_\theta(\mathbf Z),
        \alpha_\theta(\mathbf Z)^2K
    \bigr),
\]
where $T_\theta(\mathbf Z)\in\mathcal H_K$ and
$\alpha_\theta(\mathbf Z)>0$. Then:

\begin{itemize}[topsep=0pt,itemsep=2pt]
\item[(a)]
The Cameron-Martin space associated with $\mathsf W_{\mathcal Z}$ is
\[
    \mathcal H_K
    =
    \left\{
        \mathbf R=\sum_{i=1}^{\infty}r_i\mathbf E_i:
        \sum_{i=1}^{\infty}\frac{|r_i|^2}{\lambda_i}<\infty
    \right\},
\]
with norm
\begin{equation}
    \|\mathbf R\|_{\mathcal H_K}^2
    =
    \sum_{i=1}^{\infty}
    \frac{|r_i|^2}{\lambda_i}.
    \label{eq:cm_spectral_norm}
\end{equation}
\item[(b)]
The transition admits the Karhunen--Lo\`eve representation
\begin{equation}
    \mathbf Z^+
    =
    T_\theta(\mathbf Z)
    +
    \alpha_\theta(\mathbf Z)
    \sum_{i=1}^{\infty}
        \sqrt{\lambda_i}\,\xi_i\mathbf E_i,
    \qquad
    \xi_i\overset{\mathrm{i.i.d.}}{\sim}\mathsf N(0,1),
    \label{eq:transition_kl_expansion}
\end{equation}
where the series converges in $L^2(\Omega;\mathcal Z)$ and almost surely in $\mathcal Z$. Moreover, for $\mathbf Z^+\in\mathcal H_K$,
\begin{equation}
\begin{aligned}
    -\log\frac{\d P_\theta(\cdot\mid\mathbf Z)}{\d\mathsf N(0,\alpha_\theta(\mathbf Z)^2K)}(\mathbf Z^+)
    &=\frac{1}{2\alpha_\theta(\mathbf Z)^2}\|\mathbf Z^+ -T_\theta(\mathbf Z)\|_{\mathcal H_K}^2 -\frac{1}{2\alpha_\theta(\mathbf Z)^2}\|\mathbf Z^+\|_{\mathcal H_K}^2 .
\end{aligned}
\label{eq:transition_cm_density}
\end{equation}
\end{itemize}
\end{proposition}

\begin{remark}[Spectral-coordinate interpretation]
Equivalently, writing
\[
    T_\theta(\mathbf Z)=\sum_{i=1}^{\infty}t_i(\mathbf Z)\mathbf E_i,\qquad\mathbf Z^+=\sum_{i=1}^{\infty}z_i^+\mathbf E_i,
\]
the representation \eqref{eq:transition_kl_expansion} gives
\[
    z_i^+=t_i(\mathbf Z)+\alpha_\theta(\mathbf Z)\sqrt{\lambda_i}\,\xi_i,\qquad
    \operatorname{Var}(z_i^+\mid\mathbf Z)=\alpha_\theta(\mathbf Z)^2\lambda_i.
\]
Hence the covariance spectrum simultaneously determines the directions and scales of the injected transition noise and the geometry of the corresponding Cameron-Martin residual. Directions with smaller $\lambda_i$ receive less stochastic variation and a larger weight $1/\lambda_i$ in the residual norm, whereas directions with larger $\lambda_i$ are perturbed more strongly and penalized less.
\end{remark}

A particularly relevant specialization arises when the latent space is a periodic function space and the covariance operator is diagonal in the Fourier basis. In this setting, the spectral weighting induced by the Cameron-Martin norm admits a direct Sobolev interpretation.

\begin{corollary}[Sobolev geometry of latent residuals]
\label{cor:sobolev_residual_control}
Suppose
\[
    \mathcal Z=L^2(\mathbb T^{d_x};\mathbb R^{d_z})
\]
and $K$ is diagonal in the Fourier basis with eigenvalues
\[
    \lambda_k=(1+\ell^2|k|^2)^{-s}.
\]
Then
\begin{equation}
    \|\mathbf Z^+-T_\theta(\mathbf Z)\|_{\mathcal H_K}^2
    =\sum_{k,a}(1+\ell^2|k|^2)^s
    \left|\widehat{\mathbf Z}^+_a(k)-\widehat{T_\theta(\mathbf Z)}_a(k)\right|^2.
    \label{eq:sobolev_latent_residual}
\end{equation}
Thus the Cameron-Martin geometry induced by $K$ is equivalent to a length-scale-weighted Sobolev geometry of order $s$. In particular, high-frequency latent residuals receive increasingly large weight as the covariance spectrum decays.
\end{corollary}


\subsection{Error Propagation in Variational Latent Dynamics}
\label{sec:err_var_dyn}
We now study how the variational training objective affects the stability of long-horizon latent rollout. We first make explicit the one-step objective analyzed below. Let $\mu_n$ denote the joint distribution of consecutive physical states $(\mathbf U_n,\mathbf U_{n+1})$, and consider the Gaussian encoder and transition
\[
    \widetilde{\mathbf Z}_n=m_\phi(\mathbf U_n)+\Xi_n,
    \qquad
    \Xi\sim\mathsf N(0,K_\phi(\mathbf U_n)),
\]
and
\[
    \widetilde{\mathbf Z}_n^+=T_\theta(\widetilde{\mathbf Z}_n)+\alpha_\theta(\widetilde{\mathbf Z}_n)\mathbf G,
    \qquad
    \mathbf G_n\sim\mathsf N(0,K).
\]
For notational convenience, define the decoder negative log-likelihood
\[
    \ell_\psi(\mathbf U;\mathbf Z)
    := -\log\frac{\d P_\psi(\cdot\mid\mathbf Z)}{\d\mathsf W_{\mathcal U}}(\mathbf U).
\]
For a consecutive pair $(\mathbf U_n,\mathbf U_{n+1})$, define
\begin{align*}
    \mathcal L_{\mathrm{pred}}(\mathbf U_n,\mathbf U_{n+1})
    &:=\mathbb E_{\Xi_n,\mathbf G_n}
    \left[\ell_\psi(\mathbf U_{n+1};\widetilde{\mathbf Z}_n^+)\right],
    \\
    \mathcal L_{\mathrm{KL}}(\mathbf U_n,\mathbf U_{n+1})
    &:=\mathbb E_{\widetilde{\mathbf Z}_n\sim
        Q_\phi(\cdot\mid\mathbf U_n)}
    \left[D_{\mathrm{KL}}\left(
        Q_\phi(\cdot\mid\mathbf U_{n+1})\,\Vert\,P_\theta(\cdot\mid\widetilde{\mathbf Z}_n)
    \right)\right],
    \\
    \mathcal L_{\mathrm{rec}}(\mathbf U_n)
    &:=\mathbb E_{\widetilde{\mathbf Z}_n\sim Q_\phi(\cdot\mid\mathbf U_n)}
    \left[
        \ell_\psi(\mathbf U_n;\widetilde{\mathbf Z}_n)
    \right].
\end{align*}
The population training objective takes the form
\begin{equation}
    \mathcal L^{(n)}(\phi,\theta,\psi)
    =\mathbb E_{(\mathbf U_n,\mathbf U_{n+1})\sim\mu_n}
    \left[
        \mathcal L_{\mathrm{pred}}(\mathbf U_n,\mathbf U_{n+1})
        +\beta\mathcal L_{\mathrm{KL}}(\mathbf U_n,\mathbf U_{n+1})
        +\lambda_{\mathrm{rec}}\mathcal L_{\mathrm{rec}}(\mathbf U_{n+1})
    \right],
    \label{eq:functional_training_objective}
\end{equation}
with the empirical objective obtained by averaging over observed one-step pairs.

By Proposition~\ref{prop:cm_decoder_likelihood}, whenever the corresponding physical state belongs to $\mathcal H_{\mathcal U}$,
\[
    \ell_\psi(\mathbf U;\mathbf Z)
    =\frac12
    \|\mathbf U-D_\psi(\mathbf Z)\|_{\mathcal H_{\mathcal U}}^2
    + c(\mathbf U),
\]
where $c(\mathbf U)$ is independent of the model parameters. Hence the decoder likelihood terms are equivalent, for optimization purposes, to squared Cameron-Martin reconstruction errors, up to additive data-dependent constants.

The objective \eqref{eq:functional_training_objective} can be viewed as a weighted combination of the two variational bounds introduced in Section~\ref{sec:abstract_vld}. The predictive term corresponds to the negative decoder-likelihood bound in
Proposition~\ref{prop:conditional_functional_elbo}(a), while the reconstruction and KL terms correspond to the two components of the negative conditional ELBO in
Proposition~\ref{prop:conditional_functional_elbo}(b). The coefficients $\beta$ and $\lambda_{\mathrm{rec}}$ allow for relative weighting of these terms, including the scaling induced by the fixed decoder covariance.

For convenience, our subsequent rollout analysis is stated instead in the ambient physical-state norm $\|\cdot\|_{\mathcal U}$. The two geometries are compatible because the Cameron-Martin space is continuously embedded in $\mathcal U$:
\[
    \mathcal H_{\mathcal U}\hookrightarrow\mathcal U,
    \qquad
    \|\mathbf v\|_{\mathcal U}
    \leq
    C_{\mathcal U}
    \|\mathbf v\|_{\mathcal H_{\mathcal U}},
    \quad
    \mathbf v\in\mathcal H_{\mathcal U},
\]
for some $C_{\mathcal U}>0$. Thus, whenever the relevant residual lies in $\mathcal H_{\mathcal U}$, likelihood control also yields control in the ambient norm used for rollout errors. After finite-dimensional discretization, the isotropic decoder covariance used in our implementation reduces these likelihood terms to the squared Euclidean losses described in \S\ref{subsec:vamo_training}.

\paragraph{Deterministic rollout.} At inference time, we deploy the deterministic mean dynamics. Given an initial physical state $\mathbf U_0$, the rollout is initialized and propagated according to
\begin{equation}
    \widehat{\mathbf Z}_0
    =
    m_\phi(\mathbf U_0),
    \quad
    \widehat{\mathbf Z}_n
    =
    T_\theta(\widehat{\mathbf Z}_{n-1}),
    \quad
    \widehat{\mathbf U}_n
    =
    D_\psi(\widehat{\mathbf Z}_n),\qquad n=1,2,\cdots,T/\Delta t.
    \label{eq:deterministic_mean_rollout}
\end{equation}
Hence the stochasticity in \eqref{eq:functional_training_objective} is used during training, whereas the deployed solver follows the deterministic mean latent dynamics. Our analysis below explains how this stochastic training objective can nevertheless control the error accumulated by the deterministic rollout.

We identify two complementary mechanisms. First, sampling from the encoder and transition distributions replaces pointwise fitting by training over Gaussian neighborhoods in latent space, providing a noise-injection regularization effect. Second, the variational KL term directly controls the mismatch between the learned transition and the encoded distribution of the true next state. We study these two mechanisms separately below.

\subsubsection{Noise Injection and Local Stability}
\label{subsubsec:noise_injection}

We first study the effect of random sampling from the encoder and transition distributions in the predictive pathway. Compared with a deterministic latent model, the stochastic objective evaluates the decoded transition over Gaussian neighborhoods of both the encoded input and the predicted latent. This allows the training objective to probe the local sensitivity of the learned latent dynamics and decoder.

Let $(\mathbf U_0^\star,\ldots,\mathbf U_N^\star)\sim \mu$ be a reference trajectory, and define its encoder-mean latent representation by
\[
    \mathbf Z_n^\star
    :=
    m_\phi(\mathbf U_n^\star).
\]
For each $n=0,\ldots,N-1$, let
\[
    \Xi_n
    \sim
    \mathsf N(0,K_\phi(\mathbf U_n^\star)),
    \qquad
    \mathbf G_n
    \sim
    \mathsf N(0,K),
\]
be conditionally independent given the reference trajectory, and define
\[
    \widetilde{\mathbf Z}_n
    :=
    \mathbf Z_n^\star+\Xi_n,
    \qquad
    \widetilde{\mathbf Z}_n^+
    :=
    T_\theta(\widetilde{\mathbf Z}_n)
    +
    \alpha_\theta(\widetilde{\mathbf Z}_n)\mathbf G_n.
\]

We consider three population-level quantities. The noisy one-step prediction error is
\begin{equation}
    \epsilon_n
    :=
    \left(
        \mathbb E
        \left[
            \left\|
                D_\psi(\widetilde{\mathbf Z}_n^+)
                -
                \mathbf U_{n+1}^\star
            \right\|_{\mathcal U}^2
        \right]
    \right)^{1/2},
    \label{eq:noisy_prediction_error}
\end{equation}
the encoder-noise sensitivity of the latent transition is
\begin{equation}
    \eta_n
    :=
    \left(
        \mathbb E
        \left[
            \left\|
                T_\theta(\widetilde{\mathbf Z}_n)
                -
                T_\theta(\mathbf Z_n^\star)
            \right\|_{\mathcal Z}^2
        \right]
    \right)^{1/2},
    \label{eq:encoder_noise_sensitivity}
\end{equation}
and the transition-noise decoder sensitivity is
\begin{equation}
    \rho_n
    :=
    \left(
        \mathbb E
        \left[
            \left\|
                D_\psi(\widetilde{\mathbf Z}_n^+)
                -
                D_\psi(T_\theta(\widetilde{\mathbf Z}_n))
            \right\|_{\mathcal U}^2
        \right]
    \right)^{1/2}.
    \label{eq:transition_noise_robustness}
\end{equation}
Here the expectations are taken jointly over the reference trajectory and the corresponding Gaussian perturbations. The quantity $\eta_n$ directly measures the sensitivity of the latent transition to perturbations of its encoded input, whereas $\rho_n$ measures the sensitivity of the decoder to perturbations around the predicted latent state.

The following result makes this local-stability interpretation explicit by bounding the two sensitivity terms in terms of the noise magnitude and the regularity of the learned maps.

\begin{proposition}[Gaussian sensitivity bounds]
\label{prop:gaussian_sensitivity_bounds}
Let the setting and notation of \S\ref{subsubsec:noise_injection} hold.
Assume that
\[
    \operatorname{Tr}\bigl(K_\phi(\mathbf U)\bigr)
    \leq \varsigma,
    \qquad
    \alpha_\theta(\mathbf Z)\leq \bar\alpha
\]
on the relevant regions, and suppose that $T_\theta:\mathcal Z\to\mathcal Z$
and $D_\psi:\mathcal Z\to\mathcal U$ are Lipschitz with constants $\Lambda_T$ and $\Lambda_D$, respectively. Then
\begin{equation}
    \eta_n
    \leq
    \Lambda_T\sqrt{\varsigma},
    \qquad
    \rho_n
    \leq
    \Lambda_D\bar\alpha
    \sqrt{\operatorname{Tr}(K)}.
    \label{eq:lipschitz_sensitivity_bounds}
\end{equation}
Suppose, in addition, that $T_\theta$ is Fr\'echet differentiable on the relevant encoder-mean neighborhoods and satisfies the uniform second-order remainder bound
\[
    \left\|
        T_\theta(\mathbf z+\mathbf h)
        -
        T_\theta(\mathbf z)
        -
        J_{T_\theta}(\mathbf z)\mathbf h
    \right\|_{\mathcal Z}
    \leq
    \frac{M_T}{2}\|\mathbf h\|_{\mathcal Z}^2.
\]
Then
\begin{equation}
\begin{aligned}
    \eta_n
    \leq\left(\bbE\left[\Vert J_{T_\theta}(\mathbf Z^\star_n)\Vert_\mathrm{op}^2\right]\varsigma\right)^{1/2}+\frac{\sqrt{3}M_T}{2}\varsigma.
    \label{eq:tau_jacobian_bound}
\end{aligned}
\end{equation}
Likewise, suppose that $D_\psi$ is Fr\'echet differentiable on the relevant transition-noise neighborhoods and satisfies
\[
    \left\|
        D_\psi(\mathbf z+\mathbf h)
        -
        D_\psi(\mathbf z)
        -
        J_{D_\psi}(\mathbf z)\mathbf h
    \right\|_{\mathcal U}
    \leq
    \frac{M_D}{2}\|\mathbf h\|_{\mathcal Z}^2.
\]
Then
\begin{equation}
\rho_n\leq\ol\alpha\left(\bbE\left[\Vert J_{D_\psi}(T_\theta(\wt{\mathbf Z}_n))\Vert_\mathrm{op}^2\right]\right)^{1/2}\sqrt{\operatorname{Tr} K}+\frac{\sqrt{3}M_D\ol\alpha^2}{2}\operatorname{Tr} K.
    \label{eq:rho_jacobian_banach}
\end{equation}
\end{proposition}

Proposition~\ref{prop:gaussian_sensitivity_bounds} gives two complementary views of the sensitivity induced by Gaussian perturbations. The Lipschitz bounds in \eqref{eq:lipschitz_sensitivity_bounds} provide global control in terms of the overall noise magnitude, whereas the local estimates \eqref{eq:tau_jacobian_bound}-\eqref{eq:rho_jacobian_banach} show that, for small perturbations, the sensitivities are governed primarily by the local Fr\'echet derivatives of the transition and decoder. In particular, $\eta_n$ measures the response of $T_\theta$ to encoder perturbations, while $\rho_n$ measures the response of $D_\psi$ to perturbations around the predicted latent state. The higher-order terms vanish with the corresponding noise scales.

We can interpret how stochastic training acts on these sensitivity terms from an optimization perspective. The predictive objective evaluates
\[
    D_\psi\left(
        T_\theta(\mathbf Z_n^\star+\Xi_n)
        +
        \alpha_\theta(\widetilde{\mathbf Z}_n)\mathbf G_n
    \right)
\]
over Gaussian neighborhoods of the nominal latent trajectory rather than only at their centers. Large local variations of $T_\theta$ in the encoder-noise directions or of $D_\psi$ in the transition-noise directions can therefore increase the noisy prediction error. Through the reparameterized samples, back-propagation updates the model using errors evaluated throughout these neighborhoods. Thus, although the sensitivity terms are not explicitly penalized, noise-injection training implicitly encourages local stability in the regions and directions explored by the Gaussian perturbations.

\iffalse
\subsubsection{A Perspective from Noise-Injection Training}
We now interpret the functional Gaussian latent Markov model as a noise-injection training scheme. Let $\mathbf U\in\mathcal U$ denote a physical PDE state and let $\mathbf Z\in\mathcal Z$ denote its latent representation, where $\mathcal Z$ is a separable Hilbert space. In the variational model, the latent state is not obtained from a deterministic encoder alone. Instead, we use the Gaussian encoder kernel
\[
    Q_\phi(\d \mathbf Z\mid \mathbf U)
    =
    \mathsf N\big(m_\phi(\mathbf U),K_\phi(\mathbf U)\big),
\]
where $m_\phi:\mathcal U\to\mathcal Z$ is the encoder mean and $K_\phi(\mathbf U)$ is a positive,
self-adjoint, trace-class covariance operator on $\mathcal Z$. Thus an encoded latent sample can be
written as
\[
    \widetilde{\mathbf Z}
    =
    m_\phi(\mathbf U)+\Xi,
    \qquad
    \Xi\sim \mathsf N(0,K_\phi(\mathbf U)).
\]
The learned latent transition is modeled by the Gaussian Markov kernel
\[
    P_\theta(\d \mathbf Z^+\mid \mathbf Z)
    =
    \mathsf N\big(T_\theta(\mathbf Z),\alpha_\theta(\mathbf Z)^2K\big),
\]
where $T_\theta:\mathcal Z\to\mathcal Z$ is the mean transition map,
$\alpha_\theta:\mathcal Z\to[0,\infty)$ is the learned transition noise amplitude, and
$K:\mathcal Z\to\mathcal Z$ is a positive, self-adjoint, trace-class covariance operator. By the
reparameterization trick, the transition sample from the noisy encoded input
$\widetilde{\mathbf Z}$ can be written as
\[
    \widetilde{\mathbf Y}
    =
    T_\theta(\widetilde{\mathbf Z})
    +
    \alpha_\theta(\widetilde{\mathbf Z})\mathbf G,
    \qquad
    \mathbf G\sim \mathsf N(0,K).
\]
The decoder mean $D_\psi:\mathcal Z\to\mathcal U$ maps this perturbed latent prediction back to
the physical solution space. Therefore, during training, the one-step decoded prediction is evaluated
at
\[
    D_\psi(\widetilde{\mathbf Y})
    =
    D_\psi\!\left(
    T_\theta(m_\phi(\mathbf U)+\Xi)
    +
    \alpha_\theta(m_\phi(\mathbf U)+\Xi)\mathbf G
    \right).
\]
For example, with squared reconstruction loss, the stochastic one-step prediction objective contains
\[
    \mathbb E_{\Xi,\mathbf G}
    \left[
    \left\|
    D_\psi\!\left(
    T_\theta(m_\phi(\mathbf U)+\Xi)
    +
    \alpha_\theta(m_\phi(\mathbf U)+\Xi)\mathbf G
    \right)
    -
    \mathbf U^+
    \right\|_{\mathcal U}^2
    \right],
\]
where
\[
    \Xi\sim\mathsf N(0,K_\phi(\mathbf U)),
    \qquad
    \mathbf G\sim\mathsf N(0,K).
\]
Thus the variational latent model injects noise at two levels. The encoder covariance $K_\phi(\mathbf U)$ perturbs the input to the learned transition, while the transition covariance $K$ perturbs the output of the latent mean map before decoding. Consequently, the decoder is not trained only at the nominal mean prediction
\[
    D_\psi\big(T_\theta(m_\phi(\mathbf U))\big),
\]
but over a Gaussian neighborhood induced jointly by the encoder noise and the transition noise.

This viewpoint is important because the encoder mean defines the reference latent trajectory used
in the rollout analysis. Given a reference physical trajectory
$\{\mathbf U_n^\star\}_{n=0}^N$, we set
\[
    \mathbf Z_n^\star := m_\phi(\mathbf U_n^\star),
    \qquad
    \bar{\mathbf U}_n := D_\psi(\mathbf Z_n^\star).
\]
We do not require exact reconstruction, so in general
\[
    \bar{\mathbf U}_n
    =
    D_\psi(m_\phi(\mathbf U_n^\star))
    \neq
    \mathbf U_n^\star .
\]
The reconstruction defect
\[
    \delta_n
    :=
    \|\bar{\mathbf U}_n-\mathbf U_n^\star\|_{\mathcal U}
\]
therefore appears explicitly in the rollout bound.

At inference time, for simplicity and stability, we use the deterministic mean rollout. Given an
initial physical state $\mathbf U_0$, we initialize the latent state by the encoder mean,
\[
    \widehat{\mathbf Z}_0 := m_\phi(\mathbf U_0),
\]
and then propagate by
\[
    \widehat{\mathbf Z}_{n+1}
    =
    T_\theta(\widehat{\mathbf Z}_n),
    \qquad
    \widehat{\mathbf U}_n
    =
    D_\psi(\widehat{\mathbf Z}_n).
\]
Thus neither encoder noise nor transition noise is injected during rollout. The stochasticity is used
as a training-side regularizer, while the deployed trajectory follows the deterministic mean latent
dynamics.

For comparison, a deterministic latent solver uses the same inference-time structure, but with a
deterministic encoder, transition, and decoder:
\[
    \widehat{\mathbf Z}_0
    =
    m_\phi^{\mathrm{det}}(\mathbf U_0),
    \qquad
    \widehat{\mathbf Z}_{n+1}
    =
    T_\theta^{\mathrm{det}}(\widehat{\mathbf Z}_n),
    \qquad
    \widehat{\mathbf U}_n
    =
    D_\psi^{\mathrm{det}}(\widehat{\mathbf Z}_n).
\]
It is trained through a pointwise one-step mean prediction loss of the form
\[
    \left\|
    D_\psi^{\mathrm{det}}
    \big(
    T_\theta^{\mathrm{det}}(m_\phi^{\mathrm{det}}(\mathbf U))
    \big)
    -
    \mathbf U^+
    \right\|_{\mathcal U}^2 .
\]
Consequently, the main distinction is not the inference-time recursion, since both models use deterministic mean rollout. The distinction lies in the training objective: the deterministic baseline fits only the nominal encoded transition, whereas the variational model fits a Gaussian-smoothed one-step prediction map around both the encoded input state and the latent transition output.

\begin{proposition}[Rollout bound with variational training]
\label{prop:encoder-noisy-rollout-bound}
Let
\[
\{\mathbf U^\star_n\}_{n=0}^N\subset\mathcal U
\]
be a reference physical trajectory. Define the reference latent trajectory by the encoder mean
\[
\mathbf Z_n^\star := m_\phi(\mathbf U^\star_n),\quad n=0,1,\cdots,N,
\]
and define the decoded reference trajectory by
\[
\bar{\mathbf U}_n := D_\psi(\mathbf Z_n^\star),\quad n=0,1,\cdots,N.
\]
We allow imperfect reconstruction and set $\delta_n := \|\bar{\mathbf U}_n-\mathbf U_n^\star\|_{\mathcal U}$. For each $n=0,\ldots,N-1$, let
\[
\Xi_n\sim\mathsf N(0,K_\phi(\mathbf U^\star_n)),
\qquad
\mathbf G_n\sim\mathsf N(0,K),
\]
with $\Xi_n$ and $G_n$ independent. Then the inputs of transition kernel $P_\theta$ and the decoder map $D_\psi$ are
\[
\widetilde{\mathbf Z}_n := \mathbf Z_n^\star+\Xi_n,\qquad
\widetilde{\mathbf Y}_n
:=T_\theta(\widetilde{\mathbf Z}_n)+\alpha_\theta(\widetilde{\mathbf Z}_n)\mathbf G_n .
\]
We also make the following assumptions:
\begin{itemize}[itemsep=-2pt, topsep=0pt]
    \item[(i)] The encoder covariance operator $K_\phi$ has uniformly bounded trace, i.e.  there exists some $\varsigma>0$ such that $\mathrm{Tr}(K_\phi(\mathbf{U}))\leq\varsigma$ for all $\mathbf U\in\mathcal{U}$.
    \item[(ii)] The transition noise amplitude $\alpha_\theta$ is  bounded by some $\ol\alpha>0$, i.e. $\alpha_\theta(\mathbf Z)\leq\ol\alpha$ for all $\mathbf Z\in\mathcal Z$.
    \item[(iii)] The transition map $T_\theta$ and the decoder $D_\psi$ are Lipschitz continuous on their domains.
\end{itemize}
Define the noisy prediction error $\epsilon_n$, transition-noise decoder robustness $\rho_n$, and encoder-noise propagation error $\eta_n$ by
\[
\epsilon_n:=\left(\mathbb E_{\Xi_n,\mathbf G_n}
\left[\left\|D_\psi(\widetilde{\mathbf Y}_n)-\mathbf U^\star_{n+1}
\right\|_{\mathcal U}^2\right]\right)^{1/2},
\]
\[
\rho_n:=\left(\mathbb E_{\Xi_n,\mathbf G_n}
\left[\left\|D_\psi(\widetilde{\mathbf Y}_n)-D_\psi\left(T_\theta(\widetilde{\mathbf Z}_n)\right)
\right\|_{\mathcal U}^2\right]\right)^{1/2},
\]
and
\[
\eta_n:=\left(\mathbb E_{\Xi_n}
\left[\left\|D_\psi\left(T_\theta(\widetilde{\mathbf Z}_n)\right)
-D_\psi\left(T_\theta(\mathbf Z_n^\star)\right)
\right\|_{\mathcal U}^2\right]\right)^{1/2}.
\]
Then the following hold.
\begin{itemize}[topsep=0pt, itemsep=0pt]
\item[(a)]\textit{One-step deterministic mean control.}
For every $n=0,\ldots,N-1$,
\[
\left\|
D_\psi\left(T_\theta(\mathbf Z_n^\star)\right)-\mathbf U^\star_{n+1}
\right\|_{\mathcal U}
\le\epsilon_n+\rho_n+\eta_n .
\]

\item[(b)]\textit{Lipschitz robustness estimates.}
Let $\mathcal R_D\subset\mathcal Z$ be a latent region containing the relevant arguments of $D_\psi$, and let $\mathcal R_T\subset\mathcal Z$ be a latent region containing the relevant arguments of $T_\theta$. Define
\[
\Lambda_D := \|D_\psi\|_{\mathrm{Lip}(\mathcal R_D)},
\qquad
\Lambda_T := \|T_\theta\|_{\mathrm{Lip}(\mathcal R_T)}.
\]
Then
\[
\eta_n\le\Lambda_D\Lambda_T\varsigma^{1/2}.
\]
Moreover, if $\alpha_\theta(z)\le \ol\alpha$ on the relevant noisy tube, then
\[
\rho_n\le\Lambda_D\ol\alpha\sqrt{\operatorname{Tr}(K)}.
\]

\item[(c)]\textit{Local Fr\'echet-derivative robustness estimates.} 
Assume that $D_\psi$ is Fr\'echet differentiable on the relevant transition-noise neighborhoods and that, uniformly for relevant $\mathbf w$, with $\mathbf z=T_\theta(\mathbf w)$,
\[
\left\|D_\psi(\mathbf z+\mathbf h)-D_\psi(\mathbf z)-JD_\psi(\mathbf z)\mathbf h\right\|_{\mathcal U}\le\frac{M_D}{2}\|\mathbf h\|_{\mathcal Z}^2.
\]
Then
\[
\rho_n\le\ol\alpha\left(\mathbb  E_{\Xi_n}\left[\bigl\Vert JD_\psi(T_\theta(\widetilde{\mathbf Z}_n))\bigr\Vert_{\mathrm{op}}^2\right]\operatorname{Tr}(K)\right)^{1/2}+
\frac{\sqrt{3}M_D\ol\alpha^2}{2}\mathrm{Tr}(K).
\]
In addition, assume that $T_\theta$ is Fr\'echet differentiable at $\mathbf Z_n^\star$, and the second-order remainder satisfies
\[
\left\|T_{\theta}(\mathbf Z_n^\star+h)-T_{\theta}(\mathbf Z_n^\star)-JT_{\theta}(\mathbf Z_n^\star)h \right\|_{\mathcal U}
\le\frac{M_T}{2}\|h\|_{\mathcal Z}^2.
\]
Then
\begin{equation*}
\begin{aligned}
\eta_n&\le
\Vert JD_\psi(T_\theta(\mathbf Z_n^\star))\Vert_{\mathrm{op}}\Vert JT_\theta(\mathbf Z_n^\star)\Vert_{\mathrm{op}}\varsigma^{1/2}
\\
&\qquad +
\varsigma\left(\frac{\sqrt{3}}{2}M_T\Vert JD_\psi(T_\theta(\mathbf Z_n^\star))\Vert_{\mathrm{op}}+\sqrt{6}M_D\Vert JT_\theta(\mathbf Z_n^\star)\Vert_{\mathrm{op}}^2\right)+\sqrt{\frac{105}{8}}\varsigma^2M_DM_T^2.
\end{aligned}
\end{equation*}

\item[(d)]\textit{Mean-rollout bound.}
Let $\wh{\mathbf Z}_0=\mathbf Z_0^\star=m_\phi(\mathbf{U}_0^\star)$, and let the deterministic inference rollout be
\[
\widehat{\mathbf Z}_{n}=T_\theta(\widehat{\mathbf Z}_{n-1}),\qquad\widehat{\mathbf U}_n=D_\psi(\widehat{\mathbf Z}_n),\quad n=1,\cdots,N.
\]
Assume that the decoder is locally inverse-stable on a latent manifold $\mathcal M$ containing the relevant latent states,
\[
\|\mathbf z-\mathbf z'\|_{\mathcal Z}\le\Lambda_{\mathrm{inv}}\|D_\psi(\mathbf z)-D_\psi(\mathbf z')\|_{\mathcal U},
\]
where $\Lambda_{\mathrm{inv}}:=\|(D_\psi|_{\mathcal M})^{-1}\|_{\mathrm{Lip}(D_\psi(\mathcal M))}$. Then, for every $n=1,\ldots,N$,
\[
\|\widehat{\mathbf U}_n-\mathbf U_n^\star\|_{\mathcal U}
\le\Lambda_D\Lambda_{\mathrm{inv}}
\sum_{j=0}^{n-1}\Lambda_T^{n-1-j}
\left(\epsilon_j+\rho_j+\eta_j+\delta_{j+1}\right)+\delta_n .
\]
\end{itemize}
\end{proposition}

\begin{remark}[How variational training controls the four one-step terms]
\label{rem:how-training-controls-rollout-terms}
Proposition~\ref{prop:encoder-noisy-rollout-bound} shows that the deterministic rollout error is controlled by four one-step quantities: $\epsilon_n,\rho_n,\eta_n,\delta_n.$
These terms enter the bound in different ways and are affected differently by the variational
training objective.

First, the noisy prediction error $\epsilon_n$ is directly controlled by the stochastic one-step
prediction loss. Indeed, for the reference pair
$(\mathbf U_n^\star,\mathbf U_{n+1}^\star)$, the variational training pathway samples
\[
    \widetilde{\mathbf Z}_n
    =
    m_\phi(\mathbf U_n^\star)+\Xi_n,
    \qquad
    \Xi_n\sim\mathsf N(0,K_\phi(\mathbf U_n^\star)),
\]
and then
\[
    \widetilde{\mathbf Y}_n
    =
    T_\theta(\widetilde{\mathbf Z}_n)
    +
    \alpha_\theta(\widetilde{\mathbf Z}_n)\mathbf G_n,
    \qquad
    \mathbf G_n\sim\mathsf N(0,K).
\]
Thus the stochastic physical prediction loss contains
\[
    \mathbb E_{\Xi_n,\mathbf G_n}
    \left[
    \left\|
    D_\psi(\widetilde{\mathbf Y}_n)
    -
    \mathbf U_{n+1}^\star
    \right\|_{\mathcal U}^2
    \right]
    =
    \epsilon_n^2,
\]
and minimizing the noisy one-step prediction loss directly decreases $\epsilon_n$ along the training trajectory. Likewise, the reconstruction defect $\delta_n$ is controlled by the encoder-decoder reconstruction loss. Since the reference latent state is defined by the encoder mean, $\mathbf Z_n^\star=m_\phi(\mathbf U_n^\star)$, the decoded reference state is
\[
    \bar{\mathbf U}_n
    =
    D_\psi(\mathbf Z_n^\star)
    =
    D_\psi(m_\phi(\mathbf U_n^\star)).
\]
Therefore
\[
    \delta_n^2
    =
    \left\|
    D_\psi(m_\phi(\mathbf U_n^\star))
    -
    \mathbf U_n^\star
    \right\|_{\mathcal U}^2
\]
is precisely the deterministic encoder-mean reconstruction error. Thus a static reconstruction term for the autoencoder directly controls the reconstruction defect appearing in the rollout bound. If the reconstruction loss is evaluated using stochastic encoder samples $\mathbf Z_n=m_\phi(\mathbf U_n^\star)+\Xi_n$, then it controls a noisy reconstruction analogue and additionally regularizes the decoder in a neighborhood of the encoder mean.

The remaining two terms, $\rho_n$ and $\eta_n$, are controlled more implicitly. They are not usually separate loss terms. Instead, they quantify the local sensitivity of the decoded prediction map on precisely the regions explored by the injected training noise.

The transition-noise robustness term
\[
    \rho_n
    =
    \left(
    \mathbb E_{\Xi_n,\mathbf G_n}
    \left[
    \left\|
    D_\psi(\widetilde{\mathbf Y}_n)
    -
    D_\psi(T_\theta(\widetilde{\mathbf Z}_n))
    \right\|_{\mathcal U}^2
    \right]
    \right)^{1/2}
\]
measures the sensitivity of the decoder to the transition perturbation
\[
    \alpha_\theta(\widetilde{\mathbf Z}_n)\mathbf G_n
\]
around the random transition center $T_\theta(\widetilde{\mathbf Z}_n)$. Although $\rho_n$ is not penalized explicitly, the noisy prediction loss requires the decoded output to remain accurate not only at the mean transition point $T_\theta(\widetilde{\mathbf Z}_n)$, but throughout the $K$-weighted Gaussian neighborhood sampled during training. If the decoder changes sharply in this neighborhood, then the perturbed predictions
\[
    D_\psi\!\left(
    T_\theta(\widetilde{\mathbf Z}_n)
    +
    \alpha_\theta(\widetilde{\mathbf Z}_n)\mathbf G_n
    \right)
\]
will generally vary significantly, increasing the noisy prediction loss. Thus minimizing $\epsilon_n$ encourages local decoder stability in the covariance directions explored by the transition noise. This interpretation is made explicit by the Fr\'echet-derivative estimate in Proposition~\ref{prop:encoder-noisy-rollout-bound}. Locally,
\[
    \rho_n
    \lesssim
    \left(
    \mathbb E_{\Xi_n}
    \left[
    \alpha_\theta(\widetilde{\mathbf Z}_n)^2
    \operatorname{Tr}
    \left(
    JD_\psi(T_\theta(\widetilde{\mathbf Z}_n))
    K
    JD_\psi(T_\theta(\widetilde{\mathbf Z}_n))^\ast
    \right)
    \right]
    \right)^{1/2}
    +\ 
    \text{higher-order terms}.
\]
Thus $\rho_n$ is small when the decoder derivative is small in the directions emphasized by the transition covariance $K$, or more globally when the decoder has a small effective Lipschitz constant on the noisy transition region.

Similarly, the encoder-noise propagation term
\[
    \eta_n
    =
    \left(
    \mathbb E_{\Xi_n}
    \left[
    \left\|
    D_\psi(T_\theta(\mathbf Z_n^\star+\Xi_n))
    -
    D_\psi(T_\theta(\mathbf Z_n^\star))
    \right\|_{\mathcal U}^2
    \right]
    \right)^{1/2}
\]
measures the sensitivity of the decoded transition map
\[
    H_{\theta,\psi}:=D_\psi\circ T_\theta
\]
to perturbations of the encoded input state. The encoder-noise part of training forces the model to
predict accurately not only from the encoder mean
$m_\phi(\mathbf U_n^\star)$, but also from nearby latent states
$m_\phi(\mathbf U_n^\star)+\Xi_n$. Therefore, the noisy objective smooths the learned decoded
transition map around the encoded data manifold. In the local Fr\'echet-derivative view,
\[
    \eta_n
    \lesssim
    \left[
    \operatorname{Tr}
    \left(
    JH_{\theta,\psi}(\mathbf Z_n^\star)
    K_\phi(\mathbf U_n^\star)
    JH_{\theta,\psi}(\mathbf Z_n^\star)^\ast
    \right)
    \right]^{1/2}
    +
    \text{higher-order terms}.
\]
By the chain rule, $JH_{\theta,\psi}(\mathbf Z_n^\star)= JD_\psi(T_\theta(\mathbf Z_n^\star))JT_\theta(\mathbf Z_n^\star)$.
Thus $\eta_n$ is governed by the covariance-weighted sensitivity of the composed decoded transition map $D_\psi\circ T_\theta$ in the encoder-noise directions. In a more global view, this
corresponds to controlling an effective Lipschitz constant of the decoded transition map on the Gaussian neighborhood induced by $K_\phi(\mathbf U_n^\star)$.

In summary, $\epsilon_n$ and $\delta_n$ are directly targeted by the prediction and reconstruction losses, respectively. The robustness terms $\rho_n$ and $\eta_n$ are controlled indirectly: Gaussian noise injection replaces pointwise fitting by fitting over local neighborhoods in latent space. This acts as a smoothing regularizer, discouraging large decoder sensitivity and large decoded-transition sensitivity near the training trajectory. Consequently, variational training can reduce the local one-step defects and the effective sensitivity factors that are amplified by the rollout recursion.
\end{remark}

\begin{remark}[Deterministic latent rollout bound]
\label{rem:deterministic-rollout-bound}
For comparison, consider a deterministic latent solver with encoder $m_\phi^{\mathrm{det}}:\mathcal U\to\mathcal Z$, transition map $T_\theta^{\mathrm{det}}:\mathcal Z\to\mathcal Z$,
and decoder $D_\psi^{\mathrm{det}}:\mathcal Z\to\mathcal U$. Given the same reference physical trajectory
$\{\mathbf U_n^\star\}_{n=0}^N$, write
\[
    \mathbf Z_n^{\star,\mathrm{det}}:=m_\phi^{\mathrm{det}}(\mathbf U_n^\star),
    \qquad
    \bar{\mathbf U}_n^{\mathrm{det}}
    :=D_\psi^{\mathrm{det}}(\mathbf Z_n^{\star,\mathrm{det}}).
\]
Like before, define the deterministic reconstruction defect by
\[
    \delta_n^{\mathrm{det}}
    :=
    \|\bar{\mathbf U}_n^{\mathrm{det}}-\mathbf U_n^\star\|_{\mathcal U},
\]
and the deterministic one-step decoded transition
error by
\[
    \epsilon_n^{\mathrm{det}}
    :=\left\|D_\psi^{\mathrm{det}}\left(T_\theta^{\mathrm{det}}(\mathbf Z_n^{\star,\mathrm{det}})\right)
    -
    \mathbf U_{n+1}^\star
    \right\|_{\mathcal U}.
\]
At inference time, one initializes $\widehat{\mathbf Z}_0^{\mathrm{det}}
    =
    \mathbf Z_0^{\star,\mathrm{det}}$,
and rolls out deterministically by
\[
    \widehat{\mathbf Z}_{n}^{\mathrm{det}}
    =
    T_\theta^{\mathrm{det}}(\widehat{\mathbf Z}_{n-1}^{\mathrm{det}}),
    \qquad
    \widehat{\mathbf U}_n^{\mathrm{det}}
    =
    D_\psi^{\mathrm{det}}(\widehat{\mathbf Z}_n^{\mathrm{det}}),\qquad n=1,\cdots,N.
\]
Assume that $T_\theta^{\mathrm{det}}$ and $D_\psi^{\mathrm{det}}$ are Lipschitz on the relevant regions, with constants
\[
    \Lambda_T^{\mathrm{det}}
    :=
    \|T_\theta^{\mathrm{det}}\|_{\mathrm{Lip}},
    \qquad
    \Lambda_D^{\mathrm{det}}
    :=
    \|D_\psi^{\mathrm{det}}\|_{\mathrm{Lip}},
\]
and assume that $D_\psi^{\mathrm{det}}$ is locally inverse-stable on the relevant latent manifold with constant $\Lambda_{\mathrm{inv}}^{\mathrm{det}}$. Then, for every $n=1,\ldots,N$,
\begin{equation}
    \left\|
    \widehat{\mathbf U}_n^{\mathrm{det}}
    -\mathbf U_n^\star\right\|_{\mathcal U}
    \le\Lambda_D^{\mathrm{det}}
    \Lambda_{\mathrm{inv}}^{\mathrm{det}}
    \sum_{j=0}^{n-1}
    \left(\Lambda_T^{\mathrm{det}}\right)^{n-1-j}\left(\epsilon_j^{\mathrm{det}}+\delta_{j+1}^{\mathrm{det}}\right)+\delta_n^{\mathrm{det}}.\label{eq:deterministic_rollout_bound}
\end{equation}
The derivation is similar to Proposition~\ref{prop:encoder-noisy-rollout-bound} and is deferred to the appendix.

Therefore, the deterministic and variational rollout error bounds have the same amplification mechanism: one-step errors are accumulated with powers of the Lipschitz factor of the learned transition map. The difference is in the one-step quantities. The deterministic solver controls only the pointwise transition error $\epsilon_j^{\mathrm{det}}$ and reconstruction defect $\delta_j^{\mathrm{det}}$, whereas the variational solver additionally involves the robustness terms $\rho_j$ and $\eta_j$, which reflect stability under transition noise and encoder noise.
\end{remark}

\paragraph{Mechanism: reducing the base of rollout amplification.}
The bounds above show that long-horizon error is governed not only by the size of the one-step defects, but also by the Lipschitz factor of the transition map. In the variational model, the rollout error contains the geometric weights
\[
    \Lambda_T^{\,n-1-j},
    \qquad
    \Lambda_T=\|T_\theta\|_{\mathrm{Lip}},
\]
whereas in the deterministic baseline the corresponding weights are
\[
    \left(\Lambda_T^{\mathrm{det}}\right)^{n-1-j},
    \qquad
    \Lambda_T^{\mathrm{det}}
    =
    \|T_\theta^{\mathrm{det}}\|_{\mathrm{Lip}} .
\]
Thus $\Lambda_T$ and $\Lambda_T^{\mathrm{det}}$ are the bases of the exponentially growing rollout factors. Even a modest reduction in the effective transition Lipschitz factor can therefore produce a large improvement over long horizons.

In variational training, stochasticity is used to smooth the learned prediction map around the data
manifold. For a training pair $(\mathbf U,\mathbf U^+)$, the variational decoded prediction is
evaluated at
\[
    D_\psi\!\left(
    T_\theta(m_\phi(\mathbf U)+\Xi)
    +
    \alpha_\theta(m_\phi(\mathbf U)+\Xi)\mathbf G
    \right),
    \qquad
    \Xi\sim\mathsf N(0,K_\phi(\mathbf U)),
    \quad
    \mathbf G\sim\mathsf N(0,K).
\]
The encoder noise $\Xi$ perturbs the input of the transition map, while the transition noise $\mathbf G$ perturbs the latent output before decoding. Hence the variational objective does not fit only a pointwise map, but rather a Gaussian-smoothed neighborhood of the composed decoded transition $H_{\theta,\psi}:=D_\psi\circ T_\theta$.

This smoothing directly controls the noisy one-step prediction error $\epsilon_n$ and, through the reconstruction objective, the reconstruction defect $\delta_n$. It also implicitly controls the robustness terms $\rho_n$ and $\eta_n$. Locally, the Fr\'echet-derivative estimates show that
\[
    \rho_n^2
    \approx
    \alpha_\theta(\widetilde{\mathbf Z}_n)^2
    \operatorname{Tr}
    \left(
    JD_\psi(T_\theta(\widetilde{\mathbf Z}_n))
    K
    JD_\psi(T_\theta(\widetilde{\mathbf Z}_n))^\ast
    \right),
\]
while
\[
    \eta_n^2
    \approx
    \operatorname{Tr}
    \left(
    JH_{\theta,\psi}(\mathbf Z_n^\star)
    K_\phi(\mathbf U_n^\star)
    JH_{\theta,\psi}(\mathbf Z_n^\star)^\ast
    \right).
\]
By the chain rule,
\[
    JH_{\theta,\psi}(\mathbf Z_n^\star)
    =
    JD_\psi(T_\theta(\mathbf Z_n^\star))
    JT_\theta(\mathbf Z_n^\star).
\]
Therefore, transition noise penalizes sharp decoder variations near predicted latent outputs, while encoder noise penalizes sharp variations of the composed decoded transition map with respect to the transition input.

Consequently, variational training can mitigate rollout amplification in two complementary ways. First, it reduces the one-step quantities $\epsilon_n, \rho_n, \eta_n, \delta_n$ that are accumulated in the rollout bound. Second, because encoder noise probes the input sensitivity of $T_\theta$ through $D_\psi\circ T_\theta$, the smoothed objective discourages transition maps with large local expansion near the encoded trajectory. This can reduce the effective value of $\Lambda_T$ relative to the deterministic baseline $\Lambda_T^{\mathrm{det}}$, thereby reducing the base of the exponential amplification factor.
\fi

\subsubsection{Variational Alignment and Rollout Control}
The previous subsection studies how stochastic sampling regularizes the local sensitivity of the learned transition and decoder. We now turn to the variational KL term. At the distributional level, Proposition~\ref{prop:population_elbo_transition_consistency} shows that this term aligns the learned latent transition with the encoder-induced dynamics of true PDE trajectories. Under the Gaussian parameterization, this probabilistic alignment further yields quantitative control of the corresponding latent means. The following lemma makes this connection precise by bounding the discrepancy between the transition mean and the mean of a target latent distribution in terms of their KL divergence.

\begin{lemma}[KL control of latent means]
\label{lem:kl_latent_mean_control}
Let
\[
P_\theta(\cdot\mid\mathbf{Z})=\mathsf N\left(T_\theta(\mathbf{Z}),\alpha_\theta(\mathbf{Z})^2K\right),
\]
where $K$ is positive, self-adjoint, and trace-class on $\mathcal Z$. Let $Q$ be any probability measure on $\mathcal Z$ with finite second moment and mean
\[
    m_Q:=\int_{\mathcal Z} y\,\d Q(y).
\]
Assume that
\[
D_{\mathrm{KL}}\left(Q\,\Vert\,P_\theta(\cdot\,|\,\mathbf{Z})\right)<\infty .
\]
Then
\[
\left\|m_Q-T_\theta(\mathbf{Z})\right\|_{\mathcal Z}
\le\alpha_\theta(\mathbf{Z})\sqrt{
    2\|K\|_{\mathrm{op}}
    D_{\mathrm{KL}}
    \left(Q\,\Vert\,P_\theta(\cdot\,|\,\mathbf{Z})\right)
}.
\]
\end{lemma}
Lemma~\ref{lem:kl_latent_mean_control} converts the dynamic KL divergence into quantitative control of the discrepancy between latent means. In particular, when the transition distribution is Gaussian, a small KL divergence between the encoded next-state distribution and the predicted transition distribution implies that their mean states must also be close. This is the key step that connects the probabilistic alignment enforced by the variational objective to the deterministic mean dynamics used at inference. The proof relies on a Gaussian transportation inequality and is deferred to Appendix~\S\ref{app:kl-aware-rollout}.

We now introduce the remaining population quantities needed for the rollout analysis. Motivated by Lemma~\ref{lem:kl_latent_mean_control}, for the random reference trajectory, define the dynamic KL error
\[
    \kappa_n
    :=
    \left(
        \mathbb E
        \left[
            D_{\mathrm{KL}}
            \left(
                Q_\phi(\cdot\mid\mathbf U_{n+1}^\star)
                \,\Vert\,
                P_\theta(\cdot\mid\widetilde{\mathbf Z}_n)
            \right)
        \right]
    \right)^{1/2}.
\]
Thus, $\kappa_n$ measures, at the population level, how well the learned transition distribution from the perturbed current latent state matches the encoder distribution of the true next physical state. It is the direct population counterpart of the dynamic KL term appearing in the training objective.

In addition, we do not assume exact reconstruction of the physical state from its encoder mean, and therefore define
\[
    \delta_n
    :=
    \left(
        \mathbb E
        \left[
            \|D_\psi(\mathbf Z_n^\star)-\mathbf U_n^\star\|_{\mathcal U}^2
        \right]
    \right)^{1/2}.
\]
The quantity $\delta_n$ measures the discrepancy introduced when the encoder-mean latent state is decoded back to the physical space. Here the expectations are taken over the reference trajectory and, where applicable, the encoder perturbation.

Together with the sensitivity quantities studied in \S\ref{subsubsec:noise_injection}, we have the necessary ingredients for analyzing deterministic autoregressive rollout. The following proposition characterizes how the corresponding one-step discrepancies propagate and accumulate over the rollout horizon.

\begin{proposition}[Population KL-aware rollout control with predictive refinement]
\label{prop:population_kl_rollout}
Let the setting and notation of \S\ref{subsubsec:noise_injection} hold. Assume that
\[
    0<\alpha_\theta(\mathbf Z)\leq\bar\alpha
\]
on the relevant latent region, and that $T_\theta$ and $D_\psi$ are Lipschitz with constants $\Lambda_T$ and $\Lambda_D$, respectively. Then the following hold.

\begin{enumerate}[label=(\alph*),topsep=0pt,itemsep=3pt]

\item
The population latent one-step defect satisfies
\begin{equation}
    \left(
        \mathbb E
        \left[
            \left\|
                T_\theta(\mathbf Z_n^\star)
                -
                \mathbf Z_{n+1}^\star
            \right\|_{\mathcal Z}^2
        \right]
    \right)^{1/2}
    \leq
    \eta_n+\bar\alpha\sqrt{2\|K\|_{\mathrm{op}}}\,\kappa_n.
    \label{eq:population_latent_one_step}
\end{equation}

\item
For the deterministic mean rollout
\eqref{eq:deterministic_mean_rollout},
\begin{equation}
    \left(
        \mathbb E
        \left[
            \|\widehat{\mathbf U}_n-\mathbf U_n^\star\|_{\mathcal U}^2
        \right]
    \right)^{1/2}
    \leq\Lambda_D
    \sum_{j=0}^{n-1}\Lambda_T^{\,n-1-j}\left(\eta_j+\bar\alpha\sqrt{2\|K\|_{\mathrm{op}}}\,\kappa_j
    \right)
    +
    \delta_n.
    \label{eq:population_kl_physical_rollout}
\end{equation}

\item
Incorporating the noisy predictive error gives the refined bound
\begin{equation}
\begin{aligned}
\left(\mathbb E\left[\|\widehat{\mathbf U}_n-\mathbf U_n^\star\|_{\mathcal U}^2\right]\right)^{1/2}
&\leq\Lambda_D
    \sum_{j=0}^{n-2}
    \Lambda_T^{\,n-1-j}\left(\eta_j+\bar\alpha\sqrt{2\|K\|_{\mathrm{op}}}\,\kappa_j\right)+\Lambda_D\eta_{n-1}
    \\
    &\qquad
    +\min\left\{ \Lambda_D\bar\alpha\sqrt{2\|K\|_{\mathrm{op}}}\,\kappa_{n-1}+\delta_n\epsilon_{n-1}+\rho_{n-1}\right\},
\end{aligned}
\label{eq:population_refined_rollout}
\end{equation}
for every $n\geq 1$, where the sum is understood as zero when $n=1$.
\end{enumerate}
\end{proposition}

\begin{remark}[Interpretation of the rollout bound]
\label{rem:population_rollout_interpretation}
The terms appearing in Proposition~\ref{prop:population_kl_rollout} are controlled by different components of the training objective. The dynamic KL error $\kappa_n$ is the population counterpart of the latent consistency term $\mathcal L_{\mathrm{KL}}$, while the reconstruction error $\delta_n$ is associated with the encoder-decoder reconstruction objective. The noisy prediction error $\epsilon_n$ is controlled through $\mathcal L_{\mathrm{pred}}$, and the sensitivity quantities $\eta_n$ and $\rho_n$ are studied in Proposition~\ref{prop:gaussian_sensitivity_bounds}. As discussed in \S\ref{subsubsec:noise_injection}, these sensitivity terms are not explicitly optimized, but stochastic training evaluates the model over Gaussian neighborhoods and thereby implicitly encourages local stability in the corresponding perturbation directions.

Part~(b) makes explicit how the two terms in the conditional ELBO \eqref{eq:condnegelbo} enter deterministic rollout. The static reconstruction term is reflected by the reconstruction error $\delta_n$, which measures the discrepancy between the true physical state and the decoder applied to its encoder-mean representation. The dynamic KL term is reflected by $\kappa_n$, which controls the mismatch between the predicted transition mean and the encoder mean of the true next state. The additional sensitivity term $\eta_n$ accounts for the fact that the KL objective is evaluated at a perturbed encoder input, whereas deterministic rollout propagates the encoder mean. These one-step latent discrepancies are accumulated through the mean transition with geometric factors $\Lambda_T^{n-1-j}$, while $\delta_n$ enters when the propagated latent state is decoded back to the physical space. Thus, the two principal terms of the conditional ELBO reappear naturally in the rollout error bound, linking variational training to deterministic long-horizon accuracy.

Part~(c) additionally incorporates the predictive objective into the same deterministic rollout analysis. At the final prediction step, the physical error admits two complementary upper bounds: one obtained from the KL-controlled latent mean discrepancy together with reconstruction,
\[
    \Lambda_D\bar\alpha\sqrt{2\|K\|_{\mathrm{op}}}\,
    \kappa_{n-1}+\delta_n,
\]
and the other from the noisy predictive error together with transition-noise sensitivity,
\[
    \epsilon_{n-1}+\rho_{n-1}.
\]
Taking the minimum simply selects the tighter analytical estimate for the same deterministic prediction and therefore sharpens the physical-space rollout bound. In particular, this refinement reflects the complementary roles of the variational and predictive components of the training objective.

Finally we discuss the error amplification in rollouts. The geometric factors $\Lambda_T^{n-1-j}$ characterize the growth rate of the rollout error bound. If the one-step error and sensitivity terms remain uniformly bounded, then $\Lambda_T<1$ yields a uniformly bounded geometric accumulation, $\Lambda_T=1$ leads to at most linear growth with the rollout length, and $\Lambda_T>1$ allows geometric growth in the worst case. Thus, the bound highlights two ingredients for long-horizon accuracy: controlling the local one-step discrepancies and limiting their amplification through the learned latent transition.


Finally, we discuss the error amplification in rollouts and characterize the growth rate of the rollout error bound more explicitly. Suppose that the one-step error and sensitivity terms remain uniformly bounded over the rollout horizon:
\[
\tau_j \leq \bar{\tau}, \qquad
\kappa_j \leq \bar{\kappa}, \qquad
\delta_j \leq \bar{\delta}.
\]
Then Proposition~\ref{prop:population_kl_rollout} (b) implies
\[
\left(
\mathbb{E}
\|\widehat{\mathbf{U}}_n - \mathbf{U}_n^\star\|_{\mathcal{U}}^2
\right)^{1/2}
\leq
\begin{cases}
\displaystyle
\Lambda_D
\left(
\bar{\tau}
+
\bar{\alpha}\sqrt{2\|K\|_{\mathrm{op}}}\,\bar{\kappa}
\right)
\frac{1-\Lambda_T^n}{1-\Lambda_T}
+\bar{\delta},
& \Lambda_T < 1,\\[1.2em]
\displaystyle
n\,\Lambda_D
\left(
\bar{\tau}
+
\bar{\alpha}\sqrt{2\|K\|_{\mathrm{op}}}\,\bar{\kappa}
\right)
+\bar{\delta},
& \Lambda_T = 1,\\[1.2em]
\displaystyle
\Lambda_D
\left(
\bar{\tau}
+
\bar{\alpha}\sqrt{2\|K\|_{\mathrm{op}}}\,\bar{\kappa}
\right)
\frac{\Lambda_T^n - 1}{\Lambda_T - 1}
+\bar{\delta},
& \Lambda_T > 1.
\end{cases}
\]
Therefore, under uniformly controlled error and sensitivity terms, the bound remains uniformly bounded for $\Lambda_T < 1$, grows at most linearly for $\Lambda_T = 1$, and permits geometric growth for $\Lambda_T > 1$. The result thus separates two ingredients of long-horizon accuracy: the magnitude of the local one-step discrepancies and their amplification through the learned latent transition.
\end{remark}

\paragraph{Generic autoregressive amplification.} To distinguish the generic effect of autoregressive error accumulation from the mechanisms specific to variational latent dynamics, we consider a generic deterministic predictor $F_\vartheta:\mathcal U\to\mathcal U$, which is deployed recursively as
\[
\widehat{\mathbf U}_0=\mathbf U_0^\star,
\qquad
\widehat{\mathbf U}_{n+1}
=
F_\vartheta(\widehat{\mathbf U}_n).
\]
Define its population one-step prediction error by
\[
\epsilon_n^{\mathrm{dir}}
:=
\left(\mathbb E
\left[\|F_\vartheta(\mathbf U_n^\star)-\mathbf U_{n+1}^\star\|_{\mathcal U}^2\right]\right)^{1/2}.
\]
The following result isolates how these one-step errors accumulate solely through repeated application of the learned predictor.

\begin{proposition}[Generic autoregressive error amplification]
\label{prop:generic_rollout}
Suppose that $F_\vartheta$ is Lipschitz on the relevant region with constant
$\Lambda_F$. Then
\begin{equation}
\left(
\mathbb E
\left[\|\widehat{\mathbf U}_n-\mathbf U_n^\star\|_{\mathcal U}^2
\right]\right)^{1/2}
\leq\sum_{j=0}^{n-1}\Lambda_F^{\,n-1-j}\epsilon_j^{\mathrm{dir}}.
\label{eq:generic_rollout}
\end{equation}
\end{proposition}

Proposition~\ref{prop:generic_rollout} applies to deterministic autoregressive predictors such as the direct neural-operator baseline considered in our subsequent experiments. In particular, \eqref{eq:generic_rollout} shows that geometric error amplification is a generic feature of autoregressive prediction rather than a consequence of the variational formulation. Although the variational latent model evolves autoregressively in latent rather than physical space, Proposition~\ref{prop:population_kl_rollout} exhibits the same amplification mechanism through powers of $\Lambda_T$. Thus, the preceding analysis characterizes how the variational training  controls the one-step quantities that are subsequently amplified during latent rollout.

The sensitivity analysis in \S\ref{subsubsec:noise_injection} suggests a complementary role for stochastic training in the autoregressive amplification mechanism. Both the generic deterministic bound and the variational rollout bound contain geometric amplification factors like $\Lambda_F$ and $\Lambda_T$, governed respectively by the stability of the learned one-step maps. In the variational latent formulation, encoder perturbations expose the transition map to neighborhoods of the encoded trajectory, while transition perturbations probe the decoder around predicted latent states. Proposition~\ref{prop:gaussian_sensitivity_bounds} shows that these sensitivities are locally governed by the Fréchet derivatives of the corresponding maps. Sharp local fluctuations in rollout-relevant neighborhoods can therefore increase the noisy predictive loss and are implicitly discouraged through reparameterized back-propagation. In this way, variational training can regularize the effective local expansion encountered along the rollout trajectory and thereby mitigate geometric error amplification over long horizons.


\section{Variational Autoencoding Markov Operator}
\label{sec:vamo}

The variational latent dynamics developed in
\S\ref{sec:var_latent_dyn} are formulated directly at the function-space level and are therefore independent of a particular spatial discretization or neural architecture. In this section, we introduce the \emph{Variational Autoencoding Markov Operator} (VAMO), a discretized neural-operator realization of this framework for time-dependent PDEs. VAMO realizes the encoder-transition-decoder structure of \S\ref{sec:abstract_vld} using spatially resolved physical and latent fields. We first construct its deterministic backbone from a resolution-preserving encoder and decoder together with a residual neural-operator transition in \S\ref{subsec:vamo_architecture}. We then augment this backbone with structured Gaussian latent perturbations in \S\ref{subsec:structured_latent_perturbations} and formulate the resulting training objective in \S\ref{subsec:vamo_training}.

\subsection{Spatially Resolved Architecture}
\label{subsec:vamo_architecture}
To obtain a computational realization of the function-space framework, we first discretize the physical and latent fields on a finite spatial grid and specify the deterministic mean maps underlying VAMO. This subsection focuses on the encoder mean, latent mean transition, and decoder mean that form the deterministic backbone of the model; the Gaussian covariance structure and stochastic sampling used for variational training are introduced separately in \S\ref{subsec:structured_latent_perturbations}.

Let $\{x_1,x_2,\cdots,x_{N_\mathrm{res}}\}\subset\Omega$ denote a spatial grid on the domain $\Omega$ with $N_{\mathrm{res}}$ discretization points. Evaluating a physical field with $C_u$ channels and its latent counterpart with $C_z$ channels on this grid gives
\[ 
\mathbf u_n \in \mathbb R^{C_u\times N_{\mathrm{res}}}, \quad
\mathbf z_n \in \mathbb R^{C_z\times N_{\mathrm{res}}}. 
\] 
Thus, spatial discretization converts the function-space variables $\mathbf U_n\in\mathcal U$ and $\mathbf Z_n\in\mathcal Z$ into finite-dimensional arrays while preserving their spatial organization. In particular, the latent representation remains a field over the computational grid rather than being collapsed into a single latent vector. Under this discretization, the encoder, transition, and decoder mean maps from Section~\ref{sec:abstract_vld} are represented by 
\[ 
m_\phi: \mathbb R^{C_u\times N_{\mathrm{res}}} \rightarrow \mathbb R^{C_z\times N_{\mathrm{res}}}, \quad 
T_\theta: \mathbb R^{C_z\times N_{\mathrm{res}}} \rightarrow \mathbb R^{C_z\times N_{\mathrm{res}}},\quad\text{and}\quad
D_\psi: \mathbb R^{C_z\times N_{\mathrm{res}}} \rightarrow \mathbb R^{C_u\times N_{\mathrm{res}}}. 
\] 
The encoder $m_\phi$ first lifts the physical state to a spatially resolved latent representation, which is subsequently evolved by $T_\theta$ and mapped back to the physical variables through $D_\psi$. We implement the encoder and decoder using resolution-preserving convolutional residual networks, while the latent dynamics are modeled by a Fourier neural operator. Specifically, the mean transition takes the residual form 
\begin{equation} 
T_\theta(\mathbf z) = \mathbf z + \mathcal G_\theta^{\mathrm{FNO}}(\mathbf z), \label{eq:vamo_residual_transition} 
\end{equation} 
where $\mathcal G_\theta^{\mathrm{FNO}}$ is an FNO acting on the latent feature field. The residual parameterization represents one-step evolution through an increment of the current latent state, consistent with the time-discretized PDE flow map, while the Fourier layers provide nonlocal interactions across the spatial domain. The three mean maps therefore define the deterministic backbone 
\[ 
\mathbf u_n \xrightarrow{\,m_\phi\,} \mathbf z_n \xrightarrow{\,T_\theta\,} \mathbf z_{n+1} \xrightarrow{\,D_\psi\,} \mathbf u_{n+1}. 
\] 
VAMO augments this backbone with state-dependent Gaussian perturbations of the encoded and propagated latent fields. We next describe how the corresponding covariance structures are constructed on the discretized grid.

\subsection{Structured Gaussian Latent Perturbations}
\label{subsec:structured_latent_perturbations}

The deterministic backbone in \S\ref{subsec:vamo_architecture} specifies the mean evolution of the latent state. To realize the variational formulation, VAMO augments this evolution with Gaussian perturbations whose covariance preserves the spatial structure of the latent field. Rather than injecting independent noise at individual grid points, we generate correlated perturbations through fixed spectral covariance operators and modulate their amplitudes according to the current state.

\paragraph{Spectral covariance structure.}
We first describe the covariance model shared by the encoder and transition perturbations. On a periodic spatial grid, let $K_{\ell,s}$ denote the discrete analogue of the Laplacian-resolvent covariance
\[
    K_{\ell,s}
    =
    (I-\ell^2\Delta)^{-s},
\]
which is diagonal in the Fourier basis. Its square root therefore acts as
\begin{equation}
    \widehat{K_{\ell,s}^{1/2}\boldsymbol\xi}(k)
    =
    \left(1+\ell^2|k|^2\right)^{-s/2}
    \widehat{\boldsymbol\xi}(k),
    \label{eq:discrete_spectral_filter}
\end{equation}
where $\boldsymbol\xi$ is a standard Gaussian field and $k$ denotes the discrete spatial frequency. Thus, $\ell>0$ controls the correlation length of the perturbation, while $s>0$ determines the spectral decay at high frequencies. This construction is the finite-dimensional counterpart of the Laplacian-resolvent Gaussian covariance and Sobolev geometry discussed in \S\ref{sec:functional_gaussian_latent_markov}.

\paragraph{Complementary encoder and transition perturbations.}
For the Gaussian kernels, we use separate covariance operators
\[
    K_{\mathrm{enc}}=K_{\ell_{\mathrm{enc}},s},
    \quad
    K_{\mathrm{tr}}=K_{\ell_{\mathrm{tr}},s}
\]
for the encoder and latent transition, with correlation lengths chosen independently. In practice, we use $\ell_{\mathrm{enc}}<\ell_{\mathrm{tr}}$ and combine these covariance structures with different state-dependent modulations.

\begin{enumerate}[topsep=2pt]
\item\textit{State-dependent encoder distribution.} In addition to the encoder mean $m_\phi(\mathbf u)$, the encoder network produces a pointwise log-variance modulation field
\[
    \ell_\phi(\mathbf u)
    \in
    \mathbb R^{C_z\times N_{\mathrm{res}}}.
\]
We define the corresponding standard-deviation multiplier
\[
    \sigma_\phi(\mathbf u)
    :=
    \exp\!\left(\frac12\ell_\phi(\mathbf u)\right),
\]
and let $S_\phi(\mathbf u)$ denote pointwise multiplication by
$\sigma_\phi(\mathbf u)$. A latent state is generated via the reparameterization
\begin{equation}
    \mathbf z
    =
    m_\phi(\mathbf u)
    +
    S_\phi(\mathbf u)
    K_{\mathrm{enc}}^{1/2}\boldsymbol\xi,\qquad
    \boldsymbol\xi\sim\mathsf N(0,I),
    \label{eq:discrete_encoder_reparam}
\end{equation}
or equivalently,
\[
    \mathbf z
    \sim
    \mathsf N\bigl(
        m_\phi(\mathbf u),
        K_\phi(\mathbf u)
    \bigr),\quad K_\phi(\mathbf u)=S_\phi(\mathbf u)K_{\mathrm{enc}}S_\phi^*(\mathbf u).
\]
Hence, the fixed covariance operator $K_{\mathrm{enc}}$ specifies the underlying spatial correlation structure, while the learned multiplier $S_\phi(\mathbf u)$ adapts the perturbation magnitude across latent channels, spatial locations, and input states.

\item\textit{Stochastic latent transition.} The transition distribution uses the same structured-noise principle, but modulates the fixed covariance $K_{\mathrm{tr}}$ through a scalar state-dependent amplitude. Given a latent state $\mathbf z$, we sample
\begin{equation}
    \mathbf z^+
    =
    T_\theta(\mathbf z)
    +
    \alpha_\theta(\mathbf z)
    K_{\mathrm{tr}}^{1/2}\boldsymbol\xi^{\mathrm{tr}},
    \qquad
    \boldsymbol\xi^{\mathrm{tr}}
    \sim
    \mathsf N(0,I),
    \label{eq:discrete_transition_reparam}
\end{equation}
which defines
\begin{equation}
    P_\theta(\cdot\mid\mathbf z)
    =
    \mathsf N\!\left(
        T_\theta(\mathbf z),
        \alpha_\theta(\mathbf z)^2K_{\mathrm{tr}}
    \right).
    \label{eq:discrete_transition_kernel}
\end{equation}
The amplitude $\alpha_\theta(\mathbf z)>0$ is predicted by a lightweight network head and adapts the overall perturbation magnitude to the current latent state. The precise numerical parameterization of $\alpha_\theta$ and $\ell_\phi$ is deferred to Appendix~\S\ref{app:implementation_details}.
\end{enumerate}
This construction yields an intentional local-global asymmetry between the two stochastic components. The encoder combines the relatively small correlation length $\ell_{\mathrm{enc}}$ with a pointwise, spatially heterogeneous multiplier, thereby probing robustness to localized and fine-scale variations in the encoded physical state. In contrast, the transition instead combines the larger correlation length $\ell_{\mathrm{tr}}$ with a single state-dependent amplitude applied globally across the latent field, producing smoother and more spatially coherent perturbations around the evolved latent state. Thus, the encoder probes spatially heterogeneous local uncertainty, whereas the transition probes coherent variations of the predicted latent dynamics. Together, these complementary perturbation mechanisms realize the stochastic training effects analyzed in \S\ref{subsubsec:noise_injection}.

\paragraph{Decoder likelihood.}
Finally, we equip the discretized decoder with an isotropic Gaussian likelihood,
\begin{equation}
    P_\psi(\cdot\mid\mathbf z)
    =
    \mathsf N\!\left(
        D_\psi(\mathbf z),
        \sigma_u^2 I
    \right),
    \label{eq:discrete_decoder_likelihood}
\end{equation}
where $\sigma_u>0$ is fixed. This is the finite-dimensional specialization of the Gaussian decoder considered in Proposition~\ref{prop:cm_decoder_likelihood}. In the present setting, the Cameron-Martin space is the entire discretized physical space and its norm is a scaled Euclidean norm:
\[
    \|\mathbf v\|_{\mathcal H_{\sigma_u^2 I}}^2
    =
    \frac{1}{\sigma_u^2}\|\mathbf v\|_2^2.
\]
Consequently, up to an additive constant, the decoder negative
log-likelihood is
\[
    -\log P_\psi(\mathbf u\mid\mathbf z)
    =
    \frac{1}{2\sigma_u^2}
    \|\mathbf u-D_\psi(\mathbf z)\|_2^2
    +\mathrm{const}.
\]
Therefore, the Cameron-Martin reconstruction geometry of the function-space formulation reduces exactly, up to a fixed scaling, to the squared Euclidean norm loss used in the discretized model.

Unlike the encoder and transition distributions, we do not sample from the decoder likelihood. Given a latent state $\mathbf z$, prediction uses $D_\psi(\mathbf z)$ directly, which is both the mean and the \textit{maximum-likelihood} output of the isotropic Gaussian decoder. Hence the stochasticity used for variational training is confined to the latent representation and transition rather than being propagated into the physical prediction.

The fixed likelihood scale $\sigma_u>0$ determines the relative weighting between reconstruction error and latent KL regularization in the ELBO. Indeed, multiplying the negative ELBO by $2\sigma_u^2$ converts the likelihood term to an unweighted squared Euclidean loss while rescaling the KL term. Accordingly, rather than treating $\sigma_u$ as an additional tunable uncertainty parameter, we fix it and absorb this relative scaling into the KL coefficient $\beta$ in the training objective introduced below.

\subsection{Training Objective}
\label{subsec:vamo_training}
Having specified the discretized encoder, transition, and decoder
distributions, we now formulate the training objective on observed PDE trajectories as the finite-dimensional counterpart of the variational objective studied in \S\ref{sec:err_var_dyn}. Let $\{\mathbf u_n\}_{n=0}^{N}$ denote a trajectory from the training set. For each adjacent pair $(\mathbf u_n,\mathbf u_{n+1})$, we draw
\[
    \mathbf z_n\sim Q_\phi(\cdot\mid\mathbf u_n),
    \qquad
    \mathbf z_n^+\sim P_\theta(\cdot\mid\mathbf z_n),
\]
using the reparameterizations in \eqref{eq:discrete_encoder_reparam} and \eqref{eq:discrete_transition_reparam}. The resulting one-step objective contains three components:
\begin{align}
    \mathcal L_{\mathrm{pred}}^{(n)}
    &:=
    \mathbb E_{\mathbf z_n\sim Q_\phi(\cdot\mid\mathbf u_n)}
    \mathbb E_{\mathbf z_n^+\sim P_\theta(\cdot\mid\mathbf z_n)}
    \left[
        \left\|
            \mathbf u_{n+1}
            -
            D_\psi(\mathbf z_n^+)
        \right\|_2^2
    \right],
    \label{eq:vamo_pred_loss}
    \\
    \mathcal L_{\mathrm{KL}}^{(n)}
    &:=
    \mathbb E_{\mathbf z_n\sim Q_\phi(\cdot\mid\mathbf u_n)}
    \left[
        D_{\mathrm{KL}}
        \left(
            Q_\phi(\cdot\mid\mathbf u_{n+1})
            \,\Vert\,
            P_\theta(\cdot\mid\mathbf z_n)
        \right)
    \right],
    \label{eq:vamo_kl_loss}
    \\
    \mathcal L_{\mathrm{rec}}^{(n)}
    &:=
    \mathbb E_{\mathbf z_n\sim Q_\phi(\cdot\mid\mathbf u_n)}
    \left[
        \left\|
            \mathbf u_n-D_\psi(\mathbf z_n)
        \right\|_2^2
    \right].
    \label{eq:vamo_rec_loss}
\end{align}
We optimize the trajectory-averaged objective
\begin{equation}
    \mathcal L(\theta,\phi,\psi)
    =
    \frac{1}{N}
    \sum_{n=0}^{N-1}
    \left(
        \mathcal L_{\mathrm{pred}}^{(n)}
        +
        \beta\,\mathcal L_{\mathrm{KL}}^{(n)}
        +
        \lambda_{\mathrm{rec}}\,
        \mathcal L_{\mathrm{rec}}^{(n)}
    \right),
    \label{eq:implemented_vlm_loss}
\end{equation}
where the expectations are estimated by Monte Carlo samples through the reparameterized encoder and transition distributions.

The three terms correspond directly to the mechanisms analyzed in \S\ref{sec:err_var_dyn}. The predictive term trains the full stochastic pathway $\mathbf u_n\to\mathbf z_n\to\mathbf z_n^+\to\mathbf u_{n+1}$, the KL term aligns the predicted latent transition with the distribution obtained by encoding the true next state, and the reconstruction term maintains an informative latent representation of the physical field.

\paragraph{Relation to the function-space objective.} 
Equation~\eqref{eq:implemented_vlm_loss} is a finite-dimensional realization of a weighted combination of the predictive negative log-likelihood and the negative conditional ELBO. To make this relation explicit, consider
\[
    \lambda_{\mathrm{pred}}
    \mathcal L_{\mathrm{pred}}^{\mathrm{NLL}}
    +
    \lambda_{\mathrm{ELBO}}
    \left(
        \mathcal L_{\mathrm{KL}}+\mathcal L_{\mathrm{rec}}^{\mathrm{NLL}}
    \right).
\]
Under the isotropic decoder likelihood \eqref{eq:discrete_decoder_likelihood}, both likelihood terms reduce to squared Euclidean losses with the factor $(2\sigma_u^2)^{-1}$. Normalizing the coefficient of the predictive Euclidean loss to one therefore gives, up to additive constants,
\[
    \mathcal L_{\mathrm{pred}}
    +\underbrace{
    2\sigma_u^2\frac{\lambda_{\mathrm{ELBO}}}{\lambda_{\mathrm{pred}}}}_{\beta}\mathcal L_{\mathrm{KL}}
    +\underbrace{\frac{\lambda_{\mathrm{ELBO}}}{\lambda_{\mathrm{pred}}}}_{\lambda_{\mathrm{rec}}}\mathcal L_{\mathrm{rec}}.
\]
Hence, up to an overall rescaling, the relative weighting of the predictive bound and conditional ELBO in Proposition~\ref{prop:conditional_functional_elbo}, together with the decoder likelihood scale, is represented equivalently by the two coefficients $\beta$ and $\lambda_{\mathrm{rec}}$ in \eqref{eq:implemented_vlm_loss}.

Notably, the conditional ELBO in
Proposition~\ref{prop:conditional_functional_elbo} naturally contains next-state reconstruction, whereas \eqref{eq:vamo_rec_loss} uses the current state. Define the step-$n$ reconstruction term
\[
    \mathcal R_n
    :=
    \mathbb E_{\mathbf z_n\sim Q_\phi(\cdot\mid\mathbf u_n)}
    \left[
        -\log P_\psi(\mathbf u_n\mid\mathbf z_n)
    \right].
\]
Over a complete trajectory, the ELBO reconstruction terms sum to
\[
    \sum_{n=0}^{N-1}\mathcal R_{n+1}
    =
    \sum_{n=1}^{N}\mathcal R_n,
\]
whereas the implemented current-state reconstruction sums to $\sum_{n=0}^{N-1}\mathcal R_n$. The two therefore differ only through the boundary contributions $\mathcal R_0$ and $\mathcal R_N$. We use the current-state form because it reuses the latent sample $\mathbf z_n$ already drawn for the transition step, avoiding an additional reparameterized sample from $Q_\phi(\cdot\mid\mathbf u_{n+1})$ for reconstruction. This provides a modest computational efficiency gain during training while leaving all interior reconstruction contributions unchanged.

\section{Numerical Experiments}
\label{sec:experiments}

In this section, we evaluate VAMO on three fluid-dynamics benchmarks with distinct dynamical characteristics: the one-dimensional compressible Euler equations, two-dimensional compressible fluid dynamics, and forced two-dimensional incompressible Navier--Stokes equations. Together, these problems cover inviscid wave propagation, coupled compressible flow, and long-time vortex dynamics under different levels of physical dissipation. Our experiments are designed to answer three main questions:
\begin{itemize}[itemsep=0pt, topsep=2pt]
\item Can variational latent dynamics improve the accuracy and stability of long-horizon autoregressive PDE prediction?
\item Are the improvements attributable to the structured variational formulation, rather than latent compression or generic noise injection alone?
\item Does VAMO better preserve physically relevant spatial and spectral structures throughout the rollout?
\end{itemize}

For all benchmarks, the model is trained using trajectories restricted to a finite training horizon and evaluated by autoregressive rollout over a \textit{substantially longer test horizon}. During inference, the model is initialized from the exact initial condition and receives no additional ground-truth states. We report errors accumulated over the entire test rollout. This evaluation therefore measures both prediction accuracy within the temporal regime represented during training and stability beyond the training horizon.

\subsection{Benchmarks and Baselines}
\label{sec:benchmarks_baselines}

\subsubsection{Benchmarks Problems}
\label{sec:benchmarks}

We consider three representative fluid-dynamics systems, summarized in Table~\ref{tab:benchmark_summary}. We provide a brief description here and defer the complete governing equations, data-generation procedures, discretization schemes, and dataset configurations to Appendix~\S\ref{app:benchmark_details}.

\paragraph{1D compressible Euler equations.}
The one-dimensional Euler equations describe inviscid compressible flow through conservation of mass, momentum, and total energy. We represent the physical state using the primitive variables
\[
    \mathbf{u}(t,x)
    =
    \bigl(\rho(t,x),\, v(t,x),\, p(t,x)\bigr),\quad t\in[0,T],\ x\in\mathbb{T}_L.
\]
where $\rho$, $v$, and $p$ denote density, velocity, and pressure, respectively. Since the system contains no explicit viscous dissipation, prediction errors in wave speed and phase can accumulate over time, while nonlinear evolution may generate sharp gradients and discontinuous structures. This benchmark therefore tests long-horizon propagation of compressible waves in a controlled one-dimensional setting.
\paragraph{2D compressible fluid dynamics.}
We consider a two-dimensional viscous compressible-flow benchmark with a physical configuration similar to the CFD setting studied in PDEBench \citep{takamoto2022pdebench}. The state is represented by
\[
    \mathbf{u}(t,\mathbf{x})
    =
    \bigl(
        \rho(t,\mathbf{x}),
        v_x(t,\mathbf{x}),
        v_y(t,\mathbf{x}),
        p(t,\mathbf{x})
    \bigr),\quad t\in[0,T],\ \mathbf{x}\in\mathbb{T}^2.
\]
This system couples density, pressure, and the two velocity components through nonlinear conservation laws and viscous transport. It provides a challenging multi-field benchmark in which local prediction errors can propagate across physical channels and spatial locations.

\paragraph{2D incompressible Navier-Stokes equations.}
We study the forced two-dimensional incompressible Navier-Stokes equations in vorticity form on the unit torus \citep{li2022transformer}. The learned state is the scalar vorticity field
\[
    \omega(t,\mathbf{x}),\quad t\in[0,T],\ \mathbf{x}\in\mathbb{T}^2.
\]
while the corresponding velocity field is recovered nonlocally through the Biot-Savart operator. We consider several viscosity levels and two families of external forcing. Under weak viscosity, small errors in the predicted vorticity can induce nonlocal velocity errors that feed back into future advection. This benchmark is therefore particularly suited to evaluating long-time rollout stability and preservation of spatial and spectral flow structure.

\begin{table}[t]
\small
    \centering
    \caption{
    Summary of the fluid-dynamics benchmarks.
    The three systems span one-dimensional inviscid compressible waves,
    two-dimensional viscous compressible flow, and forced incompressible
    vorticity dynamics.
    }
    \label{tab:benchmark_summary}
    \begin{tabular}{llll}
        \toprule
        Benchmark
        & PDE type
        & Fields
        & Main challenge \\
        \midrule
        1D Euler
        & Compressible, inviscid
        & $(\rho,v,p)$
        & Wave propagation, phase error, shocks \\
        2D CFD
        & Compressible, viscous
        & $(\rho,\mathbf{v},p)$
        & Coupled multi-field dynamics, conservation \\
        2D Navier-Stokes
        & Incompressible, viscous
        & $\omega$
        & Nonlocal transport and long-time stability \\
        \bottomrule
    \end{tabular}
\end{table}

\subsubsection{Compared Methods}
\label{sec:compared_methods}
We compare VAMO with three baselines that isolate the effects of direct physical-space evolution, generic noise injection, and deterministic latent dynamics. All models are trained as one-step predictors and evaluated through autoregressive rollout. \begin{description}[topsep=2pt,itemsep=0pt]
\item\textbf{FNO.} We use the Fourier neural operator \citep{li2020fourier} as a direct physical-space baseline. Unlike the trajectory-to-trajectory setting considered in the original FNO experiments, our model maps the current physical snapshot to the next snapshot and is applied recursively over the test horizon. 
\item\textbf{FNO+Noise.} This baseline uses the same autoregressive FNO architecture, but injects Gaussian perturbations into the lifted feature field during training. It tests whether generic feature-space noise regularization alone can improve long-horizon stability. 
\item\textbf{FNO-AE.} This model combines a deterministic autoencoder with a residual FNO latent transition. The initial physical state is encoded once, the latent field is evolved recursively, and the resulting latent states are decoded into physical predictions. It therefore isolates deterministic latent evolution from the variational training used by VAMO. 
\item\textbf{FNO-VAMO.} This denotes the complete proposed model described in \S\ref{sec:vamo}. FNO-AE and FNO-VAMO use closely matched latent architectures, allowing their comparison to isolate the effect of the variational formulation. 
\end{description}
The roles of the four comparisons are summarized in Table~\ref{tab:comparison_logic}. 
\begin{table}[t] 
\centering 
\caption{ Experimental role of each model comparison. } \label{tab:comparison_logic} 
\begin{tabular}{ll} 
\toprule Comparison & Question addressed \\ 
\midrule 
FNO-VAMO vs.\ FNO & Does the complete framework improve long-horizon prediction? \\ 
FNO-VAMO vs.\ FNO-AE & Is variational training necessary beyond latent modeling? \\ 
FNO-VAMO vs.\ FNO+Noise & Is structured latent stochasticity better than generic noise? \\ 
FNO-AE vs.\ FNO & Does the deterministic latent architecture alone improve stability? \\ 
\bottomrule 
\end{tabular} 
\end{table} 
For a fair comparison, the models use matched FNO backbones whenever applicable, including the number of Fourier layers, retained modes, and hidden width. All methods are trained using the same trajectories, data splits, optimization budget, and model-selection protocol. Complete architectural and optimization details are provided in Appendix~\S\ref{app:implementation_details}.


\paragraph{Evaluation protocol.}
All models are trained as one-step predictors and evaluated by autoregressive rollout from the exact initial condition, without teacher forcing or ground-truth reinitialization. We report errors over the full test horizon. Let $\mathbf{u}_i^n$ and $\widehat{\mathbf{u}}_i^n$ denote the reference and predicted states for test trajectory $i$ at rollout step $n$. For a spatial norm $\|\cdot\|_\mathcal{X}$, the accumulated relative rollout error is
\[
\mathrm{Relative}\ \mathcal{X}\text{-}\mathrm{error}
=
\frac{1}{N_{\mathrm{test}}}
\sum_{i=1}^{N_{\mathrm{test}}}
\frac{
\left(
\sum_{n=1}^{T_{\mathrm{test}}}
\|\widehat{\mathbf{u}}^i_n-\mathbf{u}^i_n\|_{\mathcal X}^2
\right)^{1/2}
}{
\left(
\sum_{n=1}^{T_{\mathrm{test}}}
\|\mathbf{u}_i^n\|_{\mathcal X}^2
\right)^{1/2}
}.
\]
Relative $L^2$ error is the primary metric across all benchmarks. We additionally report relative $L^1$ error for the 1D Euler equations, relative $H^1$ error for the 2D benchmarks, and errors in enstrophy, palinstrophy, and the isotropic energy spectrum for incompressible Navier-Stokes. Complete metric definitions and numerical implementation details are provided in Appendix~\S\ref{app:benchmark_details}.


\subsection{Main Results}\label{sec:main_results}
We first compare the long-horizon rollout performance of VAMO and the three baselines across the fluid-dynamics benchmarks. For each problem, all models are trained and evaluated using the same trajectories and autoregressive protocol described in Section~\ref{sec:benchmarks_baselines}. The tables report errors accumulated over the full test horizon. We focus here on the overall quantitative comparison and defer the temporal evolution of rollout errors, qualitative predictions, and additional physical diagnostics to the subsequent sections.

\subsubsection{1D Compressible Euler Equations}
\label{sec:euler_results}

We evaluate the models on periodic 1D Euler trajectories with domain lengths $L\in\{5,10,15,20\}$ and corresponding spatial resolutions $1024$, $2048$, $3072$, and $4096$, keeping the grid spacing fixed across settings. Initial conditions are generated from normalized periodic Gaussian random fields with physical correlation length $\ell=1$. Models are trained on trajectories over $[0,1.5]$ and rolled out autoregressively to $T_{\mathrm{test}}=10$ with snapshot interval $\Delta t=0.1$. Complete data-generation and solver details are provided in Appendix~\S\ref{app:euler_data_generation}. We report full-rollout relative $L^1$ and $L^2$ errors for density, velocity, and pressure in Table~\ref{tab:euler1d_results}.

\begin{table}[t]
\centering
\caption{Results on the 1D Euler equations with different domain lengths $L$. We report rollout-aggregated relative $L^1$ and $L^2$ errors for density $\rho$, pressure $p$, velocity $u$, and their channel average. Lower is better.}
\label{tab:euler1d_results}
\resizebox{\textwidth}{!}{
\begin{tabular}{llcccccccc}
\toprule
\multirow{2}{*}{$L$}
& \multirow{2}{*}{Method}
& \multicolumn{2}{c}{Density $\rho$}
& \multicolumn{2}{c}{Pressure $p$}
& \multicolumn{2}{c}{Velocity $u$}
& \multicolumn{2}{c}{Channel Mean} \\
\cmidrule(lr){3-4}
\cmidrule(lr){5-6}
\cmidrule(lr){7-8}
\cmidrule(lr){9-10}
&
& Rel.\,$L^2$ & Rel.\,$L^1$
& Rel.\,$L^2$ & Rel.\,$L^1$
& Rel.\,$L^2$ & Rel.\,$L^1$
& Rel.\,$L^2$ & Rel.\,$L^1$ \\
\midrule
\multirow{4}{*}{$5$}
& FNO       & 0.0194 & 0.00883 & 0.0249 & 0.00951 & 0.1409 & 0.0596 & 0.0617 & 0.0260 \\
& FNO+Noise & 0.0200 & 0.00962 & 0.0259 & 0.0109 & 0.1488 & 0.0727 & 0.0649 & 0.0311 \\
& FNO-AE    & 0.0182 & $\mathbf{0.00812}$ & 0.0243 & 0.01020 & 0.1395 & 0.0682 & 0.0607 & 0.0289 \\
& FNO-VAMO  & $\mathbf{0.0171}$ & 0.0100 & $\mathbf{0.0182}$ & $\mathbf{0.00778}$ & $\mathbf{0.0997}$ & $\mathbf{0.0537}$ & $\mathbf{0.0450}$ & $\mathbf{0.0238}$ \\
\midrule
\multirow{4}{*}{$10$}
& FNO       &0.0206 & 0.00801 & 0.0277 & 0.00942 & 0.1325 & 0.0512 & 0.0602 & 0.0229 \\
& FNO+Noise & 0.0212 & 0.00835 & 0.0286 & 0.00986 & 0.1370 & 0.0540 & 0.0623 & 0.0241 \\
& FNO-AE    & 0.0180 & 0.00753 & 0.0238 & 0.00897 & 0.1237 & 0.0566 & 0.0551 & 0.0244 \\
& FNO-VAMO  & $\mathbf{0.0148}$ & $\mathbf{0.00658}$ & $\mathbf{0.0198}$ & $\mathbf{0.00772}$ & $\mathbf{0.0943}$ & $\mathbf{0.0449}$ & $\mathbf{0.0430}$ & $\mathbf{0.0197}$ \\
\midrule
\multirow{4}{*}{$15$}
& FNO       & 0.0266 & 0.00919 & 0.0369 & 0.0113 & 0.1656 & 0.0646 & 0.0763 & 0.0284 \\
& FNO+Noise  & 0.0263 & 0.00918 & 0.0366 & 0.0113 & 0.1642 & 0.0649 & 0.0757 & 0.0285 \\
& FNO-AE    & 0.0294 & 0.0119 & 0.0400 & 0.0143 & 0.1789 & 0.0696 & 0.0828 & 0.0319 \\
& FNO-VAMO  & $\mathbf{0.0205}$ & $\mathbf{0.00899}$ & $\mathbf{0.0251}$ & $\mathbf{0.00859}$ & $\mathbf{0.1144}$ & $\mathbf{0.0465}$ & $\mathbf{0.0534}$ & $\mathbf{0.0214}$ \\
\midrule
\multirow{4}{*}{$20$}
& FNO       & 0.0380 & 0.0127 & 0.0516 & 0.0159 & 0.2227 & 0.0855 & 0.1099 & 0.0380 \\
& FNO+Noise & 0.0378 & 0.0126 & 0.0511 & 0.0158 & 0.2208 & 0.0851 & 0.1032 & 0.0378 \\
& FNO-AE    & 0.5394 & 0.2259 & 0.5434 & 0.2261 & 2.1236 & 1.0253 & 1.0688 & 0.4924 \\
& FNO-VAMO  & $\mathbf{0.0194}$ & $\mathbf{0.00770}$ & $\mathbf{0.0267}$ & $\mathbf{0.00939}$ & $\mathbf{0.1172}$ & $\mathbf{0.0527}$ & $\mathbf{0.0544}$ & $\mathbf{0.0232}$ \\
\bottomrule
\end{tabular}
}
\end{table}

As shown in Table~\ref{tab:euler1d_results}, FNO-VAMO achieves the lowest channel-averaged relative $L^2$ error for all four domain lengths and the lowest channel-averaged relative $L^1$ error in every setting. Its advantage generally increases with $L$: relative to the strongest baseline, FNO-VAMO reduces the channel-mean $L^2$ error by approximately $26\%$, $22\%$, $29\%$, and $47\%$ for $L=5,10,15$, and $20$, respectively. The largest gains occur in the velocity field, which is consistently the most difficult physical channel. The comparison also distinguishes variational latent training from simpler alternatives. Adding generic noise to FNO provides no consistent benefit, while the deterministic latent model is competitive only for the shorter domains and becomes unstable at $L=20$. FNO-VAMO instead maintains nearly the same error level from $L=10$ to $L=20$, despite larger physical domain and longer-range wave interactions. This suggests that the variational formulation improves the robustness of the recursively evolved latents, rather than merely providing a better one-step fit.


\subsubsection{2D Compressible Fluid Dynamics}
\label{sec:cfd2d_results}

We evaluate the models on 2D compressible Navier-Stokes trajectories on the periodic unit torus. We consider a low-viscosity, low-Mach-number regime with $\mu=\zeta=10^{-8}$ and $M=0.1$, and sample the initial density, pressure, and velocity fields from periodic Gaussian random fields with correlation lengths in $[0.05,0.15]$. The trajectories are solved on a $256\times256$ grid and downsampled to $64\times64$ for learning. Models are trained on $t\in[0,1]$ and evaluated by autoregressive rollout to $T_{\mathrm{test}}=10$ with snapshot interval $\Delta t=0.1$. We choose the low-viscosity regime because larger viscosities rapidly damp the initial perturbations and produce nearly stationary fields, making long-horizon prediction less informative. Complete data-generation details are provided in Appendix~\S\ref{app:cfd2d_data_generation}.

\begin{table}[t]
\centering
\caption{Results on 2D compressible fluid dynamics. We report rollout-aggregated relative $L^2$ and $H^1$ errors for density $\rho$, pressure $p$, velocity $\mathbf v=(v_x,v_y)$, and their channel average. Lower is better.}
\label{tab:cfd2d_results}
\resizebox{\textwidth}{!}{
\begin{tabular}{lcccccccc}
\toprule
\multirow{2}{*}{Method}
& \multicolumn{2}{c}{Density $\rho$}
& \multicolumn{2}{c}{Pressure $p$}
& \multicolumn{2}{c}{Velocity $\mathbf v$}
& \multicolumn{2}{c}{Channel Mean} \\
\cmidrule(lr){2-3}
\cmidrule(lr){4-5}
\cmidrule(lr){6-7}
\cmidrule(lr){8-9}
& Rel.\,$L^2$ & Rel.\,$H^1$
& Rel.\,$L^2$ & Rel.\,$H^1$
& Rel.\,$L^2$ & Rel.\,$H^1$
& Rel.\,$L^2$ & Rel.\,$H^1$ \\
\midrule
FNO        & 0.1248 & 2.2440 & 0.1674 & 3.2094 & 0.8345 & 2.7868 & 0.3756 & 2.7467 \\
FNO+Noise  & 0.1422 & 2.3479 & 0.2073 & 3.8650 & 0.9055 & 3.0393 & 0.4183 & 3.0841 \\
FNO-AE     & 0.0972 & 3.3152 & 0.0331 & 1.7665 & 0.6276 & 1.9364 & 0.2526 & 2.3394 \\
FNO-VAMO   & $\mathbf{0.0290}$ & $\mathbf{0.6375}$ & $\mathbf{0.0090}$ & $\mathbf{0.2612}$ & $\mathbf{0.1316}$ & $\mathbf{0.3145}$ & $\mathbf{0.0565}$ & $\mathbf{0.4044}$ \\
\bottomrule
\end{tabular}
}
\end{table}
As shown in Table~\ref{tab:cfd2d_results}, FNO-VAMO substantially outperforms all baselines across every physical variable and both error metrics. Relative to the strongest baseline, it reduces the channel-averaged relative $L^2$ error from $0.2526$ to $0.0565$, a reduction of approximately $78\%$, and the relative $H^1$ error from $2.3394$ to $0.4044$, a reduction of approximately $83\%$. The improvement is consistent across density, pressure, and velocity, with particularly large gains for the velocity field, which is the most difficult component for all baseline models. The baseline comparisons further show that neither generic noise injection nor deterministic latent evolution is sufficient to obtain stable long-horizon predictions. FNO+Noise performs slightly worse than the direct FNO, indicating that unstructured feature perturbations alone do not improve this benchmark. FNO-AE achieves lower $L^2$ errors than direct FNO, especially for pressure, but its substantially larger density $H^1$ error indicates that improved average field accuracy does not necessarily imply preservation of spatial derivatives. In contrast, FNO-VAMO simultaneously reduces both $L^2$ and $H^1$ errors, suggesting that the variational formulation improves not only coarse field prediction but also the spatial structure of the coupled compressible dynamics.

\subsubsection{2D Incompressible Navier-Stokes Equations}
\label{sec:ns2d_results}
We consider the forced two-dimensional incompressible Navier-Stokes equations in vorticity form on the periodic unit torus. Following the standard FNO benchmark \citep{li2020fourier}, we consider viscosities $\nu\in\{10^{-3},10^{-4},10^{-5}\}$. We extend the setting by independently sampling a time-independent forcing field for each trajectory from either a Gaussian random field (\textsc{GRF}) or a sparse mixture of sinusoidal Fourier modes (\textsc{Wave}). For each viscosity, a single model is trained jointly on 1,600 trajectories, with 800 from each forcing family, and evaluated separately on the two forcing types. Trajectories are generated on a $256\times256$ grid and spectrally truncated to $64\times64$, with snapshots recorded at unit time intervals. The training horizons are $[0,10]$ for $\nu=10^{-3}$, $[0,12]$ for $\nu=10^{-4}$, and $[0,8]$ for $\nu=10^{-5}$, respectively, and the corresponding test horizons are $T=50$, $30$, and $20$. Complete data-generation details are provided in Appendix~\S\ref{app:ns2d_data_generation}.


\begin{table}[htbp]
\centering
\caption{Results on the 2D incompressible Navier-Stokes equations. We evaluate six test settings with viscosity $\nu\in\{10^{-3},10^{-4},10^{-5}\}$ and forcing type in \{GRF, \textsc{Wave}\}. The training horizons are $[0,10]$, $[0,12]$, and $[0,8]$ for $\nu=10^{-3}$, $10^{-4}$, and $10^{-5}$, respectively, while the corresponding test horizons are $[0,50]$, $[0,30]$, and $[0,20]$. We report rollout-aggregated errors. Lower is better.}
\label{tab:ns2d_main_results}
\resizebox{\textwidth}{!}{
\begin{tabular}{lllccccc}
\toprule
Data Setting
& Forcing 
& Method 
& Rel.\,$L^2$ 
& Rel.\,$H^1$ 
& $\mathcal{E}$-error 
& $\mathcal{P}$-error 
& Spec. error \\
\midrule
\multirow{8}{*}{$\begin{aligned}
&\nu=10^{-3}\\ &T_\text{train}=10 \\ &T_\text{test}=50 
\end{aligned}$}
& \multirow{4}{*}{GRF}
& Vanilla FNO         & 0.1453 & 0.8015 & 0.1151 & 3.784 & 0.0239 \\
& 
& FNO+Noise   & 0.0375 & 0.0985 & 0.0196 & 0.1773 & 0.0206 \\
&
& FNO-AE      & 22.31 & 292.0 & 2447 & $4.358\!\times\! 10^5$ & 452.0 \\
&
& FNO-VAMO    & $\mathbf{0.0203}$ & $\mathbf{0.0264}$ & $\mathbf{0.0119}$ & $\mathbf{0.0123}$ & $\mathbf{0.0138}$ \\
\cmidrule(lr){2-8}
& \multirow{4}{*}{\textsc{Wave}}
& Vanilla FNO         & 0.0698 & 0.1769 & 0.0430 & 0.4039 & 0.0116 \\
&
& FNO$+$Noise   & 0.0686 & 0.1148 & 0.0406 & 0.1778 & 0.0241 \\
&
& FNO-AE      & 5.623 & 55.84 & 565.0 & $5.753\!\times\! 10^4$ & 124.3 \\
&
& FNO-VAMO    & $\mathbf{0.0556}$ & $\mathbf{0.0687}$ & $\mathbf{0.00808}$ & $\mathbf{0.00940}$ & $\mathbf{0.0106}$ \\
\midrule
\multirow{8}{*}{$\begin{aligned}
&\nu=10^{-4}\\ &T_\text{train}=12 \\ &T_\text{test}=30 
\end{aligned}$}
& \multirow{4}{*}{GRF}
& Vanilla FNO         & 0.4068 & 3.109 & 0.7338 & 45.84 & 0.0456 \\
&
& FNO+Noise   & 0.2987 & 2.003 & 0.5866 & 35.15 & 0.0569 \\
&
& FNO-AE      & 2.691 & 31.93 & 23.19 & 3608 & 0.8164 \\
&
& FNO-VAMO    & $\mathbf{0.0403}$ & $\mathbf{0.1139}$ & $\mathbf{0.0139}$ & $\mathbf{0.0512}$ & $\mathbf{0.0146}$ \\
\cmidrule(lr){2-8}
& \multirow{4}{*}{\textsc{Wave}}
& Vanilla FNO         & 0.3570 & 1.613 & 0.6580 & 14.51 & 0.0497 \\
&
& FNO+Noise   & 0.3407 & 1.449 & 0.0666 & 15.73 & 0.0499 \\
&
& FNO-AE      & 1.188 & 10.19 & 7.213 & 580.3 & 0.4303 \\
&
& FNO-VAMO    & $\mathbf{0.1198}$ & $\mathbf{0.2212}$ & $\mathbf{0.0386}$ & $\mathbf{0.1008}$ & $\mathbf{0.0390}$ \\
\midrule
\multirow{8}{*}{$\begin{aligned}
&\nu=10^{-5}\\ &T_\text{train}=8 \\ &T_\text{test}=20 
\end{aligned}$}
& \multirow{4}{*}{GRF}
& Vanilla FNO         & 0.2987 & 2.262 & 0.4862 & 32.90 & 0.0382 \\
&
& FNO+Noise   & 0.4089 & 3.061 & 1.179 & 75.58 & 0.0909 \\
&
& FNO-AE      & 0.5991 & 1.335 & 0.2788 & 1.354 & 0.1551 \\
&
& FNO-VAMO    & $\mathbf{0.0841}$ & $\mathbf{0.3858}$ & $\mathbf{0.0211}$ & $\mathbf{0.2415}$ & $\mathbf{0.00929}$ \\
\cmidrule(lr){2-8}
& \multirow{4}{*}{\textsc{Wave}}
& Vanilla FNO & 0.2993 & 1.969 & 0.3519 & 19.14 & 0.0314 \\
&
& FNO+Noise   & 0.3450 & 2.254 & 0.5831 & 31.22 & 0.0617 \\
&
& FNO-AE      & 0.6719 & 1.158 & 0.2098 & 0.5878 & 0.2518 \\
&
& FNO-VAMO    & $\mathbf{0.1469}$ & $\mathbf{0.4509}$ & $\mathbf{0.0460}$ & $\mathbf{0.2734}$ & $\mathbf{0.0247}$ \\
\bottomrule
\end{tabular}
}
\end{table}

As shown in Table~\ref{tab:ns2d_main_results}, FNO-VAMO achieves the lowest error in all six viscosity-forcing settings and across all five evaluation metrics. The improvement is especially pronounced for $\nu=10^{-4}$ and $\nu=10^{-5}$, where weaker dissipation makes the long-horizon rollout more sensitive to accumulated errors. Relative to the strongest baseline in each setting, FNO-VAMO reduces the relative $L^2$ error by approximately $65$--$87\%$ for $\nu=10^{-4}$ and $51$--$72\%$ for $\nu=10^{-5}$. The corresponding reductions in relative $H^1$ error are approximately $85$--$94\%$ and $61$--$71\%$, respectively.

The gains also extend beyond field-level accuracy. FNO-VAMO consistently yields substantially smaller enstrophy and palinstrophy errors, indicating improved preservation of both the overall vorticity magnitude and its small-scale spatial variation. The reduction in palinstrophy error is particularly large in the $\nu=10^{-3}$ and $\nu=10^{-4}$ settings, where the strongest baseline remains one to three orders of magnitude less accurate. FNO-VAMO also achieves the lowest spectral error in every setting, showing that its advantage is not limited to pointwise prediction but extends to the distribution of energy across Fourier modes.

The baseline comparisons reveal two additional trends. First, generic noise injection is not consistently beneficial: FNO+Noise improves over vanilla FNO in some settings, particularly at $\nu=10^{-3}$ and $10^{-4}$, but degrades performance at $\nu=10^{-5}$. Second, the deterministic latent model is highly unstable for $\nu=10^{-3}$ and $10^{-4}$, with severe growth in both field and physical-statistic errors. Although FNO-AE becomes more competitive in some derivative-based metrics at $\nu=10^{-5}$, it remains substantially worse than FNO-VAMO overall. These results indicate that neither latent evolution nor unstructured noise alone is sufficient for reliable long-horizon prediction, whereas the variational latent formulation remains stable across forcing families and viscosity regimes.



\begin{table}[t]
\centering
\resizebox{\textwidth}{!}{%
\begin{tabular}{lllcccc}
\toprule
\multicolumn{3}{c}{Evaluation setting} &
\multicolumn{2}{c}{Training horizon} &
\multicolumn{2}{c}{Full horizon} \\
\cmidrule(lr){1-3}
\cmidrule(lr){4-5}
\cmidrule(lr){6-7}
Viscosity
& Forcing
& Method
& Rel.\,$L^2$
& Rel.\,$H^1$
& Rel.\,$L^2$
& Rel.\,$H^1$ \\
\midrule

\multirow{6}{*}{$\nu=10^{-4}$}
& \multirow{3}{*}{GRF}
& FNO
& $0.0123\pm0.0341$ & $0.1013\pm0.3245$ & $0.4068\pm0.5094$ & $3.1094\pm4.0204$ \\
&
& FNO-Curriculum
& $0.0115\pm0.0064$ & $0.0614\pm0.0378$ & $0.1678\pm0.2085$ & $0.9154\pm1.2697$ \\
&
& VAMO
& $0.0057\pm0.0017$ & $0.0218\pm0.0070$ & $0.0403\pm0.0562$ & $0.1139\pm0.1547$ \\
\cmidrule(lr){2-7}

& \multirow{3}{*}{\textsc{Wave}}
& FNO
& $0.0118\pm0.0363$ & $0.0828\pm0.2990$ & $0.3570\pm0.6102$ & $1.6131\pm2.7778$ \\
&
& FNO-Curriculum
& $0.0120\pm0.0102$ & $0.0508\pm0.0445$ & $0.2072\pm0.2905$ & $0.5334\pm0.7455$ \\
&
& VAMO
& $0.0058\pm0.0031$ & $0.0182\pm0.0110$ & $0.1198\pm0.1785$ & $0.2212\pm0.2680$ \\
\midrule

\multirow{6}{*}{$\nu=10^{-5}$}
& \multirow{3}{*}{GRF}
& FNO
& $0.0132\pm0.0050$ & $0.1186\pm0.0502$ & $0.2987\pm0.4330$ & $2.2621\pm3.6585$ \\
&
& FNO-Curriculum
& $0.0223\pm0.0060$ & $0.1830\pm0.0519$ & $0.1907\pm0.2064$ & $0.9638\pm1.1036$ \\
&
& VAMO
& $0.0139\pm0.0018$ & $0.1238\pm0.0230$ & $0.0841\pm0.0560$ & $0.3858\pm0.1561$ \\
\cmidrule(lr){2-7}

& \multirow{3}{*}{\textsc{Wave}}
& FNO
& $0.0122\pm0.0059$ & $0.0998\pm0.0472$ & $0.2993\pm0.4083$ & $1.9685\pm3.1219$ \\
&
& FNO-Curriculum
& $0.0208\pm0.0077$ & $0.1592\pm0.0632$ & $0.2096\pm0.1818$ & $0.7338\pm0.6364$ \\
&
& VAMO
& $0.0126\pm0.0024$ & $0.1041\pm0.0294$ & $0.1469\pm0.1127$ & $0.4509\pm0.1981$ \\
\bottomrule
\end{tabular}}
\caption{
Comparison with autoregressive curriculum training on the
two more challenging Navier--Stokes regimes.
Errors are evaluated over the training horizon
$[0,T_{\mathrm{train}}]$ and the full rollout horizon
$[0,T_{\mathrm{test}}]$.
All entries report mean $\pm$ standard deviation over 200 test trajectories.
}
\label{tab:curriculum_training}
\end{table}

\subsection{Long-Horizon Error Growth}
\label{sec:error_growth}

The trajectory-aggregated results in Section~\ref{sec:main_results} do not show how prediction errors evolve during autoregressive rollout. We therefore report per-snapshot errors for representative 1D Euler settings with $L=5$ and $L=20$ (Figures \ref{fig:eulerl5}--\ref{fig:eulerl20}), the 2D compressible-flow benchmark (Figure \ref{fig:cfd2derror}), and all six 2D incompressible Navier-Stokes settings (Figures \ref{fig:nu3grf}--\ref{fig:nu5wave}). In each figure, the shaded region denotes the temporal horizon represented during training.

Across the three PDE benchmarks, the main advantage of VAMO appears in the later stages of rollout. Its error can be comparable to, or occasionally higher than, that of the direct baselines at early in-distribution snapshots, particularly for easier settings, e.g. 1D Euler equation with $L=5$. However, the baseline errors generally grow more rapidly as the rollout proceeds, whereas VAMO exhibits substantially slower error accumulation. Consequently, VAMO achieves lower late-time error and better full-rollout accuracy.

This distinction becomes especially clear in the more challenging settings. For Euler with $L=20$, the deterministic latent model remains competitive initially but becomes unstable during the later rollout, while VAMO maintains controlled $L^1$ and $L^2$ errors. On the 2D compressible-flow benchmark, the $L^2$ and $H^1$ errors of the direct baselines grow rapidly after the training horizon, whereas VAMO remains substantially more stable. The same pattern holds for incompressible Navier--Stokes across all viscosities and both forcing families: VAMO consistently limits the growth of field-level errors and physically relevant quantities such as enstrophy and palinstrophy. These results indicate that the primary benefit of the variational formulation lies in \textit{improving long-horizon stability} rather than uniformly reducing the error at every individual prediction step.

\begin{figure}[htbp]
    \centering
    \includegraphics[width=1.0\linewidth]{figs/euler1d_L5_metric_trends.png}
    \caption{Error trend for 1D Euler Benchmark, $L=5$}
    \label{fig:eulerl5}
\end{figure}

\begin{figure}[htbp]
    \centering
    \includegraphics[width=1.0\linewidth]{figs/euler1d_L20_metric_trends.png}
    \caption{Error trend for 1D Euler Benchmark, $L=20$}
    \label{fig:eulerl20}
\end{figure}

\begin{figure}[htbp]
    \centering
    \includegraphics[width=1.0\linewidth]{figs/cfd2d_nu8_metric_trends.png}
    \caption{Error trend for 2D CFD Benchmark}
    \label{fig:cfd2derror}
\end{figure}

\begin{figure}[htbp]
    \centering
    \includegraphics[width=0.8\linewidth]{figs/ns2d_nu3_grf_metric_trends.png}
    \caption{Error trend for 2D incompressible NS Benchmark with $\nu=10^{-3}$, GRF forcing type}
    \label{fig:nu3grf}
\end{figure}

\begin{figure}[htbp]
    \centering
    \includegraphics[width=0.80\linewidth]{figs/ns2d_nu3_sinusoidal_metric_trends.png}
    \caption{Error trend for 2D incompressible NS Benchmark with $\nu=10^{-3}$, \textsc{Wave} forcing type}
    \label{fig:nu3wave}
\end{figure}

\begin{figure}[htbp]
    \centering
    \includegraphics[width=0.80\linewidth]{figs/ns2d_nu4_grf_metric_trends.png}
    \caption{Error trend for 2D incompressible NS Benchmark with $\nu=10^{-4}$, GRF forcing type}
    \label{fig:nu4grf}
\end{figure}

\begin{figure}[htbp]
    \centering
    \includegraphics[width=0.80\linewidth]{figs/ns2d_nu4_sinusoidal_metric_trends.png}
    \caption{Error trend for 2D incompressible NS Benchmark with $\nu=10^{-4}$, \textsc{Wave} forcing type}
    \label{fig:nu4wave}
\end{figure}

\begin{figure}[htbp]
    \centering
    \includegraphics[width=0.80\linewidth]{figs/ns2d_nu5_grf_metric_trends.png}
    \caption{Error trend for 2D incompressible NS Benchmark with $\nu=10^{-5}$, GRF forcing type}
    \label{fig:nu5grf}
\end{figure}

\begin{figure}[htbp]
    \centering
    \includegraphics[width=0.8\linewidth]{figs/ns2d_nu5_sinusoidal_metric_trends.png}
    \caption{Error trend for 2D incompressible NS Benchmark with $\nu=10^{-5}$, \textsc{Wave} forcing type}
    \label{fig:nu5wave}
\end{figure}

\subsection{Physical-Statistic Stability}
\label{sec:physical_diagnostics}

We further examine whether the predicted trajectories preserve physically relevant flow statistics. Figure~\ref{fig:ns2d_physical_diagnostics} compare the enstrophy and palinstrophy of 16 representative ground-truth and predicted trajectories for each viscosity and forcing family. 

Across the six settings, the baseline failures are strongly trajectory-dependent. The largest deviations and occasional explosions occur primarily for trajectories with high ground-truth enstrophy or palinstrophy. Here, high enstrophy indicates a large overall vorticity magnitude, whereas high palinstrophy indicates strong vorticity gradients and therefore sharper, more small-scale spatial structures. These high-vorticity and high-gradient regimes are particularly difficult to propagate autoregressively. In these cases, small prediction errors can generate excessive vorticity magnitude or spurious small-scale structure, which is reflected by rapidly inflated enstrophy and, more prominently, palinstrophy.

This behavior is most severe for FNO-AE, whose predicted statistics can grow by several orders of magnitude after leaving the training horizon. Vanilla FNO and FNO+Noise are more stable overall, but still tend to overestimate the physical statistics on the most energetic or small-scale-dominated trajectories. In contrast, FNO-VAMO remains stable across both moderate and high-enstrophy or high-palinstrophy trajectories and avoids the pronounced late-time inflation observed in the baselines. Although deviations remain for individual trajectories, its predictions preserve the overall magnitude and temporal evolution of the physical diagnostics across both forcing families and all viscosity regimes. These observations complement the aggregate errors in Table~\ref{tab:ns2d_main_results}: the advantage of VAMO is not limited to lower field error, but also reflects improved control of unphysical high-frequency growth during autoregressive rollout.

\begin{figure}[p]
    \centering

    \begin{subfigure}{\linewidth}
        \centering
        \includegraphics[width=\linewidth]
        {figs/ns2d_nu3_grf_physical_diagnostics.png}
        \caption{Viscosity $\nu=10^{-3}$, \textsc{GRF} forcing.}
        \label{fig:nsdiag_nu3_grf}
    \end{subfigure}

    \vspace{0.5em}

    \begin{subfigure}{\linewidth}
        \centering
        \includegraphics[width=\linewidth]
        {figs/ns2d_nu3_sinusoidal_physical_diagnostics.png}
        \caption{Viscosity $\nu=10^{-3}$, \textsc{Wave} forcing.}
        \label{fig:nsdiag_nu3_wave}
    \end{subfigure}

    \caption{
    Physical diagnostics for the 2D incompressible
    Navier-Stokes benchmark. Each panel compares enstrophy and
    palinstrophy along 16 ground-truth and predicted trajectories.
    Solid translucent curves denote ground truth, dashed curves denote
    predictions, and the shaded region marks the training horizon.
    }
    \label{fig:ns2d_physical_diagnostics}
\end{figure}

\begin{figure}[p]
    \ContinuedFloat
    \centering

    \begin{subfigure}{\linewidth}
        \centering
        \includegraphics[width=\linewidth]
        {figs/ns2d_nu4_grf_physical_diagnostics.png}
        \caption{Viscosity $\nu=10^{-4}$, \textsc{GRF} forcing.}
        \label{fig:nsdiag_nu4_grf}
    \end{subfigure}

    \vspace{0.5em}

    \begin{subfigure}{\linewidth}
        \centering
        \includegraphics[width=\linewidth]
        {figs/ns2d_nu4_sinusoidal_physical_diagnostics.png}
        \caption{Viscosity $\nu=10^{-4}$, \textsc{Wave} forcing.}
        \label{fig:nsdiag_nu4_wave}
    \end{subfigure}

    \caption[]{Physical diagnostics for the 2D incompressible
    Navier-Stokes benchmark (continued).}
\end{figure}

\begin{figure}[p]
    \ContinuedFloat
    \centering

    \begin{subfigure}{\linewidth}
        \centering
        \includegraphics[width=\linewidth]
        {figs/ns2d_nu5_grf_physical_diagnostics.png}
        \caption{Viscosity $\nu=10^{-5}$, \textsc{GRF} forcing.}
        \label{fig:nsdiag_nu5_grf}
    \end{subfigure}

    \vspace{0.5em}

    \begin{subfigure}{\linewidth}
        \centering
        \includegraphics[width=\linewidth]
        {figs/ns2d_nu5_sinusoidal_physical_diagnostics.png}
        \caption{Viscosity $\nu=10^{-5}$, \textsc{Wave} forcing.}
        \label{fig:nsdiag_nu5_wave}
    \end{subfigure}

    \caption[]{Physical diagnostics for the 2D incompressible
    Navier-Stokes benchmark (continued).}
\end{figure}

\subsection{Ablation Studies}

\section{Conclusion}

%%%%%%%%%%%%-- reference --%%%%%%%%%%%
\newpage
\bibliographystyle{ims}
\bibliography{reference}

%%%%%%%%%%%% -- appendix -- %%%%%%%%%%%
\newpage
\appendix
\section{Gaussian Measures and Cameron-Martin Spaces}
\subsection{Gaussian Measures on Banach Spaces}
\label{app:gaussian_measures_banach}
We briefly review several basic notions for Gaussian measures on Banach spaces \citep{kuo2006gaussian}. Let $\mathcal X$ be a real separable Banach space equipped with its Borel $\sigma$-algebra, and let $\mathcal X^*$ denote its dual space. A Borel probability measure $\mu$ on $\mathcal X$ is called a \textit{Gaussian measure} if, for every continuous linear functional $f\in \mathcal X^*$, the pushforward
\[
    f_\sharp \mu=\mu\circ f^{-1}
\]
is a Gaussian probability measure on $\mathbb R$, possibly degenerate. Equivalently, if a random element $u\sim \mu$, then $f(u)$ is a real-valued Gaussian random variable for every $f\in\mathcal X^*$.

The \textit{mean} of $\mu$ is defined weakly through its action on linear functionals. For each $f\in\mathcal X^*$, one write
\[
    m_\mu(f)
    :=
    \int_{\mathcal X} f(u)\,\d\mu(u).
\]
This define a linear functional on $\mathcal X^*$, hence an element of the bidual $\mathcal X^{**}$. If there exists an element $m\in\mathcal X$ such that
\[
    f(m)
    =
    \int_{\mathcal X} f(x)\,\d\mu(x),
    \qquad \forall f\in\mathcal X^*,
\]
then $m$ is called the \textit{mean element} of $\mu$. In this case we write $m = \mathbb E_{X\sim\mu}[X]$ in the weak sense. For Gaussian measures on separable Banach spaces, the mean can typically be identified with an element of the underlying Banach space $\mathcal X$.

The \textit{covariance} of $\mu$ is also naturally defined through dual pairings. Assuming that $\mu$ has mean $m\in\mathcal X$, its \textit{covariance bilinear form} is
\[
    C_\mu(f,g)
    :=
    \int_{\mathcal X}
    \bigl(f(u)-f(m)\bigr)
    \bigl(g(u)-g(m)\bigr)
    \,\d\mu(u),
    \qquad f,g\in\mathcal X^*.
\]
Equivalently, $C_\mu(f,g)=\operatorname{Cov}_{X\sim\mu}(f(X),g(X))$.
Thus $C_\mu:\mathcal X^*\times \mathcal X^*\to\mathbb R$ is a positive semidefinite bilinear form. In particular,
\[
    C_\mu(f,f)
    =
    \operatorname{Var}_\mu(f(u))
    \ge 0.
\]
The covariance bilinear form induces a \textit{covariance operator}. Abstractly, one may define a mapping $K_\mu:\mathcal X^*\to \mathcal X^{**}$
by
\[
    (K_\mu f)(g)
    :=
    C_\mu(f,g),
    \qquad f,g\in\mathcal X^*.
\]
When $K_\mu f$ can be identified with an element of $\mathcal X$, we regard the covariance operator as a mapping $K_\mu:\mathcal X^*\to \mathcal X$
satisfying
\[
    g(K_\mu f)
    =
    C_\mu(f,g),
    \qquad f,g\in\mathcal X^*.
\]
This is the natural Banach-space analogue of the finite-dimensional covariance matrix. Indeed, if $\mathcal X=\mathbb R^d$ and $\mu=\mathsf N(m,\Sigma)$, then each $f\in\mathcal X^*$ has the form $f(u)=a^\top u$ for some $a\in\mathbb R^d$, and
\[
    C_\mu(f,g)
    =
    a^\top \Sigma b,
    \qquad
    g(u)=b^\top u.
\]
In this case the covariance operator is simply multiplication by the covariance matrix $\Sigma$.


\subsection{Abstract Wiener Space and the Cameron-Martin Theorem}
\label{app:cameron_martin}

We next review the Cameron-Martin space associated with a Gaussian measure on a Banach space. A convenient framework is provided by the notion of an abstract Wiener space \citep{stroock2010probability}.

\begin{definition}[Abstract Wiener space]
\label{def:abstract_wiener_space}
Let $\mathcal H$ be a real separable Hilbert space and let $\mathcal X$ be a real separable Banach space. Suppose that $\mathcal H$ is continuously and densely embedded into $\mathcal X$ through an injective linear map
\[
    \iota:\mathcal H\hookrightarrow \mathcal X.
\]
The triple $(\mathcal H,\mathcal X,\mu)$ is called an \emph{abstract Wiener space} if $\mu$ is a centered Gaussian measure on $\mathcal X$ whose covariance structure is induced by the embedding $\iota$ in the following sense: for every $f\in\mathcal X^*$, the pushforward $f_\sharp\mu$ is a centered real Gaussian measure satisfying
\[
    \int_{\mathcal X} f(u)^2\,\d\mu(u)
    =
    \|\iota^*f\|_{\mathcal H}^2,
\]
where $\iota^*:\mathcal X^*\to\mathcal H$ is the adjoint map defined by
\[
    \langle \iota^*f,h\rangle_{\mathcal H}
    =
    f(\iota h),
    \qquad f\in\mathcal X^*,\ h\in\mathcal H.
\]
The Hilbert space $\mathcal H$ is called the \emph{Cameron-Martin space} of $\mu$.
\end{definition}
Equivalently, the covariance operator $K_\mu:\mathcal X^*\to\mathcal X$ satisfies
\[
    K_\mu=\iota\iota^*
\]
in the weak sense that
\[
    f(K_\mu g)=
    \langle \iota^* f,\iota^* g\rangle_{\mathcal H},
    \qquad f,g\in\mathcal X^*.
\]
Thus the Cameron-Martin inner product determines the covariance structure of all continuous linear observations of the Gaussian random element.

We now define the Paley-Wiener map associated with the abstract Wiener space. For any $f,g\in\mathcal X^*$, we have
\[
    \mathbb E_{X\sim\mu}\left[f(X)g(X)\right]
    =
    \langle \iota^* f,\iota^* g\rangle_{\mathcal H}.
\]
Hence the assignment $\iota^* f\mapsto f(\cdot)$
is an isometry from the subspace $\iota^*(\mathcal X^*)\subset\mathcal H$ into $L^2(\mathcal X,\mu)$. Since $\iota^*(\mathcal X^*)$ is dense in $\mathcal H$ for an abstract Wiener space, this isometry extends uniquely to all of $\mathcal H$. This extension is called the Paley-Wiener map. For $h\in\mathcal H$, we denote its image by
\[
    \langle h,\cdot\rangle^\sim
    \in L^2(\mathcal X,\mu).
\]
Thus $\langle h,u\rangle^\sim$ is a centered Gaussian random variable satisfying
\[
    \mathbb E_\mu
    \left[
        \langle h,\cdot\rangle^\sim
        \langle k,\cdot\rangle^\sim
    \right]
    =
    \langle h,k\rangle_{\mathcal H},
    \qquad h,k\in\mathcal H.
\]
In particular, for $f\in\mathcal X^*$,
\[
    \langle \iota^*f,x\rangle^\sim
    =
    f(x),
    \qquad \mu\text{-a.s.}
\]
This identity is the main point of the construction: the Paley-Wiener map extends ordinary linear observations $f(x)$ from directions of the form $\iota^*f$ to arbitrary Cameron-Martin directions $h\in\mathcal H$.

The key property of the Cameron-Martin space is that it characterizes the directions along which the Gaussian measure can be translated without changing its measure class. For $h\in\mathcal X$, define the translated measure
\[
    \mu_h(A)
    :=
    \mu(A-h),
    \qquad A\in\mathcal B(\mathcal X).
\]
Equivalently, if $u\sim\mu$, then $u+h\sim\mu_h$.

\begin{theorem}[Cameron-Martin theorem]
\label{thm:cameron_martin}
Let $(\mathcal H,\mathcal X,\mu)$ be an abstract Wiener space. For $h\in\mathcal X$, let $\mu_h$ be the translated measure defined by $\mu_h=\mu(\cdot-h)$.
Then $\mu_h\ll\mu$ if and only if $h\in\iota(\mathcal H)$. Identifying $\mathcal H$ with its image $\iota(\mathcal H)\subset\mathcal X$, if $h\in\mathcal H$, then
\[
    \frac{\d\mu_h}{\d\mu}(u)
    =
    \exp\left(
        \langle h,u\rangle^\sim
        -
        \frac12\|h\|_{\mathcal H}^2
    \right).
\]
If $h\notin\iota(\mathcal H)$, then $\mu_h\perp\mu$.
\end{theorem}

Here the notation $\langle h,u\rangle^\sim$ should not be confused with the inner product $\langle h,u\rangle_{\mathcal H}$ in the Hilbert space. In infinite dimensions, a typical sample $u\sim\mu$ does not belong to $\mathcal H$, so $\langle h,u\rangle_{\mathcal H}$ is generally not defined. The Paley-Wiener variable $\langle h,u\rangle^\sim$ is the rigorous substitute for this formal pairing.

\subsection{Construction of the Cameron-Martin Space}
\label{app:cameron_martin_construction}
Now we construct an abstract Wiener space from an arbitrary centered Gaussian measure. Let $\mu$ be a \textit{centered} Gaussian measure on a real separable Banach space $\mathcal X$, and let $\mathcal X^*$ denote the topological dual of $\mathcal X$. Each $f\in\mathcal X^*$ can be viewed as an element of
\begin{equation}
    L^2(\mathcal{X},\mu):=\left\{f:\mathcal{X}\to\bbR\ \Big|\,\int_{\mathcal{X}}\vert f(x)\vert^2\,\d\mu(x)<\infty\right\}.
\end{equation}
Let $\mathcal G_\mu$ denote the closure of $\mathcal X^*$ in $L^2(\mathcal X,\mu)$:
\[
    \mathcal G_\mu
    :=
    \overline{\mathcal X^*}^{\,L^2(\mathcal X,\mu)}.
\]
Here $\mathcal X^*$ is identified with the collection of functions $\{u\mapsto f(u):f\in\mathcal X^*\}$. The space $\mathcal G_\mu$ is a Hilbert space with inner product inherited from $L^2(\mathcal X,\mu)$:
\[
    \langle f,g\rangle_{\mathcal G_\mu}
    :=\int f(x)g(x)\,\d\mu(x),
    \qquad f,g\in\mathcal G_\mu.
\]
This space is often called the first Wiener chaos associated with $\mu$. We now map $\mathcal G_\mu$ back into the Banach space $\mathcal X$. For $f\in\mathcal G_\mu$, define
\[
R_\mu f:=\int_{\mathcal X} xf(x)\,\d\mu(x),
\]
where the integral is understood weakly, namely
\[
    g(R_\mu f)
    =\int_\mathcal{X} g(x)f(x)\,\d\mu(x),
    \qquad \forall g\in\mathcal X^*.
\]
For $f\in\mathcal X^*$, this map satisfies
\[
    R_\mu f = K_\mu f,
\]
where $K_\mu:\mathcal X^*\to\mathcal X$ is the covariance operator. The \textit{Cameron-Martin space} of $\mu$ is defined as
\[
    \mathcal H_\mu
    :=
    R_\mu(\mathcal G_\mu)\subset \mathcal X.
\]
It becomes a Hilbert space by transporting the $L^2(\mu)$ inner product from $\mathcal G_\mu$:
\[
    \langle R_\mu f,R_\mu g\rangle_{\mathcal H_\mu}
    :=
    \int_\mathcal{X} f(x)g(x)\,\d\mu(x).
\]
Equivalently, $\mathcal H_\mu$ is the completion of $K_\mu(\mathcal X^*)$ under the inner product
\[
    \langle K_\mu f,K_\mu g\rangle_{\mathcal H_\mu}
    :=\int_\mathcal{X}f(x)g(x)\,\d\mu(x).
\]
Thus $K_\mu(\mathcal X^*)$ is dense in $\mathcal H_\mu$. The inverse map
\[
    R_\mu^{-1}:\mathcal H_\mu\to \mathcal G_\mu\subset L^2(\mathcal X,\mu)
\]
is called the \textit{Paley-Wiener map}. For $h\in\mathcal H_\mu$, we denote its Paley-Wiener image by
\[
    \langle h,\cdot\rangle^\sim
    :=
    R_\mu^{-1}h.
\]
In particular, for $h=K_\mu f$ with $f\in\mathcal X^*$, one has
\[
    \langle K_\mu f,u\rangle^\sim
    =
    f(u)
    \qquad \mu\text{-a.s.}
\]
This identity explains why the Paley-Wiener map can be regarded as the rigorous version of the formal pairing $\langle h,u\rangle_{\mathcal H_\mu}$. In infinite dimensions, a typical sample $u\sim\mu$ does not belong to $\mathcal H_\mu$, so the Hilbert-space inner product $\langle h,u\rangle_{\mathcal H_\mu}$ is generally not defined.

The construction above turns $(\mathcal H_\mu,\mathcal X,\mu)$ into an abstract Wiener space: $\mathcal H_\mu$ is continuously embedded into $\mathcal X$, and $\mu$ is the Gaussian measure on $\mathcal X$ whose covariance is induced by this embedding. The inclusion map $\iota$ is simply the canonical embedding $\iota:\mathcal{H}_\mu\ni x\mapsto x\in\mathcal{X}$, with the adjoint map given by $$\iota^* f=K_\mu f,\quad f\in\mathcal{X}^*.$$  The following result is a reproduction of Theorem \ref{thm:cameron_martin} in this specific setting.

\begin{theorem}[Cameron-Martin theorem]
\label{thm:cameron_martin2}
Let $\mu$ be a centered Gaussian measure on a real, separable Banach space $\mathcal X$, and let $\mathcal H_\mu$ be its Cameron-Martin space. For $h\in\mathcal{X}$,
the translated measure $\mu_h\ll \mu$ if and only if $h\in\mathcal{H}_\mu$. Moreover, if $h\in\mathcal H_\mu$, then
\[
    \frac{\d\mu_h}{\d\mu}(u)=\exp\left(\langle h,u\rangle^\sim-\frac12\|h\|_{\mathcal H_\mu}^2\right).
\]
If $h\notin\mathcal H_\mu$, then $\mu_h\perp \mu$.
\end{theorem}

For a non-centered Gaussian measure with mean $m\in\mathcal X$ and the same covariance structure, the Cameron-Martin space is unchanged. In that case, for $h\in\mathcal H_\mu$,
\[
    \frac{\d\mathsf N(m+h,K_\mu)}{\d\mathsf N(m,K_\mu)}(u)
    =
    \exp\left(
        \langle h,u-m\rangle^\sim
        -
        \frac12\|h\|_{\mathcal H_\mu}^2
    \right).
\]
Thus the Cameron-Martin norm is the infinite-dimensional analogue of the Mahalanobis norm. In particular, shifts outside $\mathcal H_\mu$ produce singular Gaussian measures, while shifts inside $\mathcal H_\mu$ preserve absolute continuity.


\subsection{Examples: Gaussian Measures on Hilbert and Sobolev Spaces}
\label{app:gaussian_examples}

We give several examples to illustrate the preceding abstract definitions. The simplest case is a Gaussian measure on a finite-dimensional Euclidean space. If $\mathcal X=\mathbb R^d$ and
\[
    \mu=\mathsf N(0,\Sigma),
\]
where $\Sigma$ is positive definite, then the Cameron-Martin space is all of $\mathbb R^d$, equipped with the inner product
\[
    \langle h,k\rangle_{\mathcal H_\mu}
    =
    h^\top \Sigma^{-1}k.
\]
Thus the Cameron-Martin norm is the usual Mahalanobis norm,
\[
    \|h\|_{\mathcal H_\mu}^2
    =
    h^\top \Sigma^{-1}h.
\]
The Cameron-Martin theorem then reduces to the standard likelihood-ratio formula
\[
    \frac{d\mathsf N(h,\Sigma)}{d\mathsf N(0,\Sigma)}(u)
    =
    \exp\left(
        h^\top\Sigma^{-1}u
        -
        \frac12 h^\top\Sigma^{-1}h
    \right).
\]

We next consider an infinite-dimensional Hilbert space example. Let $\mathcal X=L^2(\Omega)$, where $\Omega\subset\mathbb R^d$ is a bounded domain or a periodic domain. Let $Q:L^2(\Omega)\to L^2(\Omega)$ be a positive, self-adjoint, trace-class covariance operator. Suppose that $Q$ has eigenpairs
\[
    Q e_k = \lambda_k e_k,
    \qquad \lambda_k>0,
\]
where $\{e_k\}_{k\ge 1}$ is an orthonormal basis of $L^2(\Omega)$. Then a centered Gaussian random element with covariance $Q$ admits the formal Karhunen-Loève expansion
\[
    x=
    \sum_{k=1}^\infty\sqrt{\lambda_k}\,\xi_k e_k,
    \qquad
    \xi_k\stackrel{\mathrm{i.i.d.}}{\sim}\mathsf N(0,1).
\]
Since $Q$ is trace-class, the above series converges in $L^2$ sense:
\begin{equation*}
    \lim_{N\to\infty}\bbE\left\Vert x-\sum_{k=1}^N\sqrt{\lambda_k}\,\xi_k e_k\right\Vert^2_\mathcal{X}=0.
\end{equation*}
The Cameron-Martin space is
\[
    \mathcal H_\mu
    =
    \operatorname{Range}(Q^{1/2}),
\]
with norm
\[
    \|h\|_{\mathcal H_\mu}^2
    =
    \|Q^{-1/2}h\|_{L^2}^2.
\]
Writing $h=\sum_{k\ge 1} h_k e_k$, we have $h\in\mathcal H_\mu$ if and only if $\sum_{k\ge 1}|h_k|^2/\lambda_k<\infty$, and
\[
    \|h\|_{\mathcal H_\mu}^2
    =
    \sum_{k\ge 1}\frac{|h_k|^2}{\lambda_k}.
\]
Therefore directions with significant high-frequency components are penalized more strongly when the covariance eigenvalues $\lambda_k$ decay rapidly.

A particularly important class of examples is obtained from elliptic differential operators. Let
\[
    Q=(-\Delta+\tau I)^{-s},
    \qquad \tau>0,\ s>0,
\]
with suitable boundary conditions or periodic boundary conditions. If
\[
    (-\Delta+\tau I)e_k=\rho_k e_k,
\]
then
\[
    Qe_k=\rho_k^{-s}e_k.
\]
Thus the Gaussian random field has the formal expansion
\[
    u
    =
    \sum_{k\ge 1}\rho_k^{-s/2}\xi_k e_k,
    \qquad
    \xi_k\stackrel{\mathrm{i.i.d.}}{\sim}\mathsf N(0,1).
\]
The Cameron-Martin norm is
\[
    \|h\|_{\mathcal H_\mu}^2
    =
    \|Q^{-1/2}h\|_{L^2}^2
    =
    \|(-\Delta+\tau I)^{s/2}h\|_{L^2}^2.
\]
Hence, up to equivalence of norms,
\[
    \mathcal H_\mu \simeq H^s(\Omega).
\]
In this sense, the Cameron-Martin space associated with the covariance operator $(-\Delta+\tau I)^{-s}$ is a Sobolev space of order $s$.

For example, on the torus $\Omega=\mathbb T^d$, one may write
\[
    h(x)
    =
    \sum_{k\in\mathbb Z^d} h_k e^{2\pi i k\cdot x}.
\]
Since
\[
    (-\Delta+\tau I)e^{2\pi i k\cdot x}
    =
    (4\pi^2|k|^2+\tau)e^{2\pi i k\cdot x},
\]
the Cameron-Martin norm becomes
\[
    \|h\|_{\mathcal H_\mu}^2
    =
    \sum_{k\in\mathbb Z^d}
    (4\pi^2|k|^2+\tau)^s |h_k|^2.
\]
This is equivalent to the standard Sobolev $H^s(\mathbb T^d)$ norm.

It is important to distinguish the Cameron-Martin space from the typical sample space of the Gaussian measure. If
\[
    u
    =
    \sum_{k\in\mathbb Z^d}
    (4\pi^2|k|^2+\tau)^{-s/2}\xi_k e^{2\pi i k\cdot x},
\]
then
\[
    \mathbb E\|u\|_{H^r}^2
    \asymp
    \sum_{k\in\mathbb Z^d}
    (1+|k|^2)^r(1+|k|^2)^{-s}.
\]
This sum is finite when
\[
    r<s-\frac d2.
\]
Thus a typical sample belongs to $H^r(\Omega)$ for $r<s-d/2$, up to the usual endpoint caveat, whereas the Cameron-Martin space is $H^s(\Omega)$. Therefore typical samples are rougher than Cameron-Martin directions.

As a classical example, Wiener measure on $\mathcal X=C_0([0,1])$ has Cameron-Martin space
\[
    \mathcal H_\mu
    =
    \left\{
        h\in C_0([0,1]):
        h \text{ is absolutely continuous and } h'\in L^2([0,1])
    \right\},
\]
with inner product
\[
    \langle h,k\rangle_{\mathcal H_\mu}
    =
    \int_0^1 h'(t)k'(t)\,dt.
\]
Thus $\mathcal H_\mu$ is the Sobolev space $H_0^1([0,1])$ embedded into $C_0([0,1])$. Brownian paths themselves are almost surely nowhere differentiable and hence do not belong to $\mathcal H_\mu$. This example illustrates the general principle that the Cameron-Martin space consists of admissible deterministic shifts, not typical random samples.

\section{Proofs}
\subsection{Proof of Proposition \ref{prop:conditional_functional_elbo}}
\begin{proof}[Proof of Proposition \ref{prop:conditional_functional_elbo}]
(a) Define the latent predictive kernel
\begin{equation*}
    K_{\theta,\phi}(\d\mathbf{Z}_{n+1},\d\mathbf{Z}_n\,|\,\mathbf{U}_n)=P_\theta(\d\mathbf{Z}_{n+1}\,|\,\mathbf{Z}_n)\,Q_\phi(\d\mathbf{Z}_n\,|\,\mathbf{U}_n).
\end{equation*}
Then
\begin{equation*}
    P_{\theta,\psi,\phi}(\cdot\,|\,\mathbf{U}_n)=\int_{\mathcal{Z}\times\mathcal{Z}} P_\psi(\cdot\,|\,\mathbf{Z}_{n+1})\,K_{\theta,\phi}(\d\mathbf{Z}_{n+1},\d\mathbf{Z}_n\,|\,\mathbf{U}_n),
\end{equation*}
and $\mathsf{W}_\mathcal{U}$-almost every $\mathbf{U}_{n+1}$, we have
\begin{equation*}
    \frac{\d P_{\theta,\psi,\phi}(\cdot\,|\,\mathbf{U}_n)}{\d\mathsf{W}_\mathcal{U}}(\mathbf{U}_{n+1})=\int_{\mathcal{Z}\times\mathcal{Z}} \frac{\d P_\psi(\cdot\,|\,\mathbf{Z}_{n+1})}{\d\mathsf{W}_\mathcal{U}}(\mathbf{U}_{n+1})\,K_{\theta,\phi}(\d\mathbf{Z}_{n+1},\d\mathbf{Z}_n\,|\,\mathbf{U}_n).
\end{equation*}
By Jensen's inequality,
\begin{align*}
    \log \frac{\d P_{\theta,\psi,\phi}(\cdot\,|\,\mathbf{U}_n)}{\d\mathsf{W}_\mathcal{U}}(\mathbf{U}_{n+1})&\geq\int_{\mathcal{Z}\times\mathcal{Z}} \log \frac{\d P_\psi(\cdot\,|\,\mathbf{Z}_{n+1})}{\d\mathsf{W}_\mathcal{U}}(\mathbf{U}_{n+1})\,K_{\theta,\phi}(\d\mathbf{Z}_{n+1},\d\mathbf{Z}_n\,|\,\mathbf{U}_n)\\
    &=\bbE_{\mathbf{Z}_{n+1},\mathbf{Z}_n\sim K_{\theta,\phi}(\cdot|\mathbf{U}_n)}\left[\log \frac{\d P_\psi(\cdot\,|\,\mathbf{Z}_{n+1})}{\d\mathsf{W}_\mathcal{U}}(\mathbf{U}_{n+1})\right].
\end{align*}
Expanding $K_{\theta,\phi}$ gives \eqref{eq:predlb}.
\item (b) For fixed $\mathbf Z_n\in\mathcal{Z}$, define the latent-conditioned predictive distribution
\begin{equation*}
    P_{\theta,\psi}(\d\mathbf U_{n+1}\mid \mathbf Z_n)
    =
    \int_{\mathcal Z}
    P_\psi(\d\mathbf U_{n+1}\mid \mathbf Z_{n+1})
    P_\theta(\d\mathbf Z_{n+1}\mid \mathbf Z_n).
\end{equation*}
Then we have the decomposition
\begin{equation}
    \frac{\d P_{\theta,\psi,\phi}(\cdot\,|\,\mathbf{U}_n)}{\d\mathsf{W}_\mathcal{U}}(\mathbf{U}_{n+1})=\int_{\mathcal{Z}}\frac{\d P_{\theta,\psi}(\cdot\,|\,\mathbf{Z}_n)}{\d\mathsf{W}_\mathcal{U}}(\d\mathbf{U}_{n+1})\,Q_\phi(\d\mathbf{Z}_n\,|\,\mathbf{U}_n),\quad\text{for $\mathsf{W}_\mathcal{U}$-a.e.}\ \mathbf{U}_{n+1}.\label{eq:vaedecomp}
\end{equation}
We choose a reference measure $\Lambda$ on $\mathcal{Z}$ such that both $Q_\phi(\cdot\,|\,\mathbf{U}_{n+1})$ and $P_\theta(\cdot\,|\,\mathbf{Z}_n)$ are absolutely continuous with respect to $\Lambda$, for instance, $\Lambda=Q_\phi(\cdot\,|\,\mathbf{U}_{n+1})/2+P_\theta(\cdot\,|\,\mathbf{Z}_n)/2$, and define the likelihood densities
$$p_\theta=\frac{\d P_\theta(\cdot\,|\,\mathbf{Z}_n)}{\d\Lambda},\quad q_\phi=\frac{\d Q_\theta(\cdot\,|\,\mathbf{U}_{n+1})}{\d\Lambda}.$$
Then we factorize $P_{\theta,\psi}$ and apply Jensen's inequality to obtain
\begin{align*}
    \log \frac{\d P_{\theta,\psi}(\cdot\,|\,\mathbf{Z}_n)}{\d\mathsf{W}_\mathcal{U}}&(\mathbf{U}_{n+1})
    =
    \log\int_{\mathcal Z}
    \frac{\d P_\psi(\cdot\,|\, \mathbf Z_{n+1})}{\d\mathsf{W}_\mathcal{U}}(\mathbf U_{n+1})\,
    p_\theta(\mathbf Z_{n+1}\,|\, \mathbf Z_n)\,\Lambda(\d\mathbf{Z}_{n+1})\\
    &=\log\int_{\mathcal Z}
    \frac{\d P_\psi(\cdot\,|\, \mathbf Z_{n+1})}{\d\mathsf{W}_\mathcal{U}}(\mathbf U_{n+1})\,
    \frac{p_\theta(\mathbf Z_{n+1}\,|\, \mathbf Z_n)}{q_\phi(\mathbf Z_{n+1}\,|\, \mathbf U_{n+1})}\,q_\phi(\mathbf Z_{n+1}\,|\, \mathbf U_{n+1})\Lambda(\d\mathbf{Z}_{n+1})\\
    &\geq\int_{\mathcal Z}
    \left[\log \frac{\d P_\psi(\cdot\,|\, \mathbf Z_{n+1})}{\d\mathsf{W}_\mathcal{U}}(\mathbf U_{n+1})-\log
    \frac{q_\phi(\mathbf Z_{n+1}\,|\, \mathbf U_{n+1})}{p_\theta(\mathbf Z_{n+1}\,|\, \mathbf Z_n)}\right]\! q_\phi(\mathbf Z_{n+1}\,|\,\mathbf U_{n+1})\Lambda(\d\mathbf{Z}_{n+1}),
\end{align*}
with the convention $q_\phi(\mathbf Z_{n+1}\,|\, \mathbf U_{n+1})\log q_\phi(\mathbf Z_{n+1}\,|\, \mathbf U_{n+1})=0$ outside the support of $Q_\phi(\cdot\,|\, \mathbf U_{n+1})$. For $Q_\phi(\cdot\,|\,\mathbf{U}_{n+1})\ll P_\theta(\cdot\,|\,\mathbf{Z}_n)$, the above result is equivalent to
\begin{align*}
    \log \frac{\d P_{\theta,\psi}(\cdot\,|\,\mathbf{Z}_n)}{\d\mathsf{W}_\mathcal{U}}(\mathbf{U}_{n+1})&\geq\int_{\mathcal Z}
    \left[\log \frac{\d P_\psi(\cdot\,|\, \mathbf Z_{n+1})}{\d\mathsf{W}_\mathcal{U}}(\mathbf U_{n+1})-\log
    \frac{\d Q_\phi(\cdot\,|\, \mathbf U_{n+1})}{P_\theta(\cdot\,|\, \mathbf Z_n)}(\mathbf Z_{n+1})\right] Q_\phi(\d\mathbf Z_{n+1}\,|\, \mathbf U_{n+1})\\
    &=\bbE_{\mathbf{Z}_{n+1}\sim Q_\phi(\cdot\,|\,\mathbf{U}_{n+1})}\left[ \jg{\log} \frac{\d P_\psi(\cdot\,|\, \mathbf Z_{n+1})}{\d\mathsf{W}_\mathcal{U}}(\mathbf U_{n+1})\right]-D_\mathrm{KL}\left(Q_\phi(\cdot\,|\,\mathbf{U}_{n+1})\,\|\,P_\theta(\cdot\,|\,\mathbf{Z}_n)\right).
\end{align*}
Again we use Jensen's inequality to \eqref{eq:vaedecomp}:
\begin{equation*}
    \log \frac{\d P_{\theta,\psi,\phi}(\cdot\,|\,\mathbf{U}_n)}{\d\mathsf{W}_\mathcal{U}}(\mathbf{U}_{n+1})\geq\bbE_{\mathbf{Z}_n\sim Q_\phi(\cdot|\mathbf{U}_n)}\left[\log \frac{\d P_{\theta,\psi}(\cdot\,|\,\mathbf{Z}_n)}{\d\mathsf{W}_\mathcal{U}}(\mathbf{U}_{n+1})\right].
\end{equation*}
Combining the two inequalities gives
\begin{align*}
    \log \frac{\d P_{\theta,\psi,\phi}(\cdot\,|\,\mathbf{U}_n)}{\d\mathsf{W}_\mathcal{U}}(\mathbf{U}_{n+1})&\geq\bbE_{\mathbf{Z}_{n+1}\sim Q_\phi(\cdot\,|\,\mathbf{U}_{n+1})}\left[
    \jg{\log} \frac{\d P_\psi(\cdot\,|\, \mathbf Z_{n+1})}{\d\mathsf{W}_\mathcal{U}}(\mathbf U_{n+1})\right]\\
    &\quad -\bbE_{\mathbf{Z}_n\sim Q_\phi(\cdot|\mathbf{U}_n)}\left[D_\mathrm{KL}\left(Q_\phi(\cdot\,|\,\mathbf{U}_{n+1})\,\|\,P_\theta(\cdot\,|\,\mathbf{Z}_n)\right)\right],
\end{align*}
which proves the claim.
\end{proof}
\subsection{Proof of Proposition \ref{prop:population_elbo_transition_consistency}}
\begin{proof}[Proof of Proposition \ref{prop:population_elbo_transition_consistency}]
As in the proof of Proposition~\ref{prop:conditional_functional_elbo},
we choose a common dominating measure on $\mathcal Z$ as reference and, with a slight abuse of notation, write
$Q_\phi(\mathbf Z^+\,|\, \mathbf U^+)$,
$P(\mathbf Z^+\,|\, \mathbf Z)$, and
$\Pi_\phi(\mathbf Z^+\,|\, \mathbf Z)$ for the corresponding densities with respect to the dominating measure. Expanding the KL divergence in
$\mathcal R(P)$ gives
\begin{equation*}
\begin{aligned}
\mathcal R(P)&=
\mathbb E_{(\mathbf U,\mathbf U^+)\sim\rho}
\mathbb E_{\mathbf Z\sim Q_\phi(\cdot\,|\,\mathbf U)}
\mathbb E_{\mathbf Z^+\sim Q_\phi(\cdot\,|\,\mathbf U^+)}
\left[\log
\frac{\d Q_\phi(\cdot\,|\,\mathbf U^+)
}{\d P(\cdot\,|\,\mathbf Z)}
(\mathbf Z^+)\right]
\\
&=\mathbb E_{(\mathbf Z,\mathbf Z^+,\mathbf{U}^+)\sim\Gamma_\phi}
\left[
\log \frac{\d Q_\phi(\cdot\,|\,\mathbf U^+)}{\d\Pi_\phi(\cdot\,|\,\mathbf{Z})}(\mathbf Z^+)
\right]+
\mathbb E_{(\mathbf Z,\mathbf Z^+)\sim\Gamma_\phi}
\left[
\log\frac{\d\Pi_\phi(\cdot\,|\,\mathbf{Z})}{\d P(\cdot\,|\,\mathbf Z)}(\mathbf Z^+)
\right].
\end{aligned}
\end{equation*}
The first term depends only on the encoder and the data distribution, and is therefore independent of $P$. For the second term, disintegrate $\Gamma_\phi$ as
\[
\Gamma_\phi(\d\mathbf Z,\d\mathbf Z^+,\d\mathbf{U}^+)=
\Gamma_\phi(\d\mathbf{Z},\d\mathbf{U}^+)\,Q_\phi(\d\mathbf{Z}^+|\mathbf{U}^+)
\]
Then
\begin{equation}
\begin{aligned}
&\mathbb E_{(\mathbf Z,\mathbf Z^+,\mathbf{U}^+)\sim\Gamma_\phi}
\left[
\log \frac{\d Q_\phi(\cdot\,|\,\mathbf U^+)}{\d \Pi_\phi(\cdot\,|\,\mathbf{Z})}(\mathbf Z^+)
\right]\\
&\quad=\int_{\mathcal{U}\times\mathcal{U}}\int_\mathcal{Z} D_\mathrm{KL}\left(Q_\phi(\cdot\,|\,\mathbf{U}^+)\,\|\,\Pi_\phi(\cdot\,|\,\mathbf{Z})\right) Q_\phi(\d\mathbf{Z}\,|\,\mathbf{U})\,\rho(\d\mathbf{U},\d\mathbf{U}^+).
\end{aligned}\label{eq:elbokldecompaux1}
\end{equation}
For the second term, disintegrate $\Gamma_\phi$ as
\[
\Gamma_\phi(\d\mathbf Z,\d\mathbf Z^+)=\Pi_\phi(\d\mathbf Z^+\,|\,\mathbf Z)\,\Gamma_\phi(\d\mathbf Z).
\]
Then
\begin{equation}
\begin{aligned}
\mathbb E_{(\mathbf Z,\mathbf Z^+)\sim\Gamma_\phi}
\left[
\log\frac{\d\Pi_\phi(\cdot\,|\,\mathbf{Z})}{\d P(\cdot\,|\,\mathbf Z)}(\mathbf Z^+)
\right]=\int_{\mathcal Z}
D_{\mathrm{KL}}
\left(\Pi_\phi(\cdot\,|\,\mathbf Z)\,\Vert\,P(\cdot\,|\,\mathbf Z)\right)
\Gamma_\phi(\d\mathbf Z),
\end{aligned}\label{eq:elbokldecompaux2}
\end{equation}
Combining the identities \eqref{eq:elbokldecompaux1} and \eqref{eq:elbokldecompaux2} proves \eqref{eq:elbokldecomp}.

The first term in \eqref{eq:elbokldecomp} is independent of $P$. The second term is nonnegative and equals zero if and only if $P(\cdot\,|\,\mathbf Z)=\Pi_\phi(\cdot\,|\,\mathbf Z)$ for
$\Gamma_\phi$-almost every $\mathbf Z$. Therefore $P^\star=\Pi_\phi$ is the population minimizer over all Markov kernels.
\end{proof}

\subsection{Proof of Proposition \ref{prop:cm_decoder_likelihood}}
\begin{proof}[Proof of Proposition \ref{prop:cm_decoder_likelihood}]
Since $D_\psi(\mathbf Z)\in \mathcal H_{\mathcal U}$, the
Cameron-Martin theorem [Theorem \ref{thm:cameron_martin2}] implies that the shifted Gaussian measure $\mathsf W_{\mathcal U}^{D_\psi(\mathbf Z)}$ is absolutely continuous with respect to $\mathsf W_{\mathcal U}$ and satisfies
\[
    \frac{\d\mathsf W_{\mathcal U}^{D_\psi(\mathbf Z)}}{\d\mathsf W_{\mathcal U}}(\mathbf U)=\exp\left(\langle D_\psi(\mathbf Z),\mathbf U\rangle^\sim
    -\frac12 \|D_\psi(\mathbf Z)\|_{\mathcal H_{\mathcal U}}^2
    \right).
\]
This proves \eqref{eq:cm_decoder_rn}. Taking negative logarithms gives \eqref{eq:cm_decoder_nll}. If $\mathbf U\in\mathcal H_{\mathcal U}$, then the Paley-Wiener map agrees with the Cameron-Martin inner product:
\[
    \langle D_\psi(\mathbf Z),\mathbf U\rangle^\sim
    =
    \langle D_\psi(\mathbf Z),\mathbf U\rangle_{\mathcal H_{\mathcal U}}.
\]
Hence
\[
\begin{aligned}
-\log\frac{\d P_\psi(\cdot\mid \mathbf Z)}{\d\mathsf W_{\mathcal U}}(\mathbf U)
&=\frac12
\|D_\psi(\mathbf Z)\|_{\mathcal H_{\mathcal U}}^2-
\langle D_\psi(\mathbf Z),\mathbf U\rangle_{\mathcal H_{\mathcal U}}\\
&=\frac12\|\mathbf U-D_\psi(\mathbf Z)\|_{\mathcal H_{\mathcal U}}^2-\frac12\|\mathbf U\|_{\mathcal H_{\mathcal U}}^2,
\end{aligned}
\]
which proves \eqref{eq:cm_decoder_squared_norm}.
\end{proof}

\begin{lemma}[Gaussian measure induced by a trace-class covariance]
\label{lem:trace_class_gaussian_cm_kl}
Let $\mathcal Z$ be a separable Hilbert space and let $K:\mathcal Z\to\mathcal Z$ be a positive, self-adjoint, trace-class operator. Let $\{(\lambda_i,\mathbf E_i)\}_{i\ge 1}$ be an eigensystem of $K$, where $\{\mathbf E_i\}_{i\ge 1}$ is an orthonormal basis of $\mathcal Z$,
$\lambda_i>0$, and
\[
    \operatorname{Tr}(K)
    =
    \sum_{i=1}^\infty \lambda_i
    <\infty .
\]
Let $\{\xi_i\}_{i\ge 1}$ be i.i.d. standard Gaussian random variables. Then the series
\[
    \mathbf G
    =
    \sum_{i=1}^\infty \sqrt{\lambda_i}\xi_i\mathbf E_i
\]
converges in $L^2(\Omega;\mathcal Z)$ and almost surely in $\mathcal Z$. Moreover, $\mathbf G$ is a centered $\mathcal Z$-valued Gaussian random element with covariance operator $K$, i.e.,
\[
    \mathbf G\sim \mathsf N(0,K).
\]
The Cameron-Martin space of $\mathsf N(0,K)$ is
\[
    \mathcal H_K=
    \left\{\mathbf h=\sum_{i=1}^\infty h_i\mathbf E_i:
        \sum_{i=1}^\infty \frac{|h_i|^2}{\lambda_i}<\infty
    \right\},
\]
equipped with the inner product
\[
\langle \mathbf h,\mathbf g\rangle_{\mathcal H_K} =\sum_{i=1}^\infty \frac{h_i g_i}{\lambda_i},
\]
and norm
\[
    \|\mathbf h\|_{\mathcal H_K}^2=\sum_{i=1}^\infty \frac{|h_i|^2}{\lambda_i}.
\]
Equivalently,
\[
    \mathcal H_K=\operatorname{Range}(K^{1/2}),
    \qquad
    \|\mathbf h\|_{\mathcal H_K}
    =
    \|K^{-1/2}\mathbf h\|_{\mathcal Z},
\]
where $K^{-1/2}$ is understood on $\operatorname{Range}(K^{1/2})$.
\end{lemma}

\begin{proof}
\textit{Step I: Convergence.} For each positive integer $N$, we define
\[
    \mathbf G_N
    =
    \sum_{i=1}^N\sqrt{\lambda_i}\xi_i\mathbf E_i .
\]
For $M>N$, by orthonormality of $\{\mathbf E_i\}_{i\ge 1}$ and independence of
$\{\xi_i\}_{i\ge 1}$,
\[
    \mathbb E\|\mathbf G_M-\mathbf G_N\|_{\mathcal Z}^2
    =
    \mathbb E
    \left\|
        \sum_{i=N+1}^M\sqrt{\lambda_i}\xi_i\mathbf E_i
    \right\|_{\mathcal Z}^2
    =
    \sum_{i=N+1}^M\lambda_i .
\]
Since $\sum_{i=1}^\infty \lambda_i<\infty$, the sequence
$\{\mathbf G_N\}_{N\ge 1}$ is Cauchy in $L^2(\Omega;\mathcal Z)$.
Therefore it converges in $L^2(\Omega;\mathcal Z)$ to some $\mathcal Z$-valued random element $\mathbf G$. In particular,
\[
    \mathbb E\|\mathbf G\|_{\mathcal Z}^2
    =
    \lim_{N\to\infty}\mathbb E\|\mathbf G_N\|_{\mathcal Z}^2
    =
    \sum_{i=1}^\infty \lambda_i
    =
    \operatorname{Tr}(K)
    <\infty.
\]
Moreover, $\{\mathbf G_N\}_{N\ge 1}$ is an $L^2$-bounded martingale with respect to the filtration
\[
    \mathcal F_N=\sigma(\xi_1,\ldots,\xi_N).
\]
Indeed, for $M>N$,
\[
    \mathbb E[\mathbf G_M\mid \mathcal F_N]=\mathbf G_N,
\]
and
\[
    \sup_{N\ge 1}\mathbb E\|\mathbf G_N\|_{\mathcal Z}^2
    =
    \sup_{N\ge 1}\sum_{i=1}^N\lambda_i
    \leq
    \sum_{i=1}^\infty \lambda_i
    <\infty.
\]
Therefore, by the Hilbert-space martingale convergence theorem, $\mathbf G_N$ converges almost surely and in $L^2(\Omega;\mathcal Z)$ to a
$\mathcal Z$-valued random element. The $L^2$ limit is unique, so the almost sure limit is the same random element $\mathbf G$.

\item\textit{Step II: Verify the Karhunen-Loève expansion.} We next verify that the random element $\mathbf G$ has covariance operator $K$. For any $\mathbf f\in\mathcal Z$,
\[
    \langle \mathbf G,\mathbf f\rangle_{\mathcal Z}
    =
    \sum_{i=1}^\infty \sqrt{\lambda_i}\xi_i
    \langle \mathbf E_i,\mathbf f\rangle_{\mathcal Z},
\]
where the convergence holds in $L^2(\Omega)$. Hence
$\langle \mathbf G,\mathbf f\rangle_{\mathcal Z}$ is a centered Gaussian random variable. Therefore $\mathbf G$ is a centered Gaussian random element. Moreover,
for any $\mathbf f,\mathbf g\in\mathcal Z$,
\[
\begin{aligned}
    \mathbb E
    \left[
        \langle \mathbf G,\mathbf f\rangle_{\mathcal Z}
        \langle \mathbf G,\mathbf g\rangle_{\mathcal Z}
    \right]
    &=
    \sum_{i=1}^\infty \lambda_i
    \langle \mathbf E_i,\mathbf f\rangle_{\mathcal Z}
    \langle \mathbf E_i,\mathbf g\rangle_{\mathcal Z}  \\
    &=
    \left\langle
        \sum_{i=1}^\infty \lambda_i
        \langle \mathbf f,\mathbf E_i\rangle_{\mathcal Z}\mathbf E_i,
        \mathbf g
    \right\rangle_{\mathcal Z}=
    \langle K\mathbf f,\mathbf g\rangle_{\mathcal Z}.
\end{aligned}
\]
Thus the covariance operator of $\mathbf G$ is $K$, and $\mathbf G\sim\mathsf N(0,K)$.

\item\textit{Step III: Identify the Cameron-Martin space.} Since $K$ is positive and self-adjoint,
\[
    K^{1/2}\mathbf f
    =
    \sum_{i=1}^\infty \sqrt{\lambda_i}
    \langle \mathbf f,\mathbf E_i\rangle_{\mathcal Z}\mathbf E_i .
\]
Therefore $\mathbf h\in \operatorname{Range}(K^{1/2})$ if and only if there
exists $\mathbf f=\sum_{i=1}^\infty f_i\mathbf E_i\in\mathcal Z$ such that
\[
    \mathbf h=K^{1/2}\mathbf f
    =
    \sum_{i=1}^\infty \sqrt{\lambda_i}f_i\mathbf E_i .
\]
Writing $h_i=\sqrt{\lambda_i}f_i$, this is equivalent to
\[
    \sum_{i=1}^\infty \frac{|h_i|^2}{\lambda_i}
    =
    \sum_{i=1}^\infty |f_i|^2
    <\infty .
\]
Hence
\[
    \mathcal H_K
    =
    \operatorname{Range}(K^{1/2})
    =
    \left\{
        \mathbf h=\sum_{i=1}^\infty h_i\mathbf E_i:
        \sum_{i=1}^\infty \frac{|h_i|^2}{\lambda_i}<\infty
    \right\}.
\]
The Cameron-Martin inner product is the pullback inner product through
$K^{1/2}$:
\[
    \langle \mathbf h,\mathbf g\rangle_{\mathcal H_K}
    =
    \langle K^{-1/2}\mathbf h,K^{-1/2}\mathbf g\rangle_{\mathcal Z}
    =
    \sum_{i=1}^\infty \frac{h_i g_i}{\lambda_i}.
\]
This proves the claim.
\end{proof}

\begin{proof}[Proof of Proposition \ref{prop:general_spectral_residual_control}]
Part~(a) follows directly from Lemma~\ref{lem:trace_class_gaussian_cm_kl}.
Indeed, since $K$ is positive, self-adjoint, and trace-class, the centered
Gaussian measure $\mathsf W_{\mathcal Z}=\mathsf N(0,K)$ is well-defined on
$\mathcal Z$, and its Cameron-Martin space is
\[
    \mathcal H_K
    =
    \operatorname{Range}(K^{1/2})
    =
    \left\{
        \mathbf R=\sum_{i=1}^\infty r_i\mathbf E_i:
        \sum_{i=1}^\infty \frac{|r_i|^2}{\lambda_i}<\infty
    \right\},
\]
with norm
\[
    \|\mathbf R\|_{\mathcal H_K}^2
    =
    \sum_{i=1}^\infty \frac{|r_i|^2}{\lambda_i}.
\]
We now prove part~(b). Fix $\mathbf Z$ and write
\[
    a=\alpha_\theta(\mathbf Z),
    \qquad
    m=T_\theta(\mathbf Z).
\]
Then
\[
    P_\theta(\cdot\mid \mathbf Z)=\mathsf N(m,a^2K).
\]
The Cameron-Martin space of $\mathsf N(0,a^2K)$ is the same set $\mathcal H_K$, but with scaled norm
\[
    \|\mathbf h\|_{\mathcal H_{a^2K}}^2
    =
    \frac{1}{a^2}\|\mathbf h\|_{\mathcal H_K}^2.
\]
Since $m\in\mathcal H_K$, the Cameron-Martin theorem implies that $\mathsf N(m,a^2K)$ is absolutely continuous with respect to $\mathsf N(0,a^2K)$, and for $\mathbf Z^+\in\mathcal H_K$,
\[
    \log
    \frac{
        \d \mathsf N(m,a^2K)
    }{
        \d \mathsf N(0,a^2K)
    }(\mathbf Z^+)
    =
    \frac{1}{a^2}
    \langle \mathbf Z^+,m\rangle_{\mathcal H_K}
    -
    \frac{1}{2a^2}\|m\|_{\mathcal H_K}^2.
\]
Therefore
\[
\begin{aligned}
    -
    \log
    \frac{
        \d \mathsf N(m,a^2K)
    }{
        \d \mathsf N(0,a^2K)
    }(\mathbf Z^+)
    &=
    \frac{1}{2a^2}\|m\|_{\mathcal H_K}^2
    -
    \frac{1}{a^2}
    \langle \mathbf Z^+,m\rangle_{\mathcal H_K}                                      \\
    &=
    \frac{1}{2a^2}
    \|\mathbf Z^+-m\|_{\mathcal H_K}^2
    -
    \frac{1}{2a^2}
    \|\mathbf Z^+\|_{\mathcal H_K}^2 .
\end{aligned}
\]
Substituting back $a=\alpha_\theta(\mathbf Z)$ and
$m=T_\theta(\mathbf Z)$ yields
\[
    -
    \log
    \frac{
        \d P_\theta(\cdot\mid\mathbf Z)
    }{
        \d\mathsf N(0,\alpha_\theta(\mathbf Z)^2K)
    }(\mathbf Z^+)
    =
    \frac{1}{2\alpha_\theta(\mathbf Z)^2}
    \|\mathbf Z^+-T_\theta(\mathbf Z)\|_{\mathcal H_K}^2-
    \frac{1}{2\alpha_\theta(\mathbf Z)^2}
    \|\mathbf Z^+\|_{\mathcal H_K}^2.
\]
Finally, by Lemma~\ref{lem:trace_class_gaussian_cm_kl}, the random series
\[
    \sum_{i=1}^\infty \sqrt{\lambda_i}\xi_i\mathbf E_i
\]
defines a $\mathcal Z$-valued Gaussian random element with law $\mathsf N(0,K)$. Hence
\[
    \alpha_\theta(\mathbf Z)
    \sum_{i=1}^\infty \sqrt{\lambda_i}\xi_i\mathbf E_i
    \sim
    \mathsf N(0,\alpha_\theta(\mathbf Z)^2K).
\]
Adding the mean $T_\theta(\mathbf Z)$ gives
\[
    \mathbf Z^+
    =
    T_\theta(\mathbf Z)
    +
    \alpha_\theta(\mathbf Z)
    \sum_{i=1}^\infty \sqrt{\lambda_i}\xi_i\mathbf E_i
    \sim
    \mathsf N(T_\theta(\mathbf Z),\alpha_\theta(\mathbf Z)^2K),
\]
which is exactly the transition kernel $P_\theta(\cdot\mid\mathbf Z)$. The spectral-coordinate statement in the remark follows by taking inner
products with $\mathbf E_i$.
\end{proof}

\begin{corollary}[Spectral geometry of the Laplacian-resolvent covariance on the torus]
\label{cor:laplacian_resolvent_spectral_geometry}
Let
\[
    \mathcal Z=L^2(\mathbb T^{d_x};\mathbb R^{d_z}),
\]
and let
\[
    \mathcal Z_{\mathbb C}
    =
    L^2(\mathbb T^{d_x};\mathbb C^{d_z})
\]
be its complexification. For each multi-index $m\in\mathbb Z^{d_x}$, define
\[
    \phi_m(x)=e^{2\pi i m\cdot x},
    \qquad x\in\mathbb T^{d_x}.
\]
Let $\{\mathbf e_j\}_{j=1}^{d_z}$ be the standard coordinate basis of $\mathbb R^{d_z}$, and set
\[
    \mathbf E_{m,j}(x)=\phi_m(x)\,\mathbf e_j.
\]
Then $\{\mathbf E_{m,j}\}_{m\in\mathbb Z^{d_x},1\le j\le d_z}$ is an orthonormal basis of the complexified space
\[
    \mathcal Z_{\mathbb C}
    =
    L^2(\mathbb T^{d_x};\mathbb C^{d_z}).
\]
For a real-valued field $\mathbf{Z}=(Z_1,\cdots,Z_{d_z})\in\mathcal Z$, its Fourier coefficients satisfy
\[
    \widehat Z_j(-m)=\overline{\widehat Z_j(m)},
\]
so the complex Fourier expansion is understood subject to this conjugate symmetry.

\begin{itemize}
\item[(a)]$\{\mathbf E_{m,j}\}_{m\in\mathbb Z^{d_x},\,1\le j\le d_z}$ is an orthonormal eigenbasis of $\mathcal Z_\mathbb{C}$ for $K$, with \jg{(I think your $\beta$ below (and proof) should be $s$; see notation in A.4)}
\[ K \mathbf E_{m,j} = \lambda_m \mathbf E_{m,j}, \qquad \lambda_m = \left(1+4\pi^2\ell^2|m|^2\right)^{-\beta}. \]
Equivalently, for any
\[
    Z(x)=\sum_{m\in\mathbb Z^{d_x}}\widehat Z(m)\phi_m(x),
    \qquad 
    \widehat Z(m)\in\mathbb C^{d_z},
\]
one has
\[
    KZ
    =
    \sum_{m\in\mathbb Z^{d_x}}
    \lambda_m \widehat Z(m)\phi_m.
\]
Thus $K$ is a Fourier multiplier with symbol
\[
    \lambda_m=\left(1+4\pi^2\ell^2|m|^2\right)^{-\beta}.
\]

Moreover, $K$ is positive and self-adjoint. It is trace-class if and only if
\[
    \sum_{m\in\mathbb Z^{d_x}}
    \left(1+4\pi^2\ell^2|m|^2\right)^{-\beta}
    <\infty,
\]
which holds precisely when $\beta>d_x/2$.
In this case,
\[
    \operatorname{Tr}(K)
    =
    d_z
    \sum_{m\in\mathbb Z^{d_x}}
    \left(1+4\pi^2\ell^2|m|^2\right)^{-\beta}.
\]

\item[(b)] The Cameron-Martin space associated with the Gaussian measure $\mathsf N(0,K)$ is
\[
    \mathcal H_K
    =
    \operatorname{Range}(K^{1/2}),
\]
with norm
\[
    \|\mathbf{h}\|_{\mathcal H_K}^2
    =
    \sum_{m\in\mathbb Z^{d_x}}
    \sum_{j=1}^{d_z}
    \lambda_m^{-1}
    |\widehat h_j(m)|^2=
    \sum_{m\in\mathbb Z^{d_x}}
    \sum_{j=1}^{d_z}
    \left(1+4\pi^2\ell^2|m|^2\right)^\beta
    |\widehat h_j(m)|^2.
\]
Consequently,
\[
    \mathcal H_K
    \simeq H^\beta(\mathbb T^{d_x};\mathbb R^{d_z}),
\]
with equivalence of norms.
\end{itemize}
\end{corollary}

\begin{proof}
We first recall the spectral decomposition of the periodic Laplacian. For each 
$m\in\mathbb Z^{d_x}$, 
\[
    \phi_m(x)=e^{2\pi i m\cdot x}
\]
solves the eigenvalue problem
\[
    -\Delta\phi_m=4\pi^2|m|^2\phi_m.
\]
Since $\{\phi_m\}_{m\in\mathbb Z^{d_x}}$ is an orthonormal basis of 
$L^2(\mathbb T^{d_x};\mathbb C)$, the vector-valued functions
\[
    \mathbf E_{m,j}(x)=\phi_m(x)e_j,
    \qquad m\in\mathbb Z^{d_x},\ j=1,\dots,d_z,
\]
form an orthonormal basis of $L^2(\mathbb T^{d_x};\mathbb C^{d_z})$, and hence also give the usual complex Fourier representation of $L^2(\mathbb T^{d_x};\mathbb R^{d_z})$ subject to the standard conjugate-symmetry constraint on real-valued fields. Because the Laplacian acts componentwise, we have
\[
    \Delta \mathbf E_{m,j}=(\Delta\phi_m)\mathbf e_j=-4\pi^2|m|^2\mathbf E_{m,j}.
\]
Therefore,
\[
    (I-\ell^2\Delta)\mathbf E_{m,j}
    =
    \left(1+4\pi^2\ell^2|m|^2\right)\mathbf E_{m,j}.
\]
Applying the functional calculus for the positive self-adjoint operator 
$I-\ell^2\Delta$, we obtain
\[
    K \mathbf E_{m,j}
    =(I-\ell^2\Delta)^{-\beta}\mathbf E_{m,j}
    =\left(1+4\pi^2\ell^2|m|^2\right)^{-\beta}\mathbf E_{m,j}.
\]
Thus $E_{m,j}$ is an eigenfunction of $K$ with eigenvalue
\[
    \lambda_m
    =
    \left(1+4\pi^2\ell^2|m|^2\right)^{-\beta}.
\]

It follows immediately that for any Fourier expansion
\[
    Z(x)
    =
    \sum_{m\in\mathbb Z^{d_x}}\widehat Z(m)\phi_m(x),
    \qquad \widehat Z(m)\in\mathbb C^{d_z},
\]
the action of $K$ is given by
\[
    KZ
    =
    \sum_{m\in\mathbb Z^{d_x}}
    \lambda_m\widehat Z(m)\phi_m(x).
\]
Hence $K$ is a Fourier multiplier with multiplier 
$\lambda_m$.

Since each $\lambda_m$ is real and nonnegative, $K$ is positive and self-adjoint. 
Moreover, $K$ is compact because $\lambda_m\to 0$ as $|m|\to\infty$. Its trace is the sum of its eigenvalues, counted with multiplicity. Since each spatial Fourier mode appears once for each latent coordinate direction $e_j$, we have
\[
    \operatorname{Tr}(K)
    =
    \sum_{j=1}^{d_z}
    \sum_{m\in\mathbb Z^{d_x}}\lambda_m
    =
    d_z
    \sum_{m\in\mathbb Z^{d_x}}
    \left(1+4\pi^2\ell^2|m|^2\right)^{-\beta}.
\]
For large $|m|$,
\[
    \left(1+4\pi^2\ell^2|m|^2\right)^{-\beta}
    \asymp
    |m|^{-2\beta}.
\]
The lattice sum
\[
    \sum_{m\in\mathbb Z^{d_x}} |m|^{-2\beta}
\]
converges if and only if $2\beta>d_x$. Therefore $K$ is trace-class if and only if
\[
    \beta>\frac{d_x}{2}.
\]

It remains to identify the Cameron-Martin space. Since $K$ is diagonal in the basis $\{\mathbf E_{m,j}\}$, its square root satisfies
\[
    K^{1/2}\mathbf E_{m,j}=\lambda_m^{1/2}\mathbf E_{m,j}.
\]
Thus $\mathbf h\in\operatorname{Range}(K^{1/2})$ if and only if
\[
    \mathbf h=K^{1/2}\mathbf g
\]
for some $g\in L^2(\mathbb T^{d_x};\mathbb R^{d_z})$. In Fourier coordinates, this means
\[
    \widehat h_j(m)
    =
    \lambda_m^{1/2}\widehat g_j(m).
\]
Hence
\[
    \sum_{m\in\mathbb Z^{d_x}}\sum_{j=1}^{d_z}
    \lambda_m^{-1}|\widehat h_j(m)|^2
    =
    \sum_{m\in\mathbb Z^{d_x}}\sum_{j=1}^{d_z}
    |\widehat g_j(m)|^2
    <\infty.
\]
The Cameron-Martin norm is therefore
\[
    \|\mathbf h\|_{\mathcal H_K}^2=\sum_{m\in\mathbb Z^{d_x}}\sum_{j=1}^{d_z}
    \lambda_m^{-1}|\widehat h_j(m)|^2.
\]
Substituting the expression for $\lambda_m$ gives
\[
    \|\mathbf h\|_{\mathcal H_K}^2=\sum_{m\in\mathbb Z^{d_x}}\sum_{j=1}^{d_z}
    \left(1+4\pi^2\ell^2|m|^2\right)^\beta
    |\widehat h_j(m)|^2.
\]
This is equivalent to the usual Fourier definition of the Sobolev norm on the torus:
\[
    \|\mathbf h\|_{H^\beta(\mathbb T^{d_x};\mathbb R^{d_z})}^2
    =\sum_{m\in\mathbb Z^{d_x}}\sum_{j=1}^{d_z}
    \left(1+4\pi^2|m|^2\right)^\beta|\widehat h_j(m)|^2.
\]
The two norms are equivalent because $\ell>0$ is fixed. Therefore
\[
    \mathcal H_K
    \simeq
    H^\beta(\mathbb T^{d_x};\mathbb R^{d_z}),
\]
which completes the proof.
\end{proof}


\begin{lemma}[Moment estimates]\label{lem:moment_gaussian}
Let $\mathcal Z$ be a separable Hilbert space, and let $K:\mathcal Z\to\mathcal Z$ be positive, self-adjoint, and trace-class operator. Assume the random element $\mathbf G\sim\mathsf{N}(0,K)$. Then
\begin{itemize}
    \item[(a)]$\bbE\Vert\mathbf G\Vert_\mathcal{Z}^4\leq 3(\operatorname{Tr} K)^2.$
    \item[(b)]$\bbE\Vert\mathbf G\Vert_\mathcal{Z}^8\leq 105(\operatorname{Tr} K)^4.$
\end{itemize}
\end{lemma}
\begin{proof}
Let $\{(\lambda_i,\mathbf E_i)\}_{i\ge 1}$ be an eigensystem of $K$. By Karhunen-Lo\`eve expansion,
\[
    \mathbf G
    =
    \sum_{i=1}^\infty\sqrt{\lambda_i}\xi_i\mathbf E_i,
    \quad
    \xi_i\overset{\mathrm{i.i.d.}}{\sim}\mathsf N(0,1).
\]
Thus $\|\mathbf G\|_{\mathcal Z}^2=\sum_{i=1}^\infty \lambda_i\xi_i^2$, and the cumulant generating function of $\Vert\mathbf G\Vert_\mathcal{Z}^2$ is given by
\begin{align*}
    \log\bbE\left[e^{t\Vert\mathbf G\Vert_\mathcal{Z}^2}\right]&=\sum_{i=1}^\infty\log\bbE\left[e^{t\lambda_i\xi_i^2}\right]=-\frac{1}{2}\sum_{i=1}^\infty\log(1-2t\lambda_i),\\
    &=-\frac{1}{2}\sum_{i=1}^\infty\sum_{n=1}^\infty\frac{(2t\lambda_i)^n}{n}=\sum_{n=1}^\infty\frac{2^{n-1}}{n}\left(\sum_{i=1}^\infty\lambda_i^n\right) t^n,\quad t<\frac{1}{2}.
\end{align*}
Matching the coefficients in
\begin{equation*}
    \log\bbE\left[e^{t\Vert\mathbf G\Vert_\mathcal{Z}^2}\right]=\sum_{n=1}^\infty\frac{\kappa_n}{n!}t^n,
\end{equation*}
we obtain the cumulants of $\Vert\mathbf{G}\Vert_\mathcal{Z}^2$:
\begin{equation}
\kappa_n=2^{n-1}(n-1)!\sum_{i=1}^\infty\lambda_i^n=2^{n-1}(n-1)!\,\mathrm{Tr}(K^n),\quad n=1,2,\cdots.
\end{equation}
For fixed $n\in\mathbb N$, let $\Pi(n)$ denote the set of partitions of $\{1,\cdots,n\}$. We use the moment-cumulant formula for $\Vert\mathbf{G}\Vert_\mathcal{Z}^2$:
\begin{equation*}
    \bbE\left\Vert\mathbf G\right\Vert_\mathcal{Z}^{2n}=\sum_{\pi\in\Pi(n)}\prod_{V\in\pi}\kappa_{\vert V\vert}.
\end{equation*}
Then
\begin{equation*}
    \bbE\left\Vert\mathbf G\right\Vert_\mathcal{Z}^4=\kappa_2+\kappa_1^2=2\operatorname{Tr}(K^2)+(\operatorname{Tr} K)^2.
\end{equation*}
and
\begin{align*}
    \bbE\left\Vert\mathbf G\right\Vert_\mathcal{Z}^8&=\kappa_4+4\kappa_1\kappa_3+3\kappa_2^2+6\kappa_1^2\kappa_2+\kappa_1^4\\
    &=48\operatorname{Tr}(K^4)+32\operatorname{Tr}(K)\operatorname{Tr}(K^3)+12\operatorname{Tr}(K^2)^2+12(\operatorname{Tr} K)^2\operatorname{Tr}(K^2)+\operatorname{Tr}(K)^4.
\end{align*}
Finally, since $K$ is positive trace-class, we have
\begin{equation*}
    \operatorname{Tr}(K^n)=\sum_{i=1}^\infty\lambda_i^n\leq\left(
    \sum_{i=1}^\infty\lambda_i\right)^n=(\operatorname{Tr} K)^n.
\end{equation*}
Combining the last three displays gives the desired estimates.
\end{proof}


\begin{lemma}[Local sensitivity under Gaussian perturbations]
\label{lem:local_jacobian_decoder_robustness}
Let $\mathcal Z$ and $\mathcal U$ be separable Hilbert spaces, and let $K:\mathcal Z\to\mathcal Z$ be positive, self-adjoint, and trace-class operator.
Let $\mathbf G\sim \mathsf N(0,K)$. Fix $\mathbf z\in\mathcal Z$ and $a>0$. Suppose that $F:\mathcal Z\to\mathcal U$ is Fr\'echet differentiable at $\mathbf z$, with Fr\'echet derivative
\[
    J_{F}(\mathbf z)\in\mathfrak{B}(\mathcal Z,\mathcal U),
\]
and assume that, on the relevant neighborhood of $\mathbf z$, the second-order remainder satisfies
\[
    \left\|F(\mathbf z+\mathbf h)-F(\mathbf z)-J_{F}(\mathbf z)\mathbf h\right\|_{\mathcal U}
    \le\frac{M}{2}\|\mathbf h\|_{\mathcal Z}^2 .
\]
Then
\[
    \mathbb E\left[\left\|F(\mathbf z+a\mathbf G)-F(\mathbf z)\right\|_{\mathcal U}^2
    \right]\le a^2\Vert J\Vert_\mathrm{op}^2\operatorname{Tr} K+\frac{3}{4}M^2a^4(\operatorname{Tr} K)^2.
\]
\end{lemma}

\begin{proof}
For brevity, write $J=J_{F}(\mathbf z)$. By the assumed second-order expansion, for every $\mathbf h$ in the relevant neighborhood,
\[
    F(\mathbf z+\mathbf h)-F(\mathbf z)
    =J\mathbf h+R(\mathbf h),
\]
where
\[
    \|R(\mathbf h)\|_{\mathcal U}\le\frac{M}{2}\|\mathbf h\|_{\mathcal Z}^2.
\]
Taking $\mathbf h=a\mathbf G$, we have
\[
    F(\mathbf z+a\mathbf G)-F(\mathbf z) = aJ\mathbf G+R(a\mathbf G).
\]
Therefore,
\[
\begin{aligned}
    \mathbb E\left[\|F(\mathbf z+a\mathbf G)-F(\mathbf z)\|_{\mathcal U}^2\right]=\mathbb E
    \left[\|aJ\mathbf G+R(a\mathbf G)\|_{\mathcal U}^2
    \right]\le 2a^2\mathbb E\|J\mathbf G\|_{\mathcal U}^2
    +2\mathbb E\|R(a\mathbf G)\|_{\mathcal U}^2.
\end{aligned}
\]
We first compute the linear term. Since $\mathbf G\sim \mathsf N(0,K)$ and $J$ is bounded linear, $J\mathbf G$ is a Gaussian random element in $\mathcal U$ with covariance operator $J K J^*$.
Hence
\[
    \mathbb E\|J\mathbf G\|_{\mathcal U}^2
    =\Vert J\Vert_\mathrm{op}^2\mathbb E\|\mathbf G\|_{\mathcal U}^2\leq\Vert J\Vert_\mathrm{op}^2\operatorname{Tr}(K).
\]
We next bound the remainder term. From the second-order remainder assumption,
\[
    \|R(a\mathbf G)\|_{\mathcal U}\le\frac{M}{2}a^2\|\mathbf G\|_{\mathcal Z}^2.
\]
Therefore
\[
\mathbb E\|R(a\mathbf G)\|_{\mathcal U}^2\le\frac{M^2}{4}a^4\mathbb E\|\mathbf G\|_{\mathcal Z}^4.
\]
Combining the remainder estimate and the moment estimate in Lemma \ref{lem:moment_gaussian} yields
\[
    \rho(\mathbf z,a)\le
    a^2\Vert J\Vert_\mathrm{op}^2\operatorname{Tr} K+\frac{3}{4}M^2a^4(\operatorname{Tr} K)^2.
\]
which proves the desired bound.
\end{proof}


\begin{proof}[Proof of Proposition \ref{prop:gaussian_sensitivity_bounds}]
We first prove the Lipschitz robustness estimates. By the Lipschitz continuity of $T_\theta$ on the relevant regions,\[
    \|T_\theta(\mathbf Z_n^\star+\Xi_n)
      -T_\theta(\mathbf Z_n^\star)\|_{\mathcal Z}
    \leq
    \Lambda_T\|\Xi_n\|_{\mathcal Z}.
\]
Conditionally on $\mathbf U_n^\star$, one samples  $\Xi_n\sim\mathsf N\bigl(0,K_\phi(\mathbf U_n^\star)\bigr)$.
Then we have
\[
    \mathbb E
    \left[
        \|\Xi_n\|_{\mathcal Z}^2
        \,\middle|\,
        \mathbf U_n^\star
    \right]
    =
    \operatorname{Tr}
    K_\phi(\mathbf U_n^\star).
\]
Taking expectation over the reference trajectory therefore gives
\[
    \eta_n^2
    \leq\Lambda_T^2
    \mathbb E\left[\operatorname{Tr}
    K_\phi(\mathbf U_n^\star)\right]
    \leq\Lambda_T^2\varsigma.
\]
Similarly,
\[
\|D_\psi(\widetilde{\mathbf Z}_n^+) -D_\psi(T_\theta(\widetilde{\mathbf Z}_n))\|_{\mathcal U}\leq\Lambda_D\alpha_\theta(\widetilde{\mathbf Z}_n)\|\mathbf G_n\|_{\mathcal Z}.
\]
Using $\alpha_\theta\leq\bar\alpha$ and
$\mathbb E\|\mathbf G_n\|_{\mathcal Z}^2=\operatorname{Tr}(K)$ yields
\[
    \rho_n
    \leq
    \Lambda_D\bar\alpha\sqrt{\operatorname{Tr}(K)}.
\]
This proves \eqref{eq:lipschitz_sensitivity_bounds}.

Now we derive the local estimate for $\rho_n$. By Lemma \ref{lem:local_jacobian_decoder_robustness},
\begin{align*}
    &\bbE_{\mathbf{G}_n\sim\mathsf N(0,K)}\left\Vert D_\psi\bigl(T_\theta(\wt{\mathbf Z}_n)+\alpha_\theta(\wt{\mathbf Z}_n)\mathbf G_n\bigr)-D_\psi(T_\theta(\wt{\mathbf Z}_n))\right\Vert_\mathcal{U}^2\\
    &\quad\leq\alpha_\theta(\wt{\mathbf Z}_n)^2\bigl\Vert J_{D_\psi}(T_\theta(\wt{\mathbf Z}_n))\bigr\Vert_\mathrm{op}^2\operatorname{Tr} K +\frac{3}{4}M_D^2\alpha_\theta(\wt{\mathbf Z}_n)^4(\operatorname{Tr} K)^2.
\end{align*}
Taking expectation over the reference trajectory and Gaussian perturbation $\Xi_n$, and using the fact $\alpha_\theta(\wt{\mathbf Z}_n)\leq\ol\alpha$, we have
\begin{equation*}
    \rho_n\leq\ol\alpha\left(\bbE\left[\Vert J_{D_\psi}(T_\theta(\wt{\mathbf Z}_n))\Vert_\mathrm{op}^2\right]\right)^{1/2}\sqrt{\operatorname{Tr} K}+\frac{\sqrt{3}M_D\ol\alpha^2}{2}\operatorname{Tr} K,
\end{equation*}
which is the desired bound. Similarly, \jg{applying Lemma \ref{lem:local_jacobian_decoder_robustness} conditionally on $\mathbf U_n^\star$, with $a = 1$, and $K=K_\phi(\mathbf U_n^\star)$, and using $\operatorname{Tr} K_\phi(\mathbf U_n^\star)\leq\varsigma$ pointwise before taking any expectation,
\[
    \mathbb E\left[\left\Vert T_\theta(\wt{\mathbf Z}_n)-T_\theta(\mathbf Z_n^\star)\right\Vert_{\mathcal Z}^2\,\middle|\,\mathbf U_n^\star\right]
    \leq
    \left\Vert J_{T_\theta}(\mathbf Z_n^\star)\right\Vert_\mathrm{op}^2\varsigma+\frac{3}{4}M_T^2\varsigma^2.
\]
Taking expectation over the reference trajectory and applying $\sqrt{x+y}\leq\sqrt{x}+\sqrt{y}$ for any $x,y\geq0$ gives
\begin{equation*}
    \eta_n\leq\left(\bbE\left[\Vert J_{T_\theta}(\mathbf Z^\star_n)\Vert_\mathrm{op}^2\right]\varsigma\right)^{1/2}+\frac{\sqrt{3}M_T}{2}\varsigma,
\end{equation*}
which is \eqref{eq:tau_jacobian_bound}.
}
% \begin{equation*}
%     \eta_n\leq\left(\bbE\left[\Vert J_{T_\theta}(\mathbf Z^\star_n)\Vert_\mathrm{op}^2\operatorname{Tr} K_\phi(\mathbf{U}_n^\star)\right]\right)^{1/2}+\frac{\sqrt{3}M_T}{2}\mathbb E\left[\operatorname{Tr} K_\phi(\mathbf{U}_n^\star)\right],
% \end{equation*}
% and \eqref{eq:tau_jacobian_bound} follows from $\operatorname{Tr} K_\phi(\mathbf{U}_n^\star)\leq\varsigma$.
\end{proof}

\begin{proof}[Proof of Proposition \ref{prop:generic_rollout}]
Define the population rollout error
\[
e_n
:=
\left(
\mathbb E
\left[
\|\widehat{\mathbf U}_n-\mathbf U_n^\star\|_{\mathcal U}^2
\right]
\right)^{1/2}.
\]
Since $\widehat{\mathbf U}_0=\mathbf U_0^\star$, we have $e_0=0$. For every
$n\geq0$,
\begin{align*}
\|\widehat{\mathbf U}_{n+1}-\mathbf U_{n+1}^\star\|_{\mathcal U}
&=
\|F_\vartheta(\widehat{\mathbf U}_n)-\mathbf U_{n+1}^\star\|_{\mathcal U}
\\
&\leq
\|F_\vartheta(\widehat{\mathbf U}_n)
-F_\vartheta(\mathbf U_n^\star)\|_{\mathcal U}
+
\|F_\vartheta(\mathbf U_n^\star)-\mathbf U_{n+1}^\star\|_{\mathcal U}.
\end{align*}
Taking the population $L^2$ norm and applying Minkowski's inequality together
with the Lipschitz continuity of $F_\vartheta$ gives
\[
e_{n+1}
\leq
\Lambda_F e_n+\epsilon_n^{\mathrm{dir}}.
\]
Iterating this recursion from $e_0=0$ yields
\[
e_n
\leq
\sum_{j=0}^{n-1}
\Lambda_F^{\,n-1-j}\epsilon_j^{\mathrm{dir}},
\]
which proves \eqref{eq:generic_rollout}.
\end{proof}


\iffalse
\begin{proof}[Proof of Proposition \ref{prop:encoder-noisy-rollout-bound}]
(a) We first prove the one-step estimate. Fix $n\in\{0,\cdots,N-1\}$ and write
\[
\mathbf A_n := D_\psi(T_\theta(\widetilde{\mathbf Z}_n)),
\qquad
\mathbf B_n := D_\psi(\widetilde{\mathbf Y}_n).
\]
For every realization of $(\Xi_n,G_n)$, the triangle inequality gives
\[
\|D_\psi\left(T_\theta(\mathbf Z_n^\star)\right)-\mathbf U^\star_{n+1}\|_{\mathcal U}
\le
\|D_\psi\left(T_\theta(\mathbf Z_n^\star)\right)-\mathbf A_n\|_{\mathcal U}
+
\|\mathbf A_n-\mathbf B_n\|_{\mathcal U}
+
\|\mathbf B_n-\mathbf U^\star_{n+1}\|_{\mathcal U}.
\]
The left-hand side is deterministic. Taking the $L^2$ norm with respect to random elements $(\Xi_n,\mathbf G_n)$ and using Minkowski's inequality yields
\begin{equation*}
\begin{aligned}
\|D_\psi\left(T_\theta(\mathbf Z_n^\star)\right)-\mathbf U^\star_{n+1}\|_{\mathcal U}
&\le\left(\mathbb E_{\Xi_n}\|D_\psi\left(T_\theta(\mathbf Z_n^\star)\right)-\mathbf A_n\|_{\mathcal U}^2\right)^{\frac12}+\left(\mathbb E_{\Xi_n,\mathbf G_n}\|\mathbf A_n-\mathbf B_n\|_{\mathcal U}^2\right)^{\frac12}\\
&\quad +\left(\mathbb E_{\Xi_n,\mathbf G_n}\|\mathbf B_n-\mathbf U^\star_{n+1}\|_{\mathcal U}^2\right)^{\frac12}.
\end{aligned}
\end{equation*}
By the definitions of $\eta_n$, $\rho_n$, and $\epsilon_n$, this gives
\[
\left\|D_\psi(T_\theta(\mathbf Z_n^\star))-\mathbf U^\star_{n+1}\right\|_{\mathcal U}\le\eta_n+\rho_n+\epsilon_n .
\]
This proves the first claim.

\item (b) We now prove the Lipschitz robustness estimates. By the Lipschitz continuity of $D_\psi$ and $T_\theta$ on the relevant regions,
\[
\begin{aligned}
\eta_n &=\left(\mathbb E_{\Xi_n}
\left[\left\|D_\psi(T_\theta(\mathbf Z_n^\star+\Xi_n))-D_\psi(T_\theta(\mathbf Z_n^\star))
\right\|_{\mathcal U}^2\right]\right)^{1/2} \\
&\le\Lambda_D\left(\mathbb E_{\Xi_n}\left[\left\|
T_\theta(\mathbf Z_n^\star+\Xi_n)-T_\theta(\mathbf Z_n^\star)
\right\|_{\mathcal Z}^2\right]\right)^{1/2}\le\Lambda_D\Lambda_T
\left(\mathbb E_{\Xi_n}\|\Xi_n\|_{\mathcal Z}^2\right)^{1/2}.
\end{aligned}
\]
Since $\Xi_n\sim\mathcal N(0,K_\phi(\mathbf U_n^\star))$ and $K_\phi(\mathbf U_n^\star)$ is trace-class,
\[
\mathbb E_{\Xi_n}\|\Xi_n\|_{\mathcal Z}^2=\operatorname{Tr}(K_\phi(\mathbf U_n^\star)).
\]
Hence
\[
\eta_n\le\Lambda_D\Lambda_T\sqrt{\operatorname{Tr}(K_\phi(\mathbf U_n^\star))}.
\]
Similarly,
\[
\begin{aligned}
\rho_n&=\left(\mathbb E_{\Xi_n,\mathbf G_n}
\left[\left\|D_\psi\left(T_\theta(\widetilde{\mathbf Z}_n)+\alpha_\theta(\widetilde{\mathbf Z}_n)\mathbf G_n\right)-D_\psi(T_\theta(\widetilde{\mathbf Z}_n))
\right\|_{\mathcal U}^2\right]\right)^{1/2} \\
&\le\Lambda_D\left(\mathbb E_{\Xi_n,\mathbf G_n}\left[\alpha_\theta(\widetilde{\mathbf Z}_n)^2\|\mathbf G_n\|_{\mathcal Z}^2\right]\right)^{1/2}.
\end{aligned}
\]
Because $\Xi_n$ and $G_n$ are independent,
\[
\mathbb E_{\Xi_n,\mathbf G_n}
\left[\alpha_\theta(\widetilde{\mathbf Z}_n)^2
\|\mathbf G_n\|_{\mathcal Z}^2\right]=\mathbb E_{\Xi_n}\left[\alpha_\theta(\mathbf Z_n^\star+\Xi_n)^2\right]\mathbb E_{\mathbf G_n}\left[\|\mathbf G_n\|_{\mathcal Z}^2\right].
\]
Since $\mathbf G_n\sim\mathcal N(0,K)$,
\[
\mathbb E_{\mathbf G_n}\|\mathbf G_n\|_{\mathcal Z}^2=\operatorname{Tr}(K).
\]
Therefore
\[
\rho_n\le\Lambda_D\bar\alpha_n\sqrt{\operatorname{Tr}(K)}.
\]
The bound with $\alpha_{\max}$ follows immediately if
$\alpha_\theta(z)\le\alpha_{\max}$ on the relevant noisy tube.

\item (c) We first prove the estimate for $\rho_n$. By Lemma \ref{lem:local_jacobian_decoder_robustness},
\begin{align*}
    &\left(\bbE_{\mathbf{G}_n\sim\mathsf N(0,K)}\left\Vert D_\psi\bigl(T_\theta(\wt{\mathbf Z}_n)+\alpha_\theta(\wt{\mathbf Z}_n)\mathbf G_n\bigr)-D_\psi(T_\theta(\wt{\mathbf Z}_n))\right\Vert_\mathcal{U}^2\right)^{1/2}\\
    &\quad\leq\alpha_\theta(\wt{\mathbf Z}_n)\Vert JD_\psi(T_\theta(\wt{\mathbf Z}_n))\Vert_\mathrm{op}\sqrt{\operatorname{Tr} K}+\frac{\sqrt{3}M_D}{2}\alpha_\theta(\wt{\mathbf Z}_n)^2\operatorname{Tr} K.
\end{align*}
Taking expectation with respect to $\Xi_n\sim\mathsf N(0,K_\phi)$ and using the fact $\alpha_\theta(\wt{\mathbf Z}_n)\leq\ol\alpha$, we have
\begin{equation*}
    \rho_n\leq\ol\alpha\left(\bbE_{\Xi_n}\left[\Vert JD_\psi(T_\theta(\wt{\mathbf Z}_n))\Vert_\mathrm{op}^2\right]\right)^{1/2}\sqrt{\operatorname{Tr} K}+\frac{\sqrt{3}M_D\ol\alpha^2}{2}\operatorname{Tr} K,
\end{equation*}
which is the first desired bound. Now we prove the estimate for $\eta_n$. Fix $n$ and write
\begin{equation}
    \mathbf{z}:=\mathbf Z_n^\star,\qquad\mathbf{y}:=T_\theta(\mathbf{z}),\qquad J_T:=JT_\theta(\mathbf{z}),\qquad J_D:=JD_\psi(\mathbf{y}).
\end{equation}
For a perturbation $\mathbf{h}$, define the transition remainder
\begin{equation*}
    R_T(\mathbf{h})=T_\theta(\mathbf{z}+\mathbf{h})-T_\theta(\mathbf{z})-J_T\mathbf{h}.
\end{equation*}
By assumption,
\begin{equation*}
    \Vert R_T(\mathbf{h})\Vert_\mathcal{Z}\leq\frac{M_T}{2}\Vert\mathbf{h}\Vert_\mathcal{Z}^2.
\end{equation*}
We also define
\begin{equation}
    q(\mathbf{h}):=T_\theta(\mathbf{z}+\mathbf h)-T_\theta(\mathbf z)=J_T\mathbf{h}+R_T(\mathbf{h}).\label{eq:qremainderdef}
\end{equation}
Then decoder expansion at $\mathbf{y}=T_\theta(\mathbf{z})$ writes
\begin{equation*}
    D_\psi(\mathbf{y}+q(\mathbf h))-D_\psi(\mathbf y)=J_Dq(\mathbf{h})+R_D(q(\mathbf h)),
\end{equation*}
where
\begin{equation*}
    \Vert R_D(q(\mathbf h))\Vert_\mathcal{U}\leq\frac{M_D}{2}\Vert q(\mathbf h)\Vert_\mathcal{Z}^2.
\end{equation*}
Therefore
\begin{equation*}
\begin{aligned}
    H_{\theta,\psi}(\mathbf{z}+\mathbf{h})-H_{\theta,\psi}(\mathbf z)&=D_\psi\left(T_\theta(\mathbf{z}+\mathbf h))-D_\psi(T_\theta(\mathbf z))\right)\\
    &=D_\psi\left(T_\theta(\mathbf z)+q(\mathbf{h}))-D_\psi(T_\theta(\mathbf z))\right)\\
    &=J_Dq(\mathbf h)+R_D(q(\mathbf h))\\
    &=J_DJ_T\mathbf{h}+J_DR_T(\mathbf{h})+R_D(q(\mathbf h))
\end{aligned}
\end{equation*}
By Minkowski's inequality,
\begin{align*}
    \eta_n&=\left(\bbE_{\Xi\sim\mathsf N(0,K_\phi)}\left[\left\Vert H_{\theta,\psi}(\mathbf z+\Xi)-H_{\theta,\psi}(\mathbf z)\right\Vert_\mathcal{U}^2\right]\right)^{1/2}\\
    &\leq \left(\bbE_{\Xi\sim\mathsf N(0,K_\phi)}\Vert J_DJ_T\Xi\Vert_\mathcal{U}^2\right)^{1/2}+\left(\bbE_{\Xi\sim\mathsf N(0,K_\phi)}\Vert J_DR_T(\Xi)\Vert_\mathcal{U}^2\right)^{1/2}\\
    &\qquad+\left(\bbE_{\Xi\sim\mathsf N(0,K_\phi)}\Vert R_D(q(\Xi))\Vert_\mathcal{U}^2\right)^{1/2}\\
    &\leq\left\Vert J_D\right\Vert_\mathrm{op}\left\Vert J_T\right\Vert_\mathrm{op}\left(\bbE_{\Xi\sim\mathsf N(0,K_\phi)}\Vert\Xi\Vert_\mathcal{Z}^2\right)^{1/2}+\frac{M_T}{2}\left\Vert J_D\right\Vert_\mathrm{op}\left(\bbE_{\Xi\sim\mathsf N(0,K_\phi)}\Vert\Xi\Vert_\mathcal{Z}^4\right)^{1/2}\\
    &\qquad+\frac{M_D}{2}\left(\bbE_{\Xi\sim\mathsf N(0,K_\phi)}\Vert q(\Xi)\Vert_\mathcal{Z}^4\right)^{1/2}.
\end{align*}
By \eqref{eq:qremainderdef}, Jensen's inequality, and the moment estimates in Lemma \ref{lem:moment_gaussian},
\begin{align*}
    \bbE_{\Xi\sim\mathsf N(0,K_\phi)}\Vert q(\Xi)\Vert_\mathcal{Z}^4&\leq 8\left(\bbE_{\Xi\sim\mathsf N(0,K_\phi)}\Vert J_T\Xi\Vert_\mathcal{Z}^4+\bbE_{\Xi\sim\mathsf N(0,K_\phi)}\Vert R_T(\Xi)\Vert_\mathcal{Z}^4\right)\\
    &\leq 8\Vert J_T\Vert_\mathrm{op}^4\bbE_{\Xi\sim\mathsf N(0,K_\phi)}\Vert \Xi\Vert_\mathcal{Z}^4+\frac{M_T^4}{2}\bbE_{\Xi\sim\mathsf N(0,K_\phi)}\Vert \Xi\Vert_\mathcal{Z}^8\\
    &\leq 24\Vert J_T\Vert_\mathrm{op}^4(\operatorname{Tr} K_\phi)^2+\frac{105 M_T^4}{2}(\operatorname{Tr} K_\phi)^4.
\end{align*}
Hence
\begin{equation*}
    \eta_n\leq \left\Vert J_D\right\Vert_\mathrm{op}\left\Vert J_T\right\Vert_\mathrm{op}\sqrt{\operatorname{Tr}K_\phi}+\left(\frac{\sqrt{3}M_T}{2}\left\Vert J_D\right\Vert_\mathrm{op}+\sqrt{6}M_D\Vert J_T\Vert_\mathrm{op}^2\right)\operatorname{Tr}K_\phi+\sqrt{\frac{105}{8}}M_DM_T^2(\operatorname{Tr} K_\phi)^2.
\end{equation*}
Since $\operatorname{Tr} K_\phi\leq \varsigma$, the desired bound follows.

\item (d) It remains to prove the rollout estimate. By the triangle inequality and the reconstruction defect,
\[
\begin{aligned}
\left\|
D_\psi(T_\theta(\mathbf Z_n^\star))-D_\psi(\mathbf Z_{n+1}^\star)
\right\|_{\mathcal U}
&\le
\left\|
D_\psi(T_\theta(\mathbf Z_n^\star))-\mathbf  U^\star_{n+1}
\right\|_{\mathcal U}
+
\left\|
\mathbf U^\star_{n+1}-D_\psi(\mathbf Z_{n+1}^\star)
\right\|_{\mathcal U} \\
&\le\epsilon_n+\rho_n+\eta_n+\delta_{n+1}.
\end{aligned}
\] 
From local inverse stability,
\[
\|T_\theta(\mathbf Z_n^\star)-\mathbf Z_{n+1}^\star\|_{\mathcal Z}\le\Lambda_{\mathrm{inv}}\left\|D_\psi(T_\theta(\mathbf Z_n^\star))-D_\psi(\mathbf Z_{n+1}^\star)\right\|_{\mathcal U}\le\Lambda_{\mathrm{inv}}\left(\epsilon_n+\rho_n+\eta_n+\delta_{n+1}\right).
\]
For fixed $n$, define the latent rollout error
\[
e_n := \|\widehat{\mathbf Z}_n-\mathbf Z_n^\star\|_{\mathcal Z},\quad n=1,\cdots,N.
\]
Using $\widehat{\mathbf Z}_{n+1}=T_\theta(\widehat{\mathbf Z}_n)$, the triangle inequality, and the local Lipschitz bound for $T_\theta$, we obtain
\[
\begin{aligned}
e_{n+1}
&=
\|\widehat{\mathbf Z}_{n+1}-\mathbf Z_{n+1}^\star\|_{\mathcal Z} =
\|T_\theta(\widehat{\mathbf Z}_n)-\mathbf Z_{n+1}^\star\|_{\mathcal Z} \\
&\le
\|T_\theta(\widehat{\mathbf Z}_n)-T_\theta(\mathbf Z_n^\star)\|_{\mathcal Z}
+
\|T_\theta(\mathbf Z_n^\star)-\mathbf Z_{n+1}^\star\|_{\mathcal Z} \\
&\le
\Lambda_T e_n
+
\Lambda_{\mathrm{inv}}
\left(
\epsilon_n+\rho_n+\eta_n+\delta_{n+1}
\right).
\end{aligned}
\]
Beginning from $e_0=0$ and iterating this scalar recursion yield
\[
e_n\le\Lambda_{\mathrm{inv}}
\sum_{j=0}^{n-1}\Lambda_T^{\,n-1-j}
\left(\epsilon_j+\rho_j+\eta_j+\delta_{j+1}\right),\quad n=1,\cdots,N.
\]
Finally,
\[
\begin{aligned}
\|\widehat{\mathbf U}_n-\mathbf U^\star_n\|_{\mathcal U}
&\le\|\widehat{\mathbf U}_n-\bar{\mathbf U}_n\|_{\mathcal U}+\|\bar{\mathbf U}_n^\star-\mathbf U^\star_n\|_{\mathcal U} \\
&=\|D_\psi(\widehat{\mathbf Z}_n)-D_\psi(\mathbf Z_n^\star)\|_{\mathcal U}+\delta_n \le\Lambda_D e_n+\delta_n.
\end{aligned}
\]
Substituting the estimate for $e_n$ gives
\[
\|\widehat{\mathbf U}_n-\mathbf U^\star_n\|_{\mathcal U}\le\Lambda_D\Lambda_{\mathrm{inv}}\sum_{j=0}^{n-1}\Lambda_T^{\,n-1-j}\left(\epsilon_j+\rho_j+\eta_j+\delta_{j+1}
\right)+\delta_n.
\]
This completes the proof.
\end{proof}

\begin{proof}[Derivation of bound \eqref{eq:deterministic_rollout_bound}]
Given the same reference physical trajectory
\[
    \{\mathbf U_n^\star\}_{n=0}^N\subset\mathcal U,
\]
define the deterministic reference latent trajectory by
\[
    \mathbf Z_n^{\star,\mathrm{det}}
    :=
    m_\phi^{\mathrm{det}}(\mathbf U_n^\star),
\]
and the corresponding decoded reference trajectory by
\[
    \bar{\mathbf U}_n^{\mathrm{det}}
    :=
    D_\psi^{\mathrm{det}}(\mathbf Z_n^{\star,\mathrm{det}}).
\]
As in the variational case, exact reconstruction is not assumed. We therefore define
\[
    \delta_n^{\mathrm{det}}
    :=
    \left\|
    \bar{\mathbf U}_n^{\mathrm{det}}
    -
    \mathbf U_n^\star
    \right\|_{\mathcal U}.
\]
The deterministic one-step training error is the pointwise decoded transition error
\[
    \epsilon_n^{\mathrm{det}}
    :=
    \left\|
    D_\psi^{\mathrm{det}}
    \left(
    T_\theta^{\mathrm{det}}(\mathbf Z_n^{\star,\mathrm{det}})
    \right)
    -
    \mathbf U_{n+1}^\star
    \right\|_{\mathcal U}.
\]
Consequently,
\[
\begin{aligned}
    &
    \left\|
    D_\psi^{\mathrm{det}}
    \left(
    T_\theta^{\mathrm{det}}(\mathbf Z_n^{\star,\mathrm{det}})
    \right)
    -
    \bar{\mathbf U}_{n+1}^{\mathrm{det}}
    \right\|_{\mathcal U}
    \\
    &\qquad\le
    \left\|
    D_\psi^{\mathrm{det}}
    \left(
    T_\theta^{\mathrm{det}}(\mathbf Z_n^{\star,\mathrm{det}})
    \right)
    -
    \mathbf U_{n+1}^\star
    \right\|_{\mathcal U}
    +
    \left\|
    \mathbf U_{n+1}^\star
    -
    \bar{\mathbf U}_{n+1}^{\mathrm{det}}
    \right\|_{\mathcal U}=
    \epsilon_n^{\mathrm{det}}+\delta_{n+1}^{\mathrm{det}} .
\end{aligned}
\]
At inference time, the deterministic model is rolled out by the mean recursion
\[
    \widehat{\mathbf Z}_0^{\mathrm{det}}
    =
    \mathbf Z_0^{\star,\mathrm{det}}
    =
    m_\phi^{\mathrm{det}}(\mathbf U_0^\star),
\]
\[
    \widehat{\mathbf Z}_{n+1}^{\mathrm{det}}
    =
    T_\theta^{\mathrm{det}}
    \left(
    \widehat{\mathbf Z}_{n}^{\mathrm{det}}
    \right),
    \qquad
    \widehat{\mathbf U}_{n}^{\mathrm{det}}
    =
    D_\psi^{\mathrm{det}}
    \left(
    \widehat{\mathbf Z}_{n}^{\mathrm{det}}
    \right).
\]
Let
\[
    e_n^{\mathrm{det}}
    :=
    \left\|
    \widehat{\mathbf Z}_n^{\mathrm{det}}
    -
    \mathbf Z_n^{\star,\mathrm{det}}
    \right\|_{\mathcal Z}.
\]
Since the rollout is initialized by the deterministic encoder mean, $e_0^{\mathrm{det}}=0$. The
one-step estimate above and local inverse stability imply
\[
    \left\|
    T_\theta^{\mathrm{det}}(\mathbf Z_n^{\star,\mathrm{det}})
    -
    \mathbf Z_{n+1}^{\star,\mathrm{det}}
    \right\|_{\mathcal Z}
    \le
    \Lambda_{\mathrm{inv}}^{\mathrm{det}}
    \left(
    \epsilon_n^{\mathrm{det}}
    +
    \delta_{n+1}^{\mathrm{det}}
    \right).
\]
Therefore,
\[
\begin{aligned}
    e_{n+1}^{\mathrm{det}}
    &=
    \left\|
    T_\theta^{\mathrm{det}}(\widehat{\mathbf Z}_n^{\mathrm{det}})
    -
    \mathbf Z_{n+1}^{\star,\mathrm{det}}
    \right\|_{\mathcal Z}
    \\
    &\le
    \left\|
    T_\theta^{\mathrm{det}}(\widehat{\mathbf Z}_n^{\mathrm{det}})
    -
    T_\theta^{\mathrm{det}}(\mathbf Z_n^{\star,\mathrm{det}})
    \right\|_{\mathcal Z}
    +
    \left\|
    T_\theta^{\mathrm{det}}(\mathbf Z_n^{\star,\mathrm{det}})
    -
    \mathbf Z_{n+1}^{\star,\mathrm{det}}
    \right\|_{\mathcal Z}
    \\
    &\le
    \Lambda_T^{\mathrm{det}}e_n^{\mathrm{det}}
    +
    \Lambda_{\mathrm{inv}}^{\mathrm{det}}
    \left(
    \epsilon_n^{\mathrm{det}}
    +
    \delta_{n+1}^{\mathrm{det}}
    \right).
\end{aligned}
\]
Iterating this recursion gives
\[
    e_n^{\mathrm{det}}
    \le
    \Lambda_{\mathrm{inv}}^{\mathrm{det}}
    \sum_{j=0}^{n-1}
    \left(
    \Lambda_T^{\mathrm{det}}
    \right)^{n-1-j}
    \left(
    \epsilon_j^{\mathrm{det}}
    +
    \delta_{j+1}^{\mathrm{det}}
    \right).
\]
Finally,
\[
    \left\|
    \widehat{\mathbf U}_n^{\mathrm{det}}
    -
    \bar{\mathbf U}_n^{\mathrm{det}}
    \right\|_{\mathcal U}=
    \left\|
    D_\psi^{\mathrm{det}}(\widehat{\mathbf Z}_n^{\mathrm{det}})
    -
    D_\psi^{\mathrm{det}}(\mathbf Z_n^{\star,\mathrm{det}})
    \right\|_{\mathcal U}\le
    \Lambda_D^{\mathrm{det}}e_n^{\mathrm{det}},
\]
which is a decoded-reference bound. Adding the reconstruction defect $\delta_n^{\mathrm{det}}$ gives the physical-reference bound \eqref{eq:deterministic_rollout_bound}.
\end{proof}
\fi 

\subsection{Stochastic Rollouts}
We first prove a tail bound for Gaussian random variables in Hilbert spaces. The following result is similar to the conclusion of \cite{hsu2012tail}.
\begin{lemma}[Gaussian norm tail in a Hilbert space]
\label{lem:hilbert-gaussian-tail}
Let $\mathcal Z$ be a separable Hilbert space and let
\[
    X\sim \mathsf N(0,C),
\]
where $C:\mathcal Z\to\mathcal Z$ is positive, self-adjoint, and trace-class. Then, for every
$t\ge 0$,
\[
    \mathbb P
    \left(
    \|X\|_{\mathcal Z}^2
    \ge
    \operatorname{Tr}(C)
    +
    2\sqrt{\operatorname{Tr}(C^2)t}
    +
    2\|C\|_{\mathrm{op}}t
    \right)
    \le
    e^{-t}.
\]
Equivalently, for every $q\in(0,1)$, with probability at least $1-q$,
\[
    \|X\|_{\mathcal Z}
    \le r_C(q),
\]
where
\[
    r_C(q)
    :=
    \left[
    \operatorname{Tr}(C)
    +
    2\sqrt{\operatorname{Tr}(C^2)\log(1/q)}
    +
    2\|C\|_{\mathrm{op}}\log(1/q)
    \right]^{1/2}.
\]
\end{lemma}
\begin{remark}
Using the relation $\operatorname{Tr}(C^2)\leq(\operatorname{Tr}(C))^2$ and $\Vert C\Vert_\mathrm{op}\leq\operatorname{Tr}(C)$, we have
\begin{equation*}
    \Vert X\Vert_\mathcal{Z}^2\leq\sqrt{\operatorname{Tr} C}\left(1+2\sqrt{\log(1/q)}+2\log(1/q)\right)^{1/2}
\end{equation*}
\end{remark}

\begin{proof}
Since $C$ is positive, self-adjoint, and trace-class, there exists an orthonormal basis
$\{e_i\}_{i\ge 1}$ of $\mathcal Z$ and eigenvalues $\lambda_i\ge 0$ such that
\[
    Ce_i=\lambda_i e_i,
    \qquad
    \sum_{i=1}^\infty \lambda_i=\operatorname{Tr}(C)<\infty .
\]
By the Karhunen-Lo\`eve expansion,
\[
    X=\sum_{i=1}^\infty \sqrt{\lambda_i}\xi_i e_i,
    \qquad
    \xi_i\overset{\mathrm{i.i.d.}}{\sim}\mathsf N(0,1).
\]
Therefore
\[
    \|X\|_{\mathcal Z}^2
    =
    \sum_{i=1}^\infty \lambda_i\xi_i^2
    \quad\text{a.s.}
\]
Set
\[
    S:=\|X\|_{\mathcal Z}^2,
    \qquad
    \mu:=\operatorname{Tr}(C)=\sum_i\lambda_i,
    \qquad
    v:=\operatorname{Tr}(C^2)=\sum_i\lambda_i^2,
    \qquad
    L:=\|C\|_{\mathrm{op}}=\sup_i\lambda_i .
\]
We first estimate the moment generating function of $S-\mu$. For $0\le s<1/(2L)$,
\[
\begin{aligned}
    \log \mathbb E e^{s(S-\mu)}
    &=
    \sum_{i=1}^\infty
    \left[
    -s\lambda_i
    -
    \frac12\log(1-2s\lambda_i)
    \right].
\end{aligned}
\]
Using the expansion
\[
    -\frac12\log(1-2x)-x
    =
    \sum_{k=2}^\infty \frac{2^{k-1}}{k}x^k
    \le
    \frac{x^2}{1-2x},
    \qquad 0\le x<1/2,
\]
with $x=s\lambda_i$, we obtain
\[
\begin{aligned}
    \log \mathbb E e^{s(S-\mu)}
    &\le
    \sum_{i=1}^\infty
    \frac{s^2\lambda_i^2}{1-2s\lambda_i}\le
    \frac{s^2}{1-2sL}
    \sum_{i=1}^\infty \lambda_i^2
    =
    \frac{s^2v}{1-2sL}.
\end{aligned}
\]
Thus $S-\mu$ is sub-gamma with variance factor $2v$ and scale parameter $2L$. By Chernoff's inequality, for any $a>0$ and any $s\in[0,1/(2L))$,
\[
    \mathbb P(S-\mu\ge a)
    \le
    \exp
    \left(
    -sa+\frac{s^2v}{1-2sL}
    \right).
\]
Choose
\[
    a=2\sqrt{vt}+2Lt.
\]
The standard sub-gamma optimization gives
\[
    \mathbb P
    \left(
    S-\mu\ge 2\sqrt{vt}+2Lt
    \right)
    \le
    e^{-t}.
\]
Substituting the definitions of $S,\mu,v,L$ yields
\[
    \mathbb P
    \left(
    \|X\|_{\mathcal Z}^2
    \ge
    \operatorname{Tr}(C)
    +
    2\sqrt{\operatorname{Tr}(C^2)t}
    +
    2\|C\|_{\mathrm{op}}t
    \right)
    \le
    e^{-t}.
\]
Taking $t=\log(1/q)$ gives the stated high-probability form.
\end{proof}

\begin{proposition}[High-probability autoregressive rollout bound]
\label{prop:high-prob-autoregressive-rollout}
Let the assumptions in Proposition \ref{prop:population_kl_rollout} %\jg{(repair cross-ref)} 
hold. Consider the probabilistic autoregressive rollout
\[
    \widetilde{\mathbf Z}_0
    =
    \mathbf Z_0^\star+\Xi_0,
    \qquad
    \Xi_0\sim\mathsf N(0,K_\phi(\mathbf U_0^\star)),
\]
and, for $n=0,\ldots,N-1$,
\[
    \widetilde{\mathbf Z}_{n+1}
    =
    T_\theta(\widetilde{\mathbf Z}_n)
    +
    \alpha_\theta(\widetilde{\mathbf Z}_n)\mathbf G_n,
    \qquad
    \mathbf G_n\sim\mathsf N(0,K),
\]
with
\[
    \widetilde{\mathbf U}_n:=D_\psi(\widetilde{\mathbf Z}_n).
\]
Let $b_n=\Lambda_{\mathrm{inv}}(\epsilon_j+\rho_j+\eta_j+\delta_{j+1})$ be any deterministic latent one-step defect in Proposition \ref{prop:population_kl_rollout} %\jg{(repair cross-ref)}, 
i.e.
\[
    \|T_\theta(\mathbf Z_n^\star)-\mathbf Z_{n+1}^\star\|_{\mathcal Z}
    \le b_n .
\]
Then, for any $\gamma\in(0,1)$, with probability at least $1-\gamma$, the following bound holds
simultaneously for all $n=1,\ldots,N$:
\[
\begin{aligned}
    \|\widetilde{\mathbf U}_n-\mathbf U_n^\star\|_{\mathcal U}
    &\le
    \Lambda_D
    \Bigg[
    \Lambda_T^n\,\varsigma^{1/2}C_{N,\gamma}
    +
    \sum_{j=0}^{n-1}
    \Lambda_T^{n-1-j}
    \sqrt{\operatorname{Tr}(K)}\left(
        b_j
        +
        \bar\alpha\,C_{N,\gamma}
    \right)
    \Bigg]
    +
    \delta_n,
\end{aligned}
\]
where $$C_{N,\gamma}=\left(1+2\sqrt{\log\frac{N+1}{\gamma}}+2\log\frac{N+1}{\gamma}\right)^{1/2}.$$
\end{proposition}
\begin{remark}\label{rem:probabilistic-rollout-bound}
If the rollout is initialized deterministically by the encoder mean, i.e.
$\widetilde{\mathbf Z}_0=\mathbf Z_0^\star$, then the initial encoder-noise term is omitted, and we have
\[
\begin{aligned}
    \|\widetilde{\mathbf U}_n-\mathbf U_n^\star\|_{\mathcal U}
    &\le
    \Lambda_D
    \sum_{j=0}^{n-1}
    \Lambda_T^{\,n-1-j}
    \left(
        b_j
        +
        \bar\alpha\,
        \sqrt{\operatorname{Tr}(K)}\left(1+2\sqrt{\log\frac{N}{\gamma}}+2\log\frac{N}{\gamma}\right)^{1/2}
    \right)
    +
    \delta_n ,
\end{aligned}
\]
with probability at least $1-\gamma$.
\end{remark}

\begin{proof}[Proof of Proposition \ref{prop:high-prob-autoregressive-rollout}]
Define the latent rollout error
\[
    \widetilde e_n
    :=
    \|\widetilde{\mathbf Z}_n-\mathbf Z_n^\star\|_{\mathcal Z}.
\]
Then
\[
    \widetilde e_0=\|\Xi_0\|_{\mathcal Z}.
\]
For $n=0,\ldots,N-1$, using the autoregressive rollout equation, the triangle inequality, the
Lipschitz continuity of $T_\theta$, and the bound on $\alpha_\theta$, we get
\[
\begin{aligned}
    \widetilde e_{n+1}
    &=
    \left\|
    T_\theta(\widetilde{\mathbf Z}_n)
    +
    \alpha_\theta(\widetilde{\mathbf Z}_n)\mathbf G_n
    -
    \mathbf Z_{n+1}^\star
    \right\|_{\mathcal Z}
    \\
    &\le
    \left\|
    T_\theta(\widetilde{\mathbf Z}_n)
    -
    T_\theta(\mathbf Z_n^\star)
    \right\|_{\mathcal Z}
    +
    \left\|
    T_\theta(\mathbf Z_n^\star)
    -
    \mathbf Z_{n+1}^\star
    \right\|_{\mathcal Z}
    +
    \alpha_\theta(\widetilde{\mathbf Z}_n)\|\mathbf G_n\|_{\mathcal Z}
    \\
    &\le
    \Lambda_T\widetilde e_n
    +
    b_n
    +
    \bar\alpha\|\mathbf G_n\|_{\mathcal Z}.
\end{aligned}
\]
Iterating this recursion yields
\[
    \widetilde e_n
    \le
    \Lambda_T^n\|\Xi_0\|_{\mathcal Z}
    +
    \sum_{j=0}^{n-1}
    \Lambda_T^{\,n-1-j}
    \left(
        b_j+\bar\alpha\|\mathbf G_j\|_{\mathcal Z}
    \right).
\]
For a Gaussian random element $X\sim\mathsf N(0,C)$ in a Hilbert space, let
$$r_C(q):=\left[\operatorname{Tr}(C)+2\sqrt{\operatorname{Tr}(C^2)\log(1/q)}+2\Vert C\Vert_\mathrm{op}\log(1/q)\right]^{1/2}.$$
By Lemma \ref{lem:hilbert-gaussian-tail}, we have the concentration bound
\[
    \mathbb P
    \left(
    \|X\|_{\mathcal Z}
    \le
    r_C(q)
    \right)
    \ge
    1-q.
\]
Applying this to $\Xi_0$ and to $\mathbf G_0,\ldots,\mathbf G_{N-1}$ with $q=\gamma/(N+1)$, and using a union bound, gives an event of probability at least $1-\gamma$
on which
\[
    \|\Xi_0\|_{\mathcal Z}
    \le
    r_{K_\phi(\mathbf U_0^\star)}\!\left(\frac{\gamma}{N+1}\right)\leq\varsigma^{1/2}\left(1+2\sqrt{\log\frac{N+1}{\gamma}}+2\log\frac{N+1}{\gamma}\right)^{1/2},
\]
and
\[
    \|\mathbf G_j\|_{\mathcal Z}
    \le
    r_K\!\left(\frac{\gamma}{N+1}\right),
    \qquad
    j=0,\ldots,N-1.
\]
Substituting these inequalities into the recursion gives the latent high-probability bound
\begin{equation*}
    \wt e_n\leq \Lambda_T^n
    \varsigma^{1/2}\left(1+2\sqrt{\log\frac{N+1}{\gamma}}+2\log\frac{N+1}{\gamma}\right)^{1/2}+
    \sum_{j=0}^{n-1}
    \Lambda_T^{n-1-j}
    \left(
        b_j
        +
        \bar\alpha\,
        r_K\left(\frac{\gamma}{N+1}\right)
    \right)
\end{equation*}
Finally,
\[
\begin{aligned}
    \|\widetilde{\mathbf U}_n-\mathbf U_n^\star\|_{\mathcal U}
    &\le
    \|\widetilde{\mathbf U}_n-\bar{\mathbf U}_n\|_{\mathcal U}
    +
    \|\bar{\mathbf U}_n-\mathbf U_n^\star\|_{\mathcal U}
    \\
    &=
    \|D_\psi(\widetilde{\mathbf Z}_n)-D_\psi(\mathbf Z_n^\star)\|_{\mathcal U}
    +
    \delta_n
    \le\Lambda_D\widetilde e_n+\delta_n.
\end{aligned}
\]
This proves the first claim. For Remark \ref{rem:probabilistic-rollout-bound}, if $\widetilde{\mathbf Z}_0=\mathbf Z_0^\star$, then
$\widetilde e_0=0$, so the initial encoder-noise term disappears. In that case, the union bound is
only over the $N$ transition noises, giving the version with $\gamma/N$.
\end{proof}

 \subsection{The KL-aware latent rollout bound}
\label{app:kl-aware-rollout}

We first recall a transportation inequality for Gaussian measures on Banach spaces.

\begin{theorem}[Riedel's Gaussian transportation inequality, Corollary 1.3 in \cite{riedel2017transportation}]
\label{thm:riedel-gaussian-t2}
Let $(B,\mathcal H,\gamma)$ be a Gaussian Banach space, where $B$ is a separable Banach space,
$\gamma$ is a centered Gaussian measure on $B$, and $\mathcal H$ is its Cameron--Martin space.
Then, for every probability measure $\nu\in\mathcal P(B)$,
\[
    \inf_{\pi\in\Pi(\nu,\gamma)}
    \int_{B\times B}
    \|x-y\|_B^2\,\d\pi(x,y)
    \le
    2\sigma_\gamma^2
    D_{\mathrm{KL}}(\nu\,\Vert\,\gamma),
\]
where
\[
    \sigma_\gamma^2
    :=
    \sup_{\ell\in B^\ast,\ \|\ell\|_{B^\ast}\le 1}
    \int_B \ell(x)^2\,\d\gamma(x)
    <\infty .
\]
Equivalently,
\[
    W_2^2(\nu,\gamma)
    \le
    2\sigma_\gamma^2
    D_{\mathrm{KL}}(\nu\,\Vert\,\gamma),
\]
where $W_2$ is computed with respect to the Banach norm $\|\cdot\|_B$.
\end{theorem}

\begin{lemma}[Hilbert-space specialization]
\label{lem:hilbert-gaussian-t2}
Let $\mathcal Z$ be a separable Hilbert space and let
\[
    \mu=\mathsf N(m,C),
\]
where $m\in\mathcal Z$ and $C:\mathcal Z\to\mathcal Z$ is positive, self-adjoint, and trace-class.
Then, for every $\nu\in\mathcal P_2(\mathcal Z)$,
\[
    W_2^2(\nu,\mu)
    \le
    2\|C\|_{\mathrm{op}}
    D_{\mathrm{KL}}(\nu\,\Vert\,\mu).
\]
Here $W_2$ is computed with respect to the Hilbert norm $\|\cdot\|_{\mathcal Z}$.
\end{lemma}

\begin{proof}
It suffices to consider the centered case $m=0$, since translation preserves both Wasserstein distance and relative entropy. Let
\[
    \gamma=\mathsf N(0,C).
\]
By Theorem~\ref{thm:riedel-gaussian-t2}, it remains to identify the weak variance $\sigma_\gamma^2$.

Since $\mathcal Z$ is a Hilbert space, every continuous linear functional $\ell\in\mathcal Z^\ast$ has the form
\[
    \ell_h(x)=\langle h,x\rangle_{\mathcal Z}
\]
for a unique $h\in\mathcal Z$, with
\[
    \|\ell_h\|_{\mathcal Z^\ast}=\|h\|_{\mathcal Z}.
\]
If $X\sim\mathsf N(0,C)$, then
\[
    \int_{\mathcal Z}\ell_h(x)^2\,\d\gamma(x)
    =
    \mathbb E\langle h,X\rangle_{\mathcal Z}^2
    =
    \langle Ch,h\rangle_{\mathcal Z}.
\]
Therefore
\[
\begin{aligned}
    \sigma_\gamma^2
    &=
    \sup_{\ell\in\mathcal Z^\ast,\ \|\ell\|_{\mathcal Z^\ast}\le 1}
    \int_{\mathcal Z}\ell(x)^2\,\d\gamma(x)\sup_{\|h\|_{\mathcal Z}\le 1}
    \langle Ch,h\rangle_{\mathcal Z}.
\end{aligned}
\]
Since $C$ is positive and self-adjoint,
\[
    \sup_{\|h\|_{\mathcal Z}\le 1}
    \langle Ch,h\rangle_{\mathcal Z}
    =
    \|C\|_{\mathrm{op}}.
\]
Thus
\[
    \sigma_\gamma^2=\|C\|_{\mathrm{op}}.
\]
Substituting this into Theorem~\ref{thm:riedel-gaussian-t2} gives
\[
    W_2^2(\nu,\mathsf N(0,C))
    \le
    2\|C\|_{\mathrm{op}}
    D_{\mathrm{KL}}(\nu\,\Vert\,\mathsf N(0,C)).
\]
The translated case $\mu=\mathsf N(m,C)$ follows by applying the same result to the translated measures.
\end{proof}


\begin{proof}[Proof of Lemma \ref{lem:kl_latent_mean_control}]
Let
\[
    P_\mathbf{z}:=P_\theta(\cdot\,|\,\mathbf{z})
    =
    \mathsf N(T_\theta(z),\alpha_\theta(z)^2K).
\]
By Lemma~\ref{lem:hilbert-gaussian-t2},
\[
    W_2^2(Q,P_z)\le
    2\|\alpha_\theta(\mathbf{z})^2K\|_{\mathrm{op}}
    D_{\mathrm{KL}}(Q\,\Vert\,P_z).
\]
Since
\[
    \|\alpha_\theta(\mathbf{z})^2K\|_{\mathrm{op}}
    =\alpha_\theta(\mathbf{z})^2\|K\|_{\mathrm{op}},
\]
we have
\[
    W_2(Q,P_\mathbf{z})
    \le
    \alpha_\theta(\mathbf{z})
    \sqrt{
    2\|K\|_{\mathrm{op}}
    D_{\mathrm{KL}}(Q\,\Vert\,P_\mathbf{z})
    }.
\]
It remains to observe that Wasserstein distance controls the distance between means. Indeed, for
any coupling $(Y,Y')$ of $Q$ and $P_z$,
\[
\begin{aligned}
    \|m_Q-T_\theta(\mathbf{z})\|_{\mathcal Z}
    &=
    \left\|
    \mathbb E Y-\mathbb E Y'
    \right\|_{\mathcal Z}  \\
    &\le
    \mathbb E\|Y-Y'\|_{\mathcal Z}  \\
    &\le
    \left(
    \mathbb E\|Y-Y'\|_{\mathcal Z}^2
    \right)^{1/2}.
\end{aligned}
\]
Taking the infimum over all couplings gives
\[
    \|m_Q-T_\theta(z)\|_{\mathcal Z}
    \le
    W_2(Q,P_z).
\]
Combining the two estimates proves the claim.
\end{proof}

We now apply this latent mean control to the learned transition dynamics. Let
\[
    \{\mathbf U_n^\star\}_{n=0}^N\subset\mathcal U
\]
be a reference physical trajectory and define
\[
    \mathbf Z_n^\star:=m_\phi(\mathbf U_n^\star),
    \qquad
    \bar{\mathbf U}_n:=D_\psi(\mathbf Z_n^\star),
    \qquad
    \delta_n:=\|\bar{\mathbf U}_n-\mathbf U_n^\star\|_{\mathcal U}.
\]
For each $n=0,\ldots,N-1$, let
\[
    \Xi_n\sim\mathsf N(0,K_\phi(\mathbf U_n^\star)).
\]
Define the dynamic KL quantity
\[
    \kappa_n^2
    :=
    \mathbb E_{\Xi_n}
    \left[
    D_{\mathrm{KL}}
    \left(
    Q_\phi(\cdot\mid \mathbf U_{n+1}^\star)
    \,\Vert\,
    P_\theta(\cdot\mid \mathbf Z_n^\star+\Xi_n)
    \right)
    \right],
\]
and define the latent encoder-noise sensitivity of the transition map by
\[
    \eta_n
    :=\left(\mathbb E_{\Xi_n}
    \left[\left\|T_\theta(\mathbf Z_n^\star+\Xi_n)
    -T_\theta(\mathbf Z_n^\star)
    \right\|_{\mathcal Z}^2
    \right]\right)^{1/2}.
\]

\begin{proof}[Proof of Proposition \ref{prop:population_kl_rollout}]
\textit{Step I.} We first bound the latent one-step defect
\[
    \|T_\theta(\mathbf Z_n^\star)-\mathbf Z_{n+1}^\star\|_{\mathcal Z}.
\]
For each realization of $\Xi_n$, apply Lemma~\ref{lem:kl_latent_mean_control} with
\[
    Q=Q_\phi(\cdot\mid \mathbf U_{n+1}^\star),
    \qquad
    \mathbf{Z}=\mathbf Z_n^\star+\Xi_n.
\]
Since the mean of $Q_\phi(\cdot\mid \mathbf U_{n+1}^\star)$ is
\[
    m_\phi(\mathbf U_{n+1}^\star)=\mathbf Z_{n+1}^\star,
\]
we obtain
\begin{equation}
\bigl\|\mathbf Z_{n+1}^\star
    -T_\theta(\widetilde{\mathbf Z}_n)
\bigr\|_{\mathcal Z}\leq
\alpha_\theta(\widetilde{\mathbf Z}_n)
\sqrt{
        2\|K\|_{\mathrm{op}}
        D_{\mathrm{KL}}
        \left(
            Q_\phi(\cdot\mid\mathbf U_{n+1}^\star)
            \,\Vert\,
            P_\theta(\cdot\mid\widetilde{\mathbf Z}_n)
        \right)
    }.
    \label{eq:proof_mean_control}
\end{equation}
Using $\alpha_\theta\leq\bar\alpha$, squaring and taking expectation over both the reference trajectory and $\Xi_n$ gives
\begin{equation}
\left(\mathbb E\left[\bigl\|
\mathbf Z_{n+1}^\star-T_\theta(\widetilde{\mathbf Z}_n)
\bigr\|_{\mathcal Z}^2\right]\right)^{1/2}
\leq\ol{\alpha}\sqrt{2\Vert K\Vert_\mathrm{op}}\,\kappa_n.
\label{eq:proof_kl_mean_population}
\end{equation}
Now write
\[
    T_\theta(\mathbf Z_n^\star)-\mathbf Z_{n+1}^\star
    =
    \bigl[
        T_\theta(\mathbf Z_n^\star)
        -
        T_\theta(\widetilde{\mathbf Z}_n)
    \bigr]
    +
    \bigl[
        T_\theta(\widetilde{\mathbf Z}_n)
        -
        \mathbf Z_{n+1}^\star
    \bigr].
\]
Taking the $L^2$ norm over the joint distribution and applying Minkowski's inequality together with \eqref{eq:proof_kl_mean_population} yields
\[
    \left(
        \mathbb E
        \left[
            \left\|
                T_\theta(\mathbf Z_n^\star)
                -
                \mathbf Z_{n+1}^\star
            \right\|_{\mathcal Z}^2
        \right]
    \right)^{1/2}
    \leq
    \eta_n+\ol{\alpha}\sqrt{2\Vert K\Vert_\mathrm{op}}\kappa_n,
\]
which proves \eqref{eq:population_latent_one_step}.

\item\textit{Step II.} We next derive the latent rollout bound. Define
\[
    e_n
    :=
    \left(
        \mathbb E
        \left[
            \|\widehat{\mathbf Z}_n-\mathbf Z_n^\star\|_{\mathcal Z}^2
        \right]
    \right)^{1/2}.
\]
Note that $e_0=0$ since the rollout is initialized at the encoder mean. For every realization of the reference trajectory,
\begin{align*}
    \|\widehat{\mathbf Z}_{n+1}-\mathbf Z_{n+1}^\star\|_{\mathcal Z}
    &=
    \|T_\theta(\widehat{\mathbf Z}_n)-\mathbf Z_{n+1}^\star\|_{\mathcal Z}
    \\
    &\leq
    \|T_\theta(\widehat{\mathbf Z}_n)
      -T_\theta(\mathbf Z_n^\star)\|_{\mathcal Z}
    +
    \|T_\theta(\mathbf Z_n^\star)
      -\mathbf Z_{n+1}^\star\|_{\mathcal Z}.
\end{align*}
Taking the population $L^2$ norm and using Minkowski's inequality, Lipschitz continuity of $T_\theta$, and \eqref{eq:population_latent_one_step}, we obtain
\[
    e_{n+1}
    \leq
    \Lambda_T e_n+\eta_n+\ol{\alpha}\sqrt{2\Vert K\Vert_\mathrm{op}}\,\kappa_n.
\]
Iterating this recursion from $e_0=0$ gives
\begin{equation}
    e_n
    \leq
    \sum_{j=0}^{n-1}
    \Lambda_T^{\,n-1-j}
    \left(\eta_j+\ol{\alpha}\sqrt{2\Vert K\Vert_\mathrm{op}}\,\kappa_j\right).\label{eq:population_latent_rollout}
\end{equation}
For the physical rollout, use $\widehat{\mathbf U}_n=D_\psi(\widehat{\mathbf Z}_n)$ and insert the decoded reference state $D_\psi(\mathbf Z_n^\star)$:
\begin{align*}
    \|\widehat{\mathbf U}_n-\mathbf U_n^\star\|_{\mathcal U}
    &\leq
    \|D_\psi(\widehat{\mathbf Z}_n)
      -D_\psi(\mathbf Z_n^\star)\|_{\mathcal U}
    +
    \|D_\psi(\mathbf Z_n^\star)-\mathbf U_n^\star\|_{\mathcal U}.
\end{align*}
Taking the population $L^2$ norm and using Lipschitz continuity of $D_\psi$ yields
\[
\left(\mathbb E
\|\widehat{\mathbf U}_n-\mathbf U_n^\star\|_{\mathcal U}^2
\right)^{1/2}
\leq\Lambda_D e_n+\delta_n.
\]
Substituting \eqref{eq:population_latent_rollout} proves \eqref{eq:population_kl_physical_rollout}.

\item\textit{Step III.} It remains to incorporate the predictive loss. Define the clean decoded one-step error
\[
    q_n
    :=
    \left(
        \mathbb E
        \left[
            \left\|
                D_\psi(T_\theta(\mathbf Z_n^\star))
                -
                \mathbf U_{n+1}^\star
            \right\|_{\mathcal U}^2
        \right]
    \right)^{1/2}.
\]
For every realization of the reference trajectory and the Gaussian
perturbations,
\begin{align*}
    &
    \left\|
        D_\psi(T_\theta(\mathbf Z_n^\star))
        -
        \mathbf U_{n+1}^\star
    \right\|_{\mathcal U}
    \\
    &\quad\leq
    \left\|
        D_\psi(T_\theta(\mathbf Z_n^\star))
        -
        D_\psi(T_\theta(\widetilde{\mathbf Z}_n))
    \right\|_{\mathcal U}
    \\
    &\qquad+
    \left\|
        D_\psi(T_\theta(\widetilde{\mathbf Z}_n))
        -
        D_\psi(\widetilde{\mathbf Z}_n^+)
    \right\|_{\mathcal U}
    +
    \left\|
        D_\psi(\widetilde{\mathbf Z}_n^+)
        -
        \mathbf U_{n+1}^\star
    \right\|_{\mathcal U}.
\end{align*}
Taking the joint $L^2$ norm and applying Minkowski's inequality gives
\begin{equation}
q_n\leq\Lambda_D\eta_n+\rho_n+\epsilon_n.\label{eq:population_predictive_one_step}
\end{equation}
We now obtain a second bound for the physical rollout error at time $n$. Since
\[
    \widehat{\mathbf U}_n
    =
    D_\psi(T_\theta(\widehat{\mathbf Z}_{n-1})),
\]
by Minkowski's inequality and Lipschitz continuity of $D_\psi$ and $T_\theta$, we have
\begin{align*}
    \left(
        \mathbb E
        \|\widehat{\mathbf U}_n-\mathbf U_n^\star\|_{\mathcal U}^2
    \right)^{1/2}&\leq \left(
        \mathbb E
        \|D_\psi(T_\theta(\wh{\mathbf Z}_n))-D_\psi(T_\theta(\mathbf Z_n^\star))\|_{\mathcal U}^2
    \right)^{1/2}+\left(
        \mathbb E
        \left[
            \left\|
                D_\psi(T_\theta(\mathbf Z_n^\star))
                -
                \mathbf U_{n+1}^\star
            \right\|_{\mathcal U}^2
        \right]
    \right)^{1/2}\\
    &\leq
    \Lambda_D\Lambda_T e_{n-1}+q_{n-1}.
\end{align*}
Using \eqref{eq:population_latent_rollout} at time $n-1$ and
\eqref{eq:population_predictive_one_step},
\begin{align}
    &
    \left(
        \mathbb E
        \|\widehat{\mathbf U}_n-\mathbf U_n^\star\|_{\mathcal U}^2
    \right)^{1/2}
    \nonumber\\
    &\quad\leq
    \Lambda_D
    \sum_{j=0}^{n-2}
    \Lambda_T^{\,n-1-j}
    \left(\eta_j+\bar\alpha\sqrt{2\|K\|_{\mathrm{op}}}\,\kappa_j\right)
    +
    \Lambda_D\eta_{n-1}
    +
    \epsilon_{n-1}
    +
    \rho_{n-1}.
    \label{eq:proof_predictive_route}
\end{align}
Both \eqref{eq:proof_predictive_route} and \eqref{eq:population_kl_physical_rollout} are valid upper bounds for the same population rollout error. Taking the smaller of their final terms yields \eqref{eq:population_refined_rollout}.
\end{proof}

The quantity $\kappa_n^2$ is the population analogue of the dynamic KL term in the VAMO training objective. Proposition~\ref{prop:population_kl_rollout} therefore makes explicit how the variational latent consistency term enters the rollout analysis: it controls the mismatch between the transition mean $T_\theta(\mathbf Z_n^\star+\Xi_n)$ and the encoder mean $\mathbf Z_{n+1}^\star=m_\phi(\mathbf U_{n+1}^\star)$ of the true next state. The additional term
$\eta_n$ appears because the dynamic KL is evaluated at noisy encoder inputs, whereas deterministic rollout uses the clean encoder mean $\mathbf Z_n^\star$.


\section{Benchmark Details}\label{app:benchmark_details}
\subsection{1D compressible Euler equation} 
We consider the one-dimensional compressible Euler equations, a canonical inviscid model for compressible fluid dynamics. The system describes the evolution of density, momentum, and total energy through conservation laws. In
conservative form, the equation is
\begin{equation}
    \partial_t\mathbf{U} + \partial_x \mathbf{F}(\mathbf{U}) = 0,
    \qquad x\in \mathbb{T}_L,
    \label{eq:euler1d_conservative}
\end{equation}
where the conserved variables and flux are given by
\begin{equation}
    \mathbf{U} =
    \begin{pmatrix}
        \rho \\
        \rho u \\
        E
    \end{pmatrix},
    \qquad
    \mathbf{F}(\mathbf{U}) =
    \begin{pmatrix}
        \rho u \\
        \rho u^2 + p \\
        u(E+p)
    \end{pmatrix},
    \label{eq:euler1d_flux}
\end{equation}
where $\rho(t,x)$ denotes the fluid density, $u(t,x)$ denotes the velocity, $p(t,x)$ denotes the pressure, and $E(t,x)$ denotes the total energy density. We close the system with the ideal gas equation of state
\begin{equation*}
    p = (\gamma-1)
    \left(
        E - \frac{1}{2}\rho u^2
    \right),
    \label{eq:euler1d_eos}
\end{equation*}
where  $\gamma = 1.4$ is the constant adiabatic index. Equivalently, the system can be written as
\begin{equation*}
\begin{cases}
    \partial_t \rho + \partial_x(\rho u) &= 0,\\
    \partial_t(\rho u) + \partial_x(\rho u^2+p) &= 0,\\
    \partial_t E + \partial_x\bigl(u(E+p)\bigr) &= 0.
\end{cases}
\label{eq:euler1d_expanded}
\end{equation*}
The Euler system contains no explicit viscous or thermal diffusion. Therefore, wave-like structures are not smoothed by physical dissipation, and the dynamics can develop sharp gradients, contact discontinuities, rarefaction waves, and shocks.

In our experiments, we use smooth periodic initial conditions on $\mathbb{T}_L=[0,L]$ and represent the state by the primitive variables
\begin{equation}
    \mathbf{V}(t,x) = \bigl(\rho(t,x), u(t,x), p(t,x)\bigr),\quad t\in [0,T],\ x\in\mathbb{T}_L.
\end{equation}
The learning task is to approximate the time-$h$ solution map
\begin{equation}
    \mathcal{S}_{h}:
    \bigl(\rho(t),u(t),p(t)\bigr)
    \mapsto
    \bigl(\rho(t+h),u(t+h),p(t+h)\bigr),
    \label{eq:euler1d_solution_map}
\end{equation}
which is then rolled out autoregressively. This benchmark provides a compact test of deterministic compressible wave dynamics. Compared with dissipative 1D equations such as Burgers' equation, the absence of viscosity makes long-horizon prediction more sensitive to phase errors and accumulated wave-interaction errors. At the same time, the one-dimensional setting avoids the full cost of multidimensional compressible flow, making it useful as a controlled inviscid-fluid benchmark.

\paragraph{Evaluation metrics.} For the Euler benchmark, we report rollout-aggregated relative $L^1$ and $L^2$ errors separately for density, velocity, and pressure, together with their channel average. Let $u_i^n$ and $\widehat u_i^n$ denote the reference and predicted fields for trajectory $i$ at rollout step $n$. The relative $L^1$ error is computed as
\[
    \operatorname{Rel}\text{-}L^1
    =
    \frac{1}{N_{\mathrm{test}}}
    \sum_{i=1}^{N_{\mathrm{test}}}
    \frac{
        \sum_{n=1}^{T}\|\widehat u_i^n-u_i^n\|_{L^1}
    }{
        \sum_{n=1}^{T}\|u_i^n\|_{L^1}
    },
\]
whereas the relative $L^2$ error is
\[
    \operatorname{Rel}\text{-}L^2
    =
    \frac{1}{N_{\mathrm{test}}}
    \sum_{i=1}^{N_{\mathrm{test}}}
    \frac{
        \left(
            \sum_{n=1}^{T}\|\widehat u_i^n-u_i^n\|_{L^2}^2
        \right)^{1/2}
    }{
        \left(
            \sum_{n=1}^{T}\|u_i^n\|_{L^2}^2
        \right)^{1/2}
    }.
\]
The $L^2$ metric emphasizes larger-amplitude pointwise errors, while $L^1$ is less dominated by a small number of large local discrepancies and is therefore useful when sharp fronts or discontinuous structures are present. We do not report $H^1$ error for Euler because the inviscid dynamics may develop shocks and contact discontinuities, for which derivative-based errors become highly sensitive to grid resolution, numerical shock thickness, and small spatial shifts of the discontinuity. In this setting, $L^1$ and $L^2$ provide more interpretable measures of field-level rollout accuracy.


\subsubsection{Data Generation for the 1D Euler Equations}
\label{app:euler_data_generation}
\paragraph{Periodic Gaussian-random-field initial conditions.}
For each trajectory, density, pressure, and velocity are generated from independent periodic random fields. Let
\[
    g_\rho,\quad g_p,\quad g_u
\]
denote three independently sampled zero-mean fields on $\mathbb{T}_L$. Each field is generated spectrally. For Fourier mode $k\in\mathbb{Z}$, we prescribe the spectral power
\[
    S_k
    =
    \sigma_{\mathrm{GRF}}^2
    \exp\left(
        -2\pi^2\ell_{\mathrm{GRF}}^2
        \frac{k^2}{L^2}
    \right),
\]
where
\[
    \ell_{\mathrm{GRF}}=1,
    \qquad
    \sigma_{\mathrm{GRF}}^2=1.
\]
We sample complex Gaussian Fourier coefficients according to
\[
    \widehat g_k
    =
    S_k^{1/2}\xi_k,
    \qquad
    \xi_k\sim\mathcal{N}_{\mathbb{C}}(0,1),
\]
subject to the conjugate-symmetry constraint required for a real-valued field. The zero-frequency coefficient is removed,
\[
    \widehat g_0=0,
\]
and the field is obtained through the inverse discrete Fourier transform.

Before sample-wise normalization, this construction corresponds to a stationary periodic Gaussian random field with covariance
\[
    K_L(r)
    \propto
    \sum_{k\in\mathbb{Z}\setminus\{0\}}
    \exp\left(
        -2\pi^2\ell_{\mathrm{GRF}}^2
        \frac{k^2}{L^2}
    \right)
    e^{2\pi i kr/L},
    \qquad
    r=x-y.
\]
This is the zero-mean periodic analogue of a squared-exponential covariance kernel. Each realization is subsequently centered numerically and normalized to unit root-mean-square amplitude:
\[
    g
    \leftarrow
    \frac{
        g-\frac{1}{n_x}\sum_{j=1}^{n_x}g(x_j)
    }{
        \left[
            \frac{1}{n_x}
            \sum_{j=1}^{n_x}
            \left(
                g(x_j)
                -
                \frac{1}{n_x}\sum_{m=1}^{n_x}g(x_m)
            \right)^2
        \right]^{1/2}
        +10^{-8}
    }.
\]
Because of this realization-wise normalization, the final perturbations are normalized transforms of GRF samples rather than exact samples from a single Gaussian measure. The primitive initial conditions are then constructed as
\[
\begin{cases}
\rho(0,x)
    =
    \rho_0
    \left(
        1+a_\rho g_\rho(x)
    \right),\\
    p(0,x)
    =
    p_0
    \left(
        1+a_p g_p(x)
    \right),\\
    u(0,x)
    =
    M c_0 g_u(x),
\end{cases}
\]
where
\[
    \rho_0=1,
    \qquad
    p_0=1,
    \qquad
    a_\rho=0.15,
    \qquad
    a_p=0.10,
\]
and
\[
    c_0
    =
    \sqrt{\frac{\gamma p_0}{\rho_0}}
\]
is the reference sound speed. The trajectory-level Mach amplitude is sampled independently as
\[
    M\sim\operatorname{Uniform}(0.05,0.35).
\]
Density and pressure are clipped below at
\[
    0.1\rho_0
    \quad\text{and}\quad
    0.1p_0,
\]
respectively, although the chosen perturbation amplitudes make this clipping inactive for typical normalized samples. The physical correlation length $\ell_{\mathrm{GRF}}=1$ is held fixed across all domain lengths. Therefore, increasing $L$ enlarges the physical domain while preserving the local correlation scale of the initial fields.

\paragraph{Spatial and temporal discretization.}
We use the spatial resolutions
\[
\begin{array}{c|cccc}
    L & 5 & 10 & 15 & 20\\
    \hline
    n_x & 1024 & 2048 & 3072 & 4096
\end{array}
\]
so that the grid spacing is identical in all four settings:
\[
    \Delta x
    =
    \frac{L}{n_x}
    =
    \frac{5}{1024}.
\]
The computational and stored resolutions are the same, so no spatial downsampling is applied. The trajectories are recorded every $\Delta t_{\mathrm{data}}=0.1$. The training and validation trajectories are generated on $t\in[0,1.5]$,
giving 15 transitions and 16 stored snapshots, while the test trajectories are generated on $t\in[0,10]$, giving 100 transitions and 101 stored snapshots.

\paragraph{Finite-volume solver.}
The Euler equations are solved using a first-order finite-volume method with the local Lax--Friedrichs, or Rusanov, numerical flux. At the interface between left and right states $\mathbf{U}_L$ and $\mathbf{U}_R$, the flux is
\[
    \widehat{\mathbf{F}}
    (\mathbf{U}_L,\mathbf{U}_R)
    =
    \frac12
    \left[
        \mathbf{F}(\mathbf{U}_L)
        +
        \mathbf{F}(\mathbf{U}_R)
    \right]
    -
    \frac12
    a_{\max}
    \left(
        \mathbf{U}_R-\mathbf{U}_L
    \right),
\]
where
\[
    a_{\max}
    =
    \max
    \left\{
        |u_L|+c_L,\,
        |u_R|+c_R
    \right\},
    \qquad
    c
    =
    \sqrt{\frac{\gamma p}{\rho}}.
\]
The conservative update is
\[
    \mathbf{U}_j^{m+1}
    =
    \mathbf{U}_j^m
    -
    \frac{\Delta t_m}{\Delta x}
    \left(
        \widehat{\mathbf{F}}_{j+1/2}^m
        -
        \widehat{\mathbf{F}}_{j-1/2}^m
    \right),
\]
with periodic boundary conditions.

The internal time step is selected adaptively according to
\[
    \Delta t_m
    =
    \min
    \left\{
        10^{-3},
        \frac{0.45\,\Delta x}
        {\max_j\left(|u_j^m|+c_j^m\right)},
        t_{\mathrm{next}}-t_m
    \right\},
\]
where $t_{\mathrm{next}}$ denotes the next output time. Thus, $10^{-3}$ is a maximum solver step, while the CFL condition may impose a smaller step.

After every numerical update, the conservative state is converted to primitive variables. To avoid invalid numerical states, density and pressure are lower-bounded by
\[
    \rho_{\min}=10^{-6},
    \qquad
    p_{\min}=10^{-6},
\]
and the conservative variables are reconstructed from the corrected primitive state before continuing the simulation. All numerical evolution is performed in double precision, and the stored trajectories are converted to single precision.

\paragraph{Dataset splits.}
For every domain length, we independently generate 1200 training trajectories and 300 validation trajectories over $[0,1.5]$, using random seed 123. A separate long-horizon test set of 200 trajectories is generated over $[0,10]$ using random seed 456. The test initial conditions are therefore independently sampled from the same GRF-based distribution as the training and validation initial conditions, while the temporal horizon is substantially longer. Each dataset split stores
\[
    \mathbf{u}_0\in\mathbb{R}^{3\times n_x}
\]
and
\[
    \mathbf{u}_{0:T}\in\mathbb{R}^{(N_t+1)\times 3\times n_x},
\]
where the three channels correspond to density, velocity, and pressure. A zero-valued forcing array is also stored for compatibility with the common dataset interface; the Euler dynamics themselves are unforced.

\subsection{2D compressible fluid dynamics} 
We consider the two-dimensional compressible Navier-Stokes equations following the Compressible Fluid Dynamics (CFD) setting used in PDEBench \citep{takamoto2022pdebench}. The system describes the evolution of density, velocity, pressure, and total energy for a viscous compressible fluid. Let $\rho(t,x)$ denote the mass density, $\mathbf{v}(t,x)=(v_1(t,x),v_2(t,x))$ the velocity field, $p(t,x)$ the gas pressure, and 
\[
    \epsilon = \frac{p}{\Gamma-1}
\]
the internal energy density, where $\Gamma$ is the heat capacity ratio. The governing equations are
\begin{equation}
\begin{cases}
    \partial_t \rho + \nabla\cdot(\rho \mathbf{v}) = 0,\\
    \displaystyle\rho\left(\partial_t \mathbf{v} + \mathbf{v}\cdot\nabla \mathbf{v}\right)
    =
    -\nabla p
    + \eta \Delta\mathbf{v}
    + \left(\zeta+\frac{\eta}{3}\right)\nabla(\nabla\cdot \mathbf{v}),\vspace{0.2cm}\\
    \displaystyle\partial_t\left(\epsilon+\frac{1}{2}\rho |\mathbf{v}|^2\right)
    +
    \nabla\cdot\left[\left(\epsilon+p+\frac{1}{2}\rho |\mathbf{v}|^2\right)\mathbf{v}
        - \mathbf{v}\cdot\sigma^\prime
    \right]
    =0.
\end{cases}
\label{eq:cfd2d}
\end{equation}
Here $\eta$ and $\zeta$ denote the shear and bulk viscosity coefficients, respectively, and $\sigma'$ is the viscous stress tensor. In our experiments, we focus on the two-dimensional periodic setting and represent the state using primitive variables
\[
\mathbf{U}(t,x)=\bigl(\rho(t,x), v_1(t,x), v_2(t,x), p(t,x)\bigr).
\]
The learning task is to approximate the time-$h$ solution map
\[
    \mathcal S_{h}:
    \bigl(\rho(t),v_1(t),v_2(t),p(t)\bigr)
    \mapsto
    \bigl(\rho(t+h),v_1(t+h),v_2(t+h),p(t+h)\bigr),
\]
which is then rolled out autoregressively. This benchmark tests whether a latent dynamics model can stabilize multi-field compressible flow prediction, where density, velocity, and pressure are coupled through conservation laws and viscous dissipation.

\paragraph{Evaluation metrics.} For the two-dimensional compressible-flow benchmark, we report
rollout-aggregated relative $L^2$ and $H^1$ errors for density, pressure, and velocity, together with their channel average. For
$\mathcal X\in\{L^2,H^1\}$, the metric is
\[
    \operatorname{Rel}\text{-}\mathcal X
    =
    \frac{1}{N_{\mathrm{test}}}
    \sum_{i=1}^{N_{\mathrm{test}}}
    \frac{
        \left(
            \sum_{n=1}^{T}
            \|\widehat u_i^n-u_i^n\|_{\mathcal X}^2
        \right)^{1/2}
    }{
        \left(
            \sum_{n=1}^{T}
            \|u_i^n\|_{\mathcal X}^2
        \right)^{1/2}
    },
\]
with
\[
    \|u\|_{H^1}^2
    =
    \|u\|_{L^2}^2
    +
    \|\nabla u\|_{L^2}^2.
\]
For the velocity field $\mathbf v=(v_x,v_y)$, the component norms are combined as
\[
    \|\mathbf v\|_{\mathcal X}^2
    =
    \|v_x\|_{\mathcal X}^2
    +
    \|v_y\|_{\mathcal X}^2.
\]
The relative $L^2$ error measures overall field accuracy, whereas the $H^1$ error additionally penalizes errors in spatial derivatives and is therefore more sensitive to the preservation of interfaces, gradients, and fine spatial structure. Unlike the inviscid Euler benchmark, the
viscous compressible-flow trajectories remain sufficiently regular for the derivative-sensitive $H^1$ metric to be meaningful.

\subsubsection{Data Generation for 2D Compressible Flow Benchmark}
\label{app:cfd2d_data_generation}

\paragraph{Governing equations.}
We generate trajectories from the two-dimensional compressible Navier-Stokes equations on the periodic unit torus $\mathbb{T}^2=[0,1]^2$.
The conservative state is
\[
    \mathbf{U}
    =
    \begin{pmatrix}
        \rho &
        \rho v_x &
        \rho v_y &
        E
    \end{pmatrix},
\]
where $\rho$ is the density, $\mathbf{v}=(v_x,v_y)$ is the velocity, and $E$ is the total energy density. The pressure is determined by the ideal-gas equation of state
\[
    p
    =
    (\gamma-1)
    \left(
        E-\frac{1}{2}\rho|\mathbf{v}|^2
    \right),
    \qquad
    \gamma=\frac{5}{3}.
\]
The viscous stress tensor is
\[
    \boldsymbol{\tau}
    =
    \mu
    \left(
        \nabla\mathbf{v}
        +
        \nabla\mathbf{v}^{\mathsf T}
    \right)
    +
    \left(
        \zeta-\frac{2}{3}\mu
    \right)
    (\nabla\cdot\mathbf{v})\mathbf{I},
\]
with shear and bulk viscosities
\[
    \mu=\zeta=10^{-8}.
\]
Heat conduction is included with thermal conductivity
\[
    \kappa
    =
    \frac{\mu\gamma}
    {(\gamma-1)\operatorname{Pr}},
    \qquad
    \operatorname{Pr}=0.72,
\]
which gives
\[
    \kappa
    \approx
    3.4722\times10^{-8}.
\]
\paragraph{Gaussian-random-field initial conditions.}
For each trajectory, we independently sample four periodic random fields
\[
    f_\rho,\quad
    f_{v_x},\quad
    f_{v_y},\quad
    f_p
\]
on $\mathbb{T}^2$. A single correlation length
\[
    \ell
    \sim
    \operatorname{Uniform}(0.05,0.15)
\]
is sampled for each trajectory and shared across its four physical channels. Conditional on $\ell$, each field is generated independently through a Gaussian spectral filter. For a Fourier mode $\mathbf{k}=(k_x,k_y)\in\mathbb{Z}^2,$
the spectral power is
\[
    S_\ell(\mathbf{k})
    =
    \exp\left(
        -2\pi^2\ell^2|\mathbf{k}|^2
    \right).
\]
We sample complex Gaussian coefficients
\[
    \widehat f(\mathbf{k})
    =
    S_\ell(\mathbf{k})^{1/2}
    \xi_{\mathbf{k}},
\]
where $\xi_{\mathbf{k}}$ are standard complex Gaussian variables subject to the conjugate-symmetry condition required for a real-valued field. The zero mode is removed,
\[
    \widehat f(\mathbf{0})=0,
\]
and the field is obtained by an inverse two-dimensional Fourier transform.

Before realization-wise normalization, this construction corresponds to a centered stationary periodic Gaussian random field with covariance
\[
    K_\ell(\mathbf{x}-\mathbf{y})
    \propto
    \sum_{\mathbf{k}\in\mathbb{Z}^2\setminus\{\mathbf{0}\}}
    \exp\left(
        -2\pi^2\ell^2|\mathbf{k}|^2
    \right)
    e^{2\pi i\mathbf{k}\cdot(\mathbf{x}-\mathbf{y})}.
\]
Each sampled field is subsequently divided by its empirical spatial standard deviation. Thus, the final perturbations have zero spatial mean and unit empirical variance, but are normalized transforms of GRF samples rather than exact samples from a fixed Gaussian measure.

The background density is
\[
    \rho_0=1,
\]
and the background pressure is chosen as
\[
    p_0=\frac{1}{\gamma}.
\]
Consequently, the reference sound speed is
\[
    c_0
    =
    \sqrt{\frac{\gamma p_0}{\rho_0}}
    =
    1.
\]
The primitive initial conditions are
\[
    \rho(0,\mathbf{x})
    =
    \rho_0
    +
    a_\rho f_\rho(\mathbf{x}),
    \qquad
    a_\rho=0.1,
\]
\[
    p(0,\mathbf{x})
    =
    p_0
    \left(
        1+a_p f_p(\mathbf{x})
    \right),
    \qquad
    a_p=0.05,
\]
and
\[
    v_x(0,\mathbf{x})
    =
    Mc_0 f_{v_x}(\mathbf{x}),
    \qquad
    v_y(0,\mathbf{x})
    =
    Mc_0 f_{v_y}(\mathbf{x}),
\]
with fixed Mach amplitude
\[
    M=0.1.
\]
Density and pressure are lower-bounded by $10^{-6}$ to avoid invalid numerical states.

Although the physical parameters are similar to the low-Mach compressible-flow setting considered in PDEBench, the trajectories used in this work are generated independently. In particular, we use Gaussian-random-field initial conditions rather than the mixed sinusoidal initial fields used in the original benchmark.

\paragraph{Numerical discretization.}
The equations are solved on a uniform $256\times256$
periodic grid. Inviscid fluxes are computed using an HLLC approximate Riemann solver with second-order MUSCL reconstruction of the primitive variables. The reconstructed slopes are limited using the minmod limiter. Viscous stresses and heat-conduction terms are discretized using second-order centered finite differences.

Time integration is performed using the two-stage strong-stability-preserving Runge-Kutta method. The internal time step is selected adaptively according to both acoustic and diffusive stability conditions. In particular, the advective time-step restriction is
\[
    \Delta t_{\mathrm{adv}}
    =
    0.45
    \frac{
        \min(\Delta x,\Delta y)
    }{
        \max_{\mathbf{x}}
        \left\{
            |v_x|+c,\,
            |v_y|+c
        \right\}
    },
\]
where
\[
    c=\sqrt{\frac{\gamma p}{\rho}}
\]
is the local sound speed. A diffusive restriction is also imposed using the largest viscous or thermal diffusion coefficient. The actual step is chosen as the minimum of these restrictions and the remaining time before the next recorded snapshot.

After each Runge-Kutta stage, the state is converted to primitive variables, density and pressure positivity floors are applied, and the conservative variables are reconstructed. Numerical evolution is performed in double precision, and recorded states are stored in single precision.

After solving, each trajectory is spatially downsampled from
$256\times256$ to $64\times64$. The learning models are trained and evaluated only on the resulting $64\times64$ fields.

\paragraph{Temporal sampling and dataset splits.}
Snapshots are recorded every $\Delta t_{\mathrm{data}}=0.1$, The training and validation trajectories cover $t\in[0,1]$, giving 10 one-step transitions and 11 stored snapshots per trajectory. We generate 1200 training trajectories and 300 validation trajectories using random seed 42.

The test set is generated independently using random seed 2042 and contains $200$ trajectories over $t\in[0,10]$.
Each test trajectory therefore contains 100 transitions and 101 stored snapshots. The training, validation, and test initial conditions follow the same GRF distribution; the out-of-distribution aspect of evaluation is the substantially longer temporal horizon.

A zero-valued forcing field is stored for compatibility with the common dataset interface. The compressible-flow equations used here are unforced.

\paragraph{Choice of physical parameters.}
We select the low-viscosity regime
\[
    \mu=\zeta=10^{-8}
\]
to maintain nontrivial compressible dynamics throughout the long test horizon. In preliminary experiments, larger viscosity coefficients rapidly dissipated the random initial perturbations, and the density, pressure, and velocity fields became nearly spatially uniform after a short evolution time. Such trajectories provide only a weak test of long-horizon autoregressive stability because predicting the late-time states approaches prediction of an approximately stationary field.

The selected viscosity therefore makes the benchmark deliberately more challenging: spatial structures and coupled density-pressure-velocity interactions remain present well beyond the training interval. We use the fixed low Mach number $M=0.1$ to remain in a stable low-Mach compressible regime while avoiding excessively strong shocks that could dominate the comparison through numerical-solver artifacts.



\subsection{2D incompressible Navier-Stokes equation} 
We consider the two-dimensional incompressible Navier-Stokes equations on the periodic domain $\mathbb{T}^2$. In velocity form, the forced incompressible Navier-Stokes equations are
\begin{equation}
    \partial_t\mathbf{u} + \mathbf{u}\cdot\nabla\mathbf{u}
    = -\nabla p + \nu \Delta\mathbf{u} + f,\quad \nabla\cdot\mathbf{u} = 0,
\label{eq:ns2d_velocity}
\end{equation}
where $\mathbf{u}(t,x)\in\mathbb{R}^2$ is the velocity field, $p(t,x)$ is the pressure, $\nu>0$ is the kinematic viscosity, and $f$ is an external body force. The incompressibility constraint enforces volume preservation of the flow and determines the pressure as a Lagrange multiplier.

Following the standard vorticity formulation, we define the scalar vorticity
\begin{equation}
    \omega = \nabla\times\mathbf{u}
    =\partial_x u_2 - \partial_y u_1 .
\end{equation}
Taking the curl of \eqref{eq:ns2d_velocity} eliminates the pressure term and gives the forced vorticity equation
\begin{equation}
    \partial_t \omega + \mathbf{u}\cdot\nabla \omega
    =
    \nu \Delta \omega + f,
    \qquad x\in\mathbb{T}^2,
    \label{eq:ns2d_vorticity}
\end{equation}
where $f=\nabla\times f_u$ is the vorticity forcing. The velocity field is recovered from vorticity through the stream function $\psi$:
\begin{equation}
    \mathbf{u} = \nabla^\perp \psi,
    \qquad
    -\Delta \psi = \omega,
    \label{eq:ns2d_biot_savart}
\end{equation}
where $\nabla^\perp \psi = (\partial_y\psi,-\partial_x\psi)$. Equivalently,
$u$ is obtained from $\omega$ by a nonlocal Biot--Savart operator.

In our experiments, we use the inhomogeneous setting with a static forcing field $f(x)$, so that the dynamics are governed by
\begin{equation}
    \partial_t \omega + \mathbf{u}\cdot\nabla \omega
    =
    \nu \Delta \omega + f(x),
    \qquad
    \mathbf{u} = \nabla^\perp(-\Delta)^{-1}\omega .
    \label{eq:ns2d_forced}
\end{equation}
The learning target is the time-$\Delta t$ solution map
\begin{equation}
    \mathcal{S}_{\Delta t}:
    \omega(t) \mapsto \omega(t+\Delta t),
    \label{eq:ns2d_solution_map}
\end{equation}
conditioned on the fixed forcing field. During evaluation, the predicted vorticity is rolled out autoregressively.

This benchmark is challenging because the predicted vorticity determines the velocity field that transports future vorticity. Thus, small off-manifold errors in $\omega$ can induce nonlocal velocity errors, which then feed back into the advection term $\mathbf{u}\cdot\nabla\omega$. Under weak viscosity and persistent forcing, these errors may accumulate over long horizons and generate localized high-frequency artifacts. We therefore evaluate not only pointwise rollout accuracy, but also spectral and derivative-sensitive diagnostics such as enstrophy and palinstrophy.

\paragraph{Evaluation metrics.} For the incompressible Navier-Stokes benchmark, we report relative $L^2$ and $H^1$ errors of the vorticity field using the same rollout-aggregated definition in the 2D compressible flow benchmark. These two metrics respectively measure field-level accuracy and derivative-level accuracy. We further evaluate enstrophy and palinstrophy:
\[
    \mathsf{Ens}(\omega)
    =
    \frac12\|\omega\|_{L^2}^2,
    \qquad
    \mathsf{Pal}(\omega)
    =
    \frac12\|\nabla\omega\|_{L^2}^2.
\]
Enstrophy measures the overall magnitude of the vorticity field, while palinstrophy emphasizes vorticity gradients and is therefore particularly sensitive to sharp local structures and spurious high-frequency content. For either scalar diagnostic $\mathsf{q}\in\{\mathsf{Ens},\mathsf{Pal}\}$, we compute
\[
    \operatorname{RelErr}_\mathsf{q}
    =
    \frac{1}{N_{\mathrm{test}}}
    \sum_{i=1}^{N_{\mathrm{test}}}
    \frac{
        \left(
            \sum_{n=1}^{T}
            |\mathsf{q}(\widehat\omega_i^n)-\mathsf{q}(\omega_i^n)|^2
        \right)^{1/2}
    }{
        \left(
            \sum_{n=1}^{T}
            |\mathsf{q}(\omega_i^n)|^2
        \right)^{1/2}
    }.
\]
In addition, we compare the isotropic energy spectra of the predicted and reference flows. Let $E_i^n(k)$ and $\widehat E_i^n(k)$ denote the radially aggregated spectral energies at wavenumber shell $k$. The spectral error is
\[
    \operatorname{SpecErr}
    =
    \frac{1}{N_{\mathrm{test}}}
    \sum_{i=1}^{N_{\mathrm{test}}}
    \frac{
        \left(
            \sum_{n=1}^{T}
            \sum_k
            |\widehat E_i^n(k)-E_i^n(k)|^2
        \right)^{1/2}
    }{
        \left(
            \sum_{n=1}^{T}
            \sum_k
            |E_i^n(k)|^2
        \right)^{1/2}
    }.
\]
This metric evaluates whether the predicted flow distributes energy correctly across spatial scales, complementing the pointwise, derivative-sensitive, and integral physical diagnostics.

\subsubsection{Data Generation for 2D Incompressible Navier-Stokes Benchmark}
\label{app:ns2d_data_generation}
\paragraph{Initial-vorticity distribution.}
For all viscosities and forcing families, the initial vorticity is sampled on the mean-zero subspace according to
\[
    \omega_0
    \sim
    \mathsf{N}\left(
        0,\,
        7^3(-\Delta+49I)^{-5/2}
    \right),
\]
which has the same spectral decay as the Gaussian prior used in \citep{li2020fourier} but a larger global amplitude. This choice reduces the mismatch between the weak initial vorticity fields and the higher-amplitude states produced later by the persistent forcing. When the smaller covariance scale is used, the initial fields are substantially more damped, and models trained only on the early portion of each trajectory experience a larger distribution shift during long-horizon rollout. In additional experiments, this smaller scale degraded all methods, while VAMO remained the best-performing model. We therefore use the larger scale in the main benchmark to avoid introducing an unnecessarily severe amplitude shift between the training and evaluation regimes. Equivalently, the standard deviation of Fourier mode
$\mathbf{k}\in\mathbb{Z}^2\setminus\{\mathbf{0}\}$ is proportional to
\[
    7^{3/2}
    \left(
        4\pi^2|\mathbf{k}|^2+49
    \right)^{-5/4}.
\]
The zero Fourier mode is removed so that every initial vorticity field has zero spatial mean. Initial conditions are sampled on a $256\times256$ grid and subsequently spectrally truncated to the $64\times64$ resolution used for learning.

\paragraph{Trajectory-dependent forcing fields.}
Unlike the original FNO benchmark, which uses a fixed forcing function, we independently sample one static forcing field for each trajectory. We consider two forcing families.
\begin{itemize}[topsep=0pt, itemsep=0pt]
\item For the \textsc{GRF} family, the forcing is sampled spectrally from a periodic squared-exponential Gaussian random field. Conditional on a length scale $\ell$, its spectral power is proportional to
\[
    \exp\left(
        -2\pi^2\ell^2|\mathbf{k}|^2
    \right).
\]
The zero mode is removed, and each realization is rescaled to a randomly sampled target $L^\infty$ amplitude.

\item For the \textsc{Wave} family, the forcing is a sparse mixture of periodic
Fourier modes:
\[
    f(\mathbf{x})
    =
    \sum_{j=1}^{J}
    a_j
    \cos\left(
        2\pi\mathbf{k}_j\cdot\mathbf{x}
        +
        \phi_j
    \right).
\]
The number of terms is sampled uniformly from $J\in\{1,2,3,4\}$. The wavevectors have nonzero integer components with maximum magnitude three, the phases are sampled uniformly from $[0,2\pi)$, and the coefficients are damped proportionally to $|\mathbf{k}_j|^{-2}$. Each realization is projected to zero mean and rescaled to a random target $L^\infty$ amplitude.
\end{itemize}


The forcing distributions depend slightly on viscosity. For
$\nu=10^{-4}$ and $\nu=10^{-5}$, GRF length scales are sampled from $[0.05,0.15]$, and both forcing families are rescaled to target $L^\infty$ amplitudes in $[0.08,0.20]$ for GRF forcing and $[0.08,0.16]$ for Wave forcing. For $\nu=10^{-3}$, stronger viscous dissipation causes trajectories to approach a slowly varying regime more rapidly. To retain nontrivial long-time dynamics, we therefore use stronger and more spatially correlated forcing: GRF length scales are sampled from $[0.05,0.25]$, and both GRF and Wave forcing amplitudes are sampled from $[0.10,0.30]$. Table~\ref{tab:ns2d_forcing_settings} summarizes the forcing
distributions.

\begin{table}[t]
    \centering
    \caption{
    Forcing distributions used for the two-dimensional incompressible Navier-Stokes datasets. Amplitude ranges refer to the target $L^\infty$ norm of each sampled forcing field.
    }
    \label{tab:ns2d_forcing_settings}
    \begin{tabular}{ccccc}
        \toprule
        Viscosity
        & Forcing
        & Length scale
        & Number of modes
        & Amplitude range \\
        \midrule
        $10^{-3}$
        & \textsc{GRF}
        & $[0.05,0.25]$
        & --
        & $[0.10,0.30]$ \\
        $10^{-3}$
        & \textsc{Wave}
        & --
        & $1$--$4$
        & $[0.10,0.30]$ \\
        $10^{-4},10^{-5}$
        & \textsc{GRF}
        & $[0.05,0.15]$
        & --
        & $[0.08,0.20]$ \\
        $10^{-4},10^{-5}$
        & \textsc{Wave}
        & --
        & $1$--$4$
        & $[0.08,0.16]$ \\
        \bottomrule
    \end{tabular}
\end{table}

\paragraph{Numerical solver.}
Trajectories are generated on a $256\times256$ periodic grid using a pseudospectral solver. The velocity and vorticity gradients are computed spectrally, while the nonlinear advection term is evaluated in physical space and transformed back to Fourier space. A two-thirds spectral mask is applied to de-alias the nonlinear term.

Diffusion is advanced using a Crank-Nicolson discretization, while the advection and forcing terms are treated explicitly. For an internal time step $\delta t$, the Fourier-space update has the form
\[
    \left(
        1+\frac{\delta t\,\nu}{2}|\mathbf{k}|^2
    \right) 
    \widehat\omega^{\,m+1}_{\mathbf{k}}
    =
    \left(
        1-\frac{\delta t\,\nu}{2}|\mathbf{k}|^2
    \right)
    \widehat\omega^{\,m}_{\mathbf{k}}
    -
    \delta t\,
    \widehat{
        \mathbf{v}^m\cdot\nabla\omega^m
    }_{\mathbf{k}}
    +
    \delta t\,\widehat f_{\mathbf{k}}.
\]
We use internal time step $\delta t=10^{-4}$. \jg{This matches the standard choice from the original FNO Navier-Stokes benchmark \citep{li2020fourier}; the resulting advective CFL number stays below $0.03$ across all viscosities considered here, well within the stability limit of the explicit advection sub-step.} Snapshots are recorded every one physical time unit. The zero vorticity mode is removed after every step to preserve the mean-free periodic formulation.

After simulation, the vorticity trajectories and forcing fields are downsampled from $256\times256$ to $64\times64$ through Fourier truncation. This retains the resolved low-frequency modes while removing unresolved high-frequency components.

\paragraph{Training and evaluation datasets.}
For each viscosity, we train a unified model on both forcing families. The training set contains 1,600 trajectories, with 800 \textsc{GRF}-forced and 800 \textsc{Wave}-forced trajectories. The validation set contains 400 trajectories, split evenly between the two families. At test time, the same model is evaluated separately on GRF- and Wave-forced trajectories.

The training horizons are $T_{\mathrm{train}}=10$, $12$, and $8$ for $\nu=10^{-3}$, $10^{-4}$, and $10^{-5}$, respectively. The corresponding long-horizon test sets extend to $T_{\mathrm{test}}=50$, $30$, and $20$. Since snapshots are stored at unit intervals, these settings contain 10, 12, and 8 one-step training transitions and 50, 30, and 20 autoregressive test transitions, respectively.

The initial-vorticity and forcing fields are sampled independently across training, validation, and test trajectories. The test distributions use the same viscosity-specific initial-condition and forcing priors as the training data; the out-of-distribution component is the substantially
longer temporal rollout.



\section{Implementation Details}\label{app:implementation_details}

This appendix gives the architectural and optimization details for the
models compared in Section~\ref{sec:experiments}. All models are implemented
in PyTorch and trained as one-step predictors on adjacent snapshots. At test
time, all methods are rolled out autoregressively from the exact initial
condition, using deterministic mean predictions.

\subsection{Common Data Handling and Model Selection}
For each benchmark, the stored trajectories are converted into one-step
training pairs
\[
    (\mathbf u_n,\mathbf f,\mathbf u_{n+1}),
\]
where $\mathbf f$ is the trajectory-dependent forcing field when present and a
zero-valued compatibility channel for the unforced compressible benchmarks. The
same one-step pairs, trajectory validation loader, and train/validation split
are used for all compared methods on a fixed benchmark.

All models are optimized with AdamW, learning rate $10^{-4}$, weight decay
$10^{-5}$, and gradient clipping with maximum norm $1$. We use a StepLR
scheduler with decay factor $0.5$ every 100 epochs. The number of training
epochs is benchmark-specific and reported in Table~\ref{tab:implementation_hyperparameters}.
Validation is performed after each epoch. The checkpoint used for reporting is
selected by the validation autoregressive rollout relative $L^2$ error when a
trajectory validation loader is available; otherwise it is selected by the
one-step validation loss.

For multichannel compressible states, the one-step and reconstruction losses
use channel-weighted mean-squared error. Unless user-specified weights are
provided, the channel weights are set to the inverse empirical variance of each
physical channel on the training trajectories and normalized to have unit mean.
This prevents high-variance channels from dominating the optimization objective.
For the scalar vorticity benchmark no channel reweighting is needed.

\subsection{FNO Backbone}
The direct FNO baseline maps the current physical state to the next physical
state,
\[
    \widehat{\mathbf u}_{n+1}
    =
    \mathbf u_n+\mathcal F_\theta(\mathbf u_n,\mathbf f,\mathbf x),
\]
where $\mathbf x$ denotes the periodic grid-coordinate channels. The residual
form is used in all reported experiments. The model first lifts the concatenated
input channels by a $1\times1$ convolution, applies a stack of Fourier layers
with pointwise local convolutions and GELU activations, and projects back to the
physical channels by two $1\times1$ convolutions. The benchmark-specific layer
counts, widths, and retained modes are reported in
Table~\ref{tab:implementation_hyperparameters} where currently populated.

\subsection{Deterministic Latent Baseline}
FNO-AE uses the same spatial latent representation as VAMO but removes all
stochastic components. The encoder is a resolution-preserving residual
convolutional network
\[
    E_\phi:\mathbf u_n\mapsto \mathbf z_n
    \in\mathbb R^{C_z\times N_{\mathrm{res}}},
\]
with residual convolutional blocks. For multichannel states, the encoder input
contains both the physical fields and centered finite-difference gradient
features; for scalar vorticity the same construction is applied channelwise. The
decoder has the mirrored residual-convolutional structure and output dimension
equal to the number of physical state channels. In the saved Navier--Stokes
runs, the encoder and decoder use four residual blocks, hidden width 64, and
$C_z=16$.

The deterministic latent transition is a residual FNO,
\[
    \mathbf z_{n+1}
    =
    \mathbf z_n+\mathcal G_\theta(\mathbf z_n,\mathbf f,\mathbf x),
\]
using the same retained modes and layer counts as the corresponding VAMO latent
transition. The deterministic latent objective is
\[
    \mathcal L_{\mathrm{AE}}
    =
    \mathcal L_{\mathrm{step}}
    +\lambda_{\mathrm{rec}}\mathcal L_{\mathrm{rec}}
    +\lambda_{\mathrm{spec}}\mathcal L_{\mathrm{spec}},
\]
where $\mathcal L_{\mathrm{step}}$ is the one-step prediction loss and
$\mathcal L_{\mathrm{rec}}$ reconstructs the current state from the encoded
latent. We use $\lambda_{\mathrm{rec}}=1$. The spectral Sobolev term is used
only for the incompressible Navier--Stokes deterministic latent baseline, with
$\lambda_{\mathrm{spec}}=1$ and Sobolev exponent $s=1$; it is set to zero for
the compressible benchmarks.

\subsection{VAMO Architecture and Structured Noise}
VAMO uses the same encoder, decoder, and residual latent FNO mean transition as
FNO-AE, but the encoder outputs both a mean field and a log-variance field:
\[
    (\mu_\phi(\mathbf u),\ell_\phi(\mathbf u))
    =
    \mathrm{Enc}_\phi(\mathbf u).
\]
During training, the posterior log-variance is clamped to $[-8,2]$ before
sampling. A posterior latent sample is generated by
\[
    \mathbf z_n
    =
    \mu_\phi(\mathbf u_n)
    +
    \exp\!\left(\frac12\ell_\phi(\mathbf u_n)\right)
    \odot \varepsilon_{\mathrm{enc}},
\]
where $\varepsilon_{\mathrm{enc}}$ is a spectrally filtered standard Gaussian
field. In 1D and 2D, respectively, this filtering is implemented by FFT as
\[
    \widehat{\varepsilon}(k)
    =
    \left(1+\ell^2 |2\pi k|^2\right)^{-s/2}\widehat{\xi}(k),
    \qquad
    \widehat{\varepsilon}(\mathbf k)
    =
    \left(1+\ell^2 |2\pi\mathbf k|^2\right)^{-s/2}
    \widehat{\xi}(\mathbf k),
\]
with orthonormal FFT normalization. This is the discrete implementation of the
Laplacian-resolvent covariance described in
Section~\ref{subsec:structured_latent_perturbations}.

The latent prior mean is the residual FNO transition
\[
    \mu_p=T_\theta(\mathbf z_n),
\]
and the transition noise amplitude is a scalar predicted separately for each
sample. The amplitude head consists of two convolutional layers with GELU
activations, global spatial averaging, and a final $1\times1$ convolution. In
the multichannel implementation used for the reported Euler, compressible-flow,
and incompressible Navier--Stokes experiments, the final scalar is passed
through a softplus and then shifted and clipped:
\[
    \alpha_\theta(\mathbf z,\mathbf f)
    =
    \min\{\alpha_{\max},\, \alpha_{\min}
    +\operatorname{softplus}(a_\theta(\mathbf z,\mathbf f))\}.
\]
This is the parameterization used by the \texttt{\_mc.py} model and training
scripts. The stochastic transition sample is then
\[
    \mathbf z_{n+1}
    =
    \mu_p
    +
    \alpha_\theta(\mathbf z_n,\mathbf f)\varepsilon_{\mathrm{tr}},
\]
where $\varepsilon_{\mathrm{tr}}$ is generated by the same spectral filter with
the transition correlation length. During training, the transition sampler adds
a small spectral variance floor $10^{-2}$ before taking the square root of the
filter variance. This prevents the highest retained Fourier modes from becoming
nearly deterministic. Inference uses the deterministic mean only: the initial
state is encoded as $\mu_\phi(\mathbf u_0)$, the latent is advanced by
$T_\theta$, and each latent state is decoded by $D_\psi$.

The VAMO training loss is
\[
    \mathcal L_{\mathrm{VAMO}}
    =
    \mathcal L_{\mathrm{step}}
    +\beta_{\mathrm{KL}}\mathcal L_{\mathrm{KL}}
    +\lambda_{\mathrm{rec}}\mathcal L_{\mathrm{rec}} .
\]
The prediction and reconstruction terms are the same channel-weighted MSE terms
used for the deterministic latent baseline. The KL term compares the encoded
posterior at the next state with the predicted latent prior. In the numerical
implementation this KL is evaluated using the encoder's pointwise diagonal
posterior variance and the transition's scalar per-sample variance
$\alpha_\theta^2$; the spectral covariance is used for reparameterized sampling
of the stochastic latent states.

\begin{table}[t]
    \centering
    \caption{
    Main implementation hyperparameters. ``Enc. length'' and ``tr. length''
    denote the encoder and transition spectral-noise correlation lengths,
    respectively. The Euler and 2D compressible-flow entries are left blank here
    as placeholders. The Navier--Stokes entries are read from the saved
    experiment configurations in \texttt{grad\_flow\_l2/ns2d\_per/results}.
    The spectral decay exponent is $s=2$ for the listed VAMO and FNO+Noise
    runs.
    }
    \label{tab:implementation_hyperparameters}
    \begin{tabular}{lccccccc}
        \toprule
        Benchmark
        & batch
        & epochs
        & latent ch.
        & FNO layers
        & modes
        & $\beta_{\mathrm{KL}}$
        & enc./tr. length \\
        \midrule
        1D Euler
        & & & & & & & \\
        2D CFD
        & & & & & & & \\
        2D Navier--Stokes, $\nu=10^{-3}$
        & 8 & 500 & 16 & 6 & $16\times16$
        & $10^{-3}$
        & $0.01/0.05$ \\
        2D Navier--Stokes, $\nu=10^{-4},10^{-5}$
        & 8 & 500 & 16 & 6 & $16\times16$
        & $10^{-2}$
        & $0.01/0.05$ \\
        \bottomrule
    \end{tabular}
\end{table}

For VAMO, we use $\lambda_{\mathrm{rec}}=1$ throughout. The Euler and 2D
compressible-flow transition-amplitude settings are left unspecified here. For
incompressible Navier--Stokes, the saved runs use $\alpha_{\min}=10^{-2}$ and
$\alpha_{\max}=5$. The softplus component is initialized to $0.10$, so the
initial effective amplitude before clipping is $0.11$.

\subsection{FNO+Noise Baseline}
FNO+Noise uses the same architecture, optimizer, data loaders, and model
selection rule as the direct FNO baseline. The only change is that during
training we add spectrally filtered Gaussian noise to the lifted FNO feature
field before the Fourier layers:
\[
    h \leftarrow h+\sigma_{\mathrm{lift}}\varepsilon_{\mathrm{lift}},
\]
where $\varepsilon_{\mathrm{lift}}$ is sampled with the same FFT filter family
as above. In the saved incompressible Navier--Stokes FNO+Noise runs, the lift
noise uses $\sigma_{\mathrm{lift}}=0.1$, correlation length $0.05$, and spectral
decay $s=2$. This noise is disabled during
validation and test rollout. The baseline therefore tests generic feature-space
noise injection without a latent posterior, learned transition uncertainty, or
KL alignment term.

\subsection{Rollout Safeguards}
The reported autoregressive evaluations use deterministic model outputs. During
validation/model selection, the training scripts optionally clip a single-step
increment in $L^\infty$ norm to avoid numerical overflow from already-diverged
checkpoints. The default clip thresholds are $5$ for the Euler models, $10$ for
the direct and deterministic-latent 2D models, and $1$ for the 2D
compressible-flow VAMO. This safeguard is applied only to the computed model
increment and is shared across the corresponding model-selection protocol; it
does not introduce stochastic forcing or teacher forcing.


\section{Additional Experimental Results}

\begin{figure}[p]
    \centering
    \includegraphics[width=\linewidth]
    {figs/euler1d_L5_sample_0040_trajectory_curves.png}
    \caption{
    Representative autoregressive rollout for the 1D Euler equations with $L=5$. Rows show density $\rho$, velocity $u$, and pressure $p$, while columns correspond to selected rollout times. The solid gray curve denotes the ground truth, and the colored dashed curves denote predictions from the evaluated models. All methods remain accurate at early times, whereas the deterministic latent baseline develops increasingly pronounced overshoots and oscillatory errors around sharp wave structures during the later rollout. VAMO remains closely aligned with the reference trajectory over the full prediction horizon.
    }
    \label{fig:euler1d_L5_qualitative}
\end{figure}

\begin{figure}[p]
    \centering
    \includegraphics[width=\linewidth]
    {figs/euler1d_L10_sample_0000_trajectory_curves.png}
    \caption{
    Representative autoregressive rollout for the 1D Euler equations with $L=10$. Rows show density $\rho$, velocity $u$, and pressure $p$, while columns correspond to selected rollout times. The solid gray curve denotes the ground truth, and the colored dashed curves denote predictions from the evaluated models. All methods remain accurate at early times, whereas the deterministic latent baseline develops increasingly pronounced overshoots and oscillatory errors around sharp wave structures during the later rollout. VAMO remains closely aligned with the reference trajectory over the full prediction horizon.
    }
    \label{fig:euler1d_L10_qualitative}
\end{figure}

\begin{figure}[p]
    \centering
    \includegraphics[width=\linewidth]
    {figs/euler1d_L15_sample_0011_trajectory_curves.png}
    \caption{
    Representative autoregressive rollout for the 1D Euler equations with $L=15$. Rows show density $\rho$, velocity $u$, and pressure $p$, while columns correspond to selected rollout times. The solid gray curve denotes the ground truth, and the colored dashed curves denote predictions from the evaluated models. All methods remain accurate at early times, whereas the deterministic latent baseline develops increasingly pronounced overshoots and oscillatory errors around sharp wave structures during the later rollout. VAMO remains closely aligned with the reference trajectory over the full prediction horizon.
    }
    \label{fig:euler1d_L15_qualitative}
\end{figure}

\begin{figure}[p]
    \centering
    \includegraphics[width=\linewidth]
    {figs/euler1d_L20_sample_0006_trajectory_curves.png}
    \caption{
    Representative autoregressive rollout for the 1D Euler equations with $L=20$. Rows show density $\rho$, velocity $u$, and pressure $p$, while columns correspond to selected rollout times. The solid gray curve denotes the ground truth, and the colored dashed curves denote predictions from the evaluated models. All methods remain accurate at early times, whereas the deterministic latent baseline develops increasingly pronounced overshoots and oscillatory errors around sharp wave structures during the later rollout. VAMO remains closely aligned with the reference trajectory over the full prediction horizon.
    }
    \label{fig:euler1d_L20_qualitative}
\end{figure}

\newpage 
\begin{figure}[ht]
    \centering
    \includegraphics[width=0.96\linewidth]
    {figs/cfd2d_sample_0010_trajectory_comparison.png}
    \caption{
    Representative autoregressive rollout for the 2D compressible fluid-dynamics benchmark. Rows compare the ground truth with predictions from Vanilla FNO, FNO+Noise, FNO-AE, and FNO-VAMO for density, the two velocity components, and pressure.
    }
    \label{fig:cfd2d_qualitative}
\end{figure}

\begin{figure}[p]
    \centering

    \begin{subfigure}[t]{\linewidth}
        \centering
        \includegraphics[width=\linewidth]
        {figs/ns2d_nu3_grf_sample_0177_trajectory.png}
        \caption{$\nu=10^{-3}$, \textsc{GRF} forcing, Sample 177.}
        \label{fig:ns2d_nu3_grf_trajectory}
    \end{subfigure}

    \vspace{0.8em}

    \begin{subfigure}[t]{\linewidth}
        \centering
        \includegraphics[width=\linewidth]
        {figs/ns2d_nu3_sinusoidal_sample_0070_trajectory.png}
        \caption{$\nu=10^{-3}$, \textsc{Wave} forcing, Sample 70.}
        \label{fig:ns2d_nu3_wave_trajectory}
    \end{subfigure}

    \caption{
    Representative autoregressive rollouts for the 2D incompressible Navier-Stokes equations with $\nu=10^{-3}$. Each panel shows the initial condition, the trajectory-specific forcing, the ground-truth vorticity, and predictions from the four evaluated methods at selected times. All predicted snapshots use the corresponding ground-truth color scale.
    }
    \label{fig:ns2d_nu3_qualitative}
\end{figure}

\begin{figure}[p]
    \centering

    \begin{subfigure}[t]{\linewidth}
        \centering
        \includegraphics[width=\linewidth]
        {figs/ns2d_nu4_grf_sample_0000_trajectory.png}
        \caption{$\nu=10^{-4}$, \textsc{GRF} forcing, Sample 0.}
        \label{fig:ns2d_nu4_grf_trajectory}
    \end{subfigure}

    \vspace{0.8em}

    \begin{subfigure}[t]{\linewidth}
        \centering
        \includegraphics[width=\linewidth]
        {figs/ns2d_nu4_sinusoidal_sample_0136_trajectory.png}
        \caption{$\nu=10^{-4}$, \textsc{Wave} forcing, Sample 136.}
        \label{fig:ns2d_nu4_wave_trajectory}
    \end{subfigure}

    \caption{
    Representative autoregressive rollouts for the 2D incompressible Navier-Stokes equations with $\nu=10^{-4}$. Each panel shows the initial condition, the trajectory-specific forcing, the ground-truth vorticity, and predictions from the four evaluated methods at selected times. All predicted snapshots use the corresponding ground-truth color scale.
    }
    \label{fig:ns2d_nu4_qualitative}
\end{figure}

% ------------------------------------------------------------
% nu = 10^{-5}
% ------------------------------------------------------------
\begin{figure}[p]
    \centering

    \begin{subfigure}[t]{\linewidth}
        \centering
        \includegraphics[width=\linewidth]
        {figs/ns2d_nu5_grf_sample_0000_trajectory.png}
        \caption{$\nu=10^{-5}$, \textsc{GRF} forcing, Sample 0.}
        \label{fig:ns2d_nu5_grf_trajectory}
    \end{subfigure}

    \vspace{0.8em}

    \begin{subfigure}[t]{\linewidth}
        \centering
        \includegraphics[width=\linewidth]
        {figs/ns2d_nu5_sinusoidal_sample_0136_trajectory.png}
        \caption{$\nu=10^{-5}$, \textsc{Wave} forcing, Sample 136.}
        \label{fig:ns2d_nu5_wave_trajectory}
    \end{subfigure}

    \caption{
    Representative autoregressive rollouts for the 2D incompressible Navier--Stokes equations with $\nu=10^{-5}$. Each panel shows the initial condition, the trajectory-specific forcing, the ground-truth vorticity, and predictions from the four evaluated methods at selected times. All predicted snapshots use the corresponding ground-truth color scale. 
    }
    \label{fig:ns2d_nu5_qualitative}
\end{figure}

\end{document}